% VI : rédiger un CV — Français Tle Tech — Cours signé OnBuch+
% Programme officiel MINESEC. Compiler avec : tectonic N18-vi-rediger-cv.tex
\newcommand{\DOCMATIERE}{Français}
\newcommand{\DOCNIVEAU}{Tle Tech}
\newcommand{\DOCMODULE}{Français – classe de Terminale technique}
\newcommand{\DOCLECON}{N18}
\newcommand{\DOCTITRE}{VI : rédiger un CV}
% OnBuch+ — préambule « cours signés » ENRICHI (compilé avec tectonic / XeLaTeX)
% Système visuel « Pop » d'OnBuch+ : fond crème (jamais blanc), bordures encre
% épaisses, ombres « dures » décalées, titres Archivo Black en capitales.
% Version MATHÉMATIQUES : graphes (pgfplots/tikz), boîtes pédagogiques
% (définition, propriété, méthode, attention, le savais-tu, exercice résolu,
% exercice, TP, bilan), page de garde et signature OnBuch+ ancrée en bas.
%
% Macros attendues dans main.tex AVANT \input{preamble} :
%   \DOCMATIERE  \DOCNIVEAU  \DOCMODULE  \DOCLECON (numéro)  \DOCTITRE
\documentclass[11pt]{article}

\usepackage{fontspec}
\usepackage[french]{babel}
\usepackage[a4paper,margin=2cm,top=3.4cm,bottom=2.8cm,headheight=1.4cm,footskip=1.3cm]{geometry}
\usepackage{amsmath,amssymb,amsfonts}
\usepackage{mathtools} % \xmapsto, \coloneqq... souvent utilisés par les modèles
\usepackage{cancel} % simplifications barrées (bilans de forces)
% \abs{z} (module, valeur absolue) et \norm{u} (norme), utilisés par les modèles
\providecommand{\abs}{}\let\abs\relax\DeclarePairedDelimiter{\abs}{\lvert}{\rvert}
\providecommand{\norm}{}\let\norm\relax\DeclarePairedDelimiter{\norm}{\lVert}{\rVert}
\usepackage[table]{xcolor} % [table] : fournit aussi \rowcolors (alternance de lignes)
\usepackage{pagecolor}
\usepackage{tikz}
\usetikzlibrary{arrows.meta,positioning,calc,shapes.geometric,shapes.misc,shapes.arrows,decorations.pathmorphing,decorations.pathreplacing,decorations.markings,patterns,shadows,mindmap,trees,fit,backgrounds,matrix,intersections,angles,quotes}
\usepackage{pgfplots}
\usepackage[version=4]{mhchem} % équations nucléaires \ce{^{235}_{92}U}
\usepackage[european,siunitx]{circuitikz} % schémas électriques
\pgfplotsset{compat=1.18}
\usepgfplotslibrary{fillbetween,groupplots} % aires entre courbes (calcul intégral)
\usepackage{enumitem}
\usepackage{fancyhdr}
% [most] charge toutes les bibliothèques tcolorbox (listings, minted, raster,
% documentation...) et gonfle inutilement la mémoire TeX sur un document long
% (~300 boîtes) : on ne charge que ce qui est réellement utilisé.
\usepackage[skins,breakable]{tcolorbox}
\usepackage{graphicx}
\usepackage{colortbl}
\usepackage{booktabs}
\usepackage{tabularx}
\usepackage{multirow}
\usepackage{diagbox} % case d'en-tête diagonale des tableaux à double entrée
% Colonnes extensibles centrée / gauche / droite (C, L, R), souvent utilisées par les modèles
\newcolumntype{C}{>{\centering\arraybackslash}X}
\newcolumntype{L}{>{\raggedright\arraybackslash}X}
\newcolumntype{R}{>{\raggedleft\arraybackslash}X}
\usepackage{array}
\usepackage{multicol}
\usepackage{siunitx}
% parse-numbers=false : accepte aussi \SI{2,5 \times 10^{-3}}{...}
\sisetup{locale=FR,output-decimal-marker={,},group-minimum-digits=4,parse-numbers=false}
\DeclareSIUnit\ohm{\ensuremath{\symup{\Omega}}} % U+2126 absent de la police
\DeclareSIUnit\division{div} % réglages d'oscilloscope (ms/div, V/div)
\DeclareSIUnit\tr{tr} % tours (vitesse de rotation, ex. \SI{1500}{\tr\per\minute})
\DeclareSIUnit\molar{M} % concentration molaire (ex. \SI{3}{\molar}, \SI{1}{\micro\molar})
\DeclareSIUnit\an{an} % année (le modèle confond parfois avec \year, un compteur TeX) — ex. \SI{5}{\kilo\meter\per\an}
\DeclareSIUnit\FCFA{FCFA} % franc CFA (ex. \SI{500}{\FCFA\per\kilo\gram})
\DeclareSIUnit\degreeBrix{\textdegree Brix} % teneur en sucre d'un sirop/jus (réfractométrie)
\DeclareSIUnit\month{mois} % le modèle utilise parfois \month, un compteur TeX, comme unité de durée
\DeclareSIUnit\microliter{\micro\liter} % le modèle écrit parfois \microliter en un seul token
\DeclareSIUnit\million{millions} % dénombrement (ex. \SI{15}{\million\per\mL} spermatozoïdes)
\usepackage{qrcode}
% \degree : commande fréquemment utilisée par les modèles pour un angle ou une
% température en dehors de \SI{}{} (non fournie par les paquets chargés).
\providecommand{\degree}{^{\circ}}
% \qed (fin de démonstration), utilisé par les modèles sans amsthm
\providecommand{\qed}{\hfill$\square$}
% \barre : double barre des valeurs interdites dans un tableau de variations (tkz-tab)
\providecommand{\barre}{\ensuremath{\|}}
% environnement proof (amsthm non chargé) utilisé par les modèles
\ifdefined\proof\else\newenvironment{proof}[1][Démonstration]{\par\noindent\textit{#1.}\ }{\qed\par}\fi
\usepackage{lastpage}
\usepackage{hyperref}
\hypersetup{hidelinks}

% ============================================================
% Palette « Pop » OnBuch+ — mêmes hex que la classe P (plus_ui.dart)
% ============================================================
\definecolor{poporange}{HTML}{F59321}
\definecolor{popdark}{HTML}{E07A0C}
\definecolor{popcream}{HTML}{FFF6E8}
\definecolor{popink}{HTML}{1C1714}
\definecolor{poppurple}{HTML}{7A5AE0}
\definecolor{popgreen}{HTML}{2E9E6A}
\definecolor{poppink}{HTML}{E0457B}
\definecolor{popblue}{HTML}{2F7DE1}
\definecolor{popgold}{HTML}{C98A12}
\definecolor{popmuted}{HTML}{8A7F76}
\definecolor{popline}{HTML}{EADFD2}
\definecolor{popgray}{HTML}{8A7F76}
\colorlet{popgrey}{popgray}
% Alias tolérants (le modèle nomme parfois les couleurs différemment)
\colorlet{popwhite}{white}
\colorlet{popblack}{popink}
\colorlet{popred}{poppink}
\colorlet{popyellow}{popgold}
% Teintes claires (fonds de tableaux/encadrés) — le modèle nomme parfois
% les couleurs avec un suffixe L (ex. popblueL) sans les avoir définies.
\colorlet{poporangeL}{poporange!20}
\colorlet{popdarkL}{popdark!20}
\colorlet{poppurpleL}{poppurple!20}
\colorlet{popgreenL}{popgreen!20}
\colorlet{poppinkL}{poppink!20}
\colorlet{popblueL}{popblue!20}
\colorlet{popgoldL}{popgold!20}
\colorlet{popredL}{popred!20}
\colorlet{popgrayL}{popgray!20}
% Teintes claires pour fonds de tableaux / zones
\colorlet{popblueL}{popblue!10!white}
\colorlet{poporangeL}{poporange!14!white}
\colorlet{popgreenL}{popgreen!12!white}
\colorlet{poppurpleL}{poppurple!12!white}
\colorlet{poppinkL}{poppink!12!white}
\colorlet{popgoldL}{popgold!14!white}

\pagecolor{popcream}

% ============================================================
% Typographie
% ============================================================
\defaultfontfeatures{Ligatures=TeX}
\setmainfont{PlusJakartaSans-Medium.ttf}[Path=../../../fonts/,BoldFont=PlusJakartaSans-Bold.ttf]
% Maths en sans-serif (Fira Math) avec chiffres et lettres droites de Plus
% Jakarta, pour que les formules se fondent dans le texte.
\usepackage[math-style=ISO,bold-style=ISO]{unicode-math}
\setmathfont{FiraMath-Regular.otf}[Path=../../../fonts/]
\setmathfont{PlusJakartaSans-Medium.ttf}[Path=../../../fonts/,range={up/num,up/{Latin,latin},\mathup}]
\usepackage{icomma} % virgule décimale sans espace en mode math
% Caractères mathématiques Unicode absents des polices de texte (Plus Jakarta,
% Archivo Black) : rendus par leur équivalent mathématique, en texte comme en
% formule — sinon ils s'affichent en carré vide (ex. « Arithmétique dans ℤ »).
\usepackage{newunicodechar}
\newunicodechar{ℝ}{\ensuremath{\mathbb{R}}}
\newunicodechar{ℤ}{\ensuremath{\mathbb{Z}}}
\newunicodechar{ℂ}{\ensuremath{\mathbb{C}}}
\newunicodechar{ℕ}{\ensuremath{\mathbb{N}}}
\newunicodechar{ℚ}{\ensuremath{\mathbb{Q}}}
\newunicodechar{𝔻}{\ensuremath{\mathbb{D}}}
\newunicodechar{α}{\ensuremath{\alpha}}
\newunicodechar{β}{\ensuremath{\beta}}
\newunicodechar{γ}{\ensuremath{\gamma}}
\newunicodechar{δ}{\ensuremath{\delta}}
\newunicodechar{ε}{\ensuremath{\varepsilon}}
\newunicodechar{θ}{\ensuremath{\theta}}
\newunicodechar{λ}{\ensuremath{\lambda}}
\newunicodechar{μ}{\ensuremath{\mu}}
\newunicodechar{σ}{\ensuremath{\sigma}}
\newunicodechar{τ}{\ensuremath{\tau}}
\newunicodechar{φ}{\ensuremath{\varphi}}
\newunicodechar{ψ}{\ensuremath{\psi}}
\newunicodechar{ω}{\ensuremath{\omega}}
\newunicodechar{Σ}{\ensuremath{\Sigma}}
% majuscules produites par \MakeUppercase dans les titres (α-aminés -> Α)
\newunicodechar{Α}{\ensuremath{\alpha}}
\newunicodechar{Β}{\ensuremath{\beta}}
\newunicodechar{Γ}{\ensuremath{\Gamma}}
\newunicodechar{Δ}{\ensuremath{\Delta}}
\newunicodechar{Ε}{\ensuremath{\varepsilon}}
\newunicodechar{Θ}{\ensuremath{\Theta}}
\newunicodechar{Λ}{\ensuremath{\Lambda}}
\newunicodechar{Μ}{\ensuremath{\mu}}
\newunicodechar{Τ}{\ensuremath{\tau}}
\newunicodechar{Φ}{\ensuremath{\Phi}}
\newunicodechar{Ψ}{\ensuremath{\Psi}}
\newunicodechar{Ω}{\ensuremath{\Omega}}
\newunicodechar{↦}{\ensuremath{\mapsto}}
\newunicodechar{∈}{\ensuremath{\in}}
\newunicodechar{∉}{\ensuremath{\notin}}
\newunicodechar{∧}{\ensuremath{\wedge}}
\newunicodechar{⇔}{\ensuremath{\Leftrightarrow}}
\newunicodechar{⟺}{\ensuremath{\Longleftrightarrow}}
\newunicodechar{⇒}{\ensuremath{\Rightarrow}}
\newunicodechar{⟹}{\ensuremath{\Longrightarrow}}
\newunicodechar{∘}{\ensuremath{\circ}}
\newunicodechar{∪}{\ensuremath{\cup}}
\newunicodechar{∩}{\ensuremath{\cap}}
\newunicodechar{≡}{\ensuremath{\equiv}}
\newunicodechar{⊂}{\ensuremath{\subset}}
\newunicodechar{⊥}{\ensuremath{\perp}}
\newunicodechar{∃}{\ensuremath{\exists}}
\newunicodechar{∀}{\ensuremath{\forall}}
\newunicodechar{∛}{\ensuremath{\sqrt[3]{\,}}}
\newunicodechar{∜}{\ensuremath{\sqrt[4]{\,}}}
\newunicodechar{⃗}{\ensuremath{{}^{\to}}}
\newunicodechar{ⁿ}{\ifmmode^{n}\else\textsuperscript{\ensuremath{n}}\fi}
\newunicodechar{ˣ}{\ifmmode^{x}\else\textsuperscript{\ensuremath{x}}\fi}
\newunicodechar{ᵇ}{\ifmmode^{b}\else\textsuperscript{\ensuremath{b}}\fi}
\newunicodechar{ʳ}{\ifmmode^{r}\else\textsuperscript{\ensuremath{r}}\fi}
\newunicodechar{ᵘ}{\ifmmode^{u}\else\textsuperscript{\ensuremath{u}}\fi}
\newunicodechar{ʸ}{\ifmmode^{y}\else\textsuperscript{\ensuremath{y}}\fi}
\newunicodechar{ᵃ}{\ifmmode^{a}\else\textsuperscript{\ensuremath{a}}\fi}
\newunicodechar{ᵉ}{\ifmmode^{e}\else\textsuperscript{\ensuremath{e}}\fi}
\newunicodechar{ᵏ}{\ifmmode^{k}\else\textsuperscript{\ensuremath{k}}\fi}
\newunicodechar{ᵖ}{\ifmmode^{p}\else\textsuperscript{\ensuremath{p}}\fi}
\newunicodechar{ᴺ}{\ifmmode^{N}\else\textsuperscript{\ensuremath{N}}\fi}
\newunicodechar{ᶜ}{\ifmmode^{c}\else\textsuperscript{\ensuremath{c}}\fi}
\newunicodechar{ʰ}{\ifmmode^{h}\else\textsuperscript{\ensuremath{h}}\fi}
\newunicodechar{ᵠ}{\ifmmode^{q}\else\textsuperscript{\ensuremath{q}}\fi}
\newunicodechar{ₙ}{\ifmmode_{n}\else\textsubscript{\ensuremath{n}}\fi}
\newunicodechar{ₐ}{\ifmmode_{a}\else\textsubscript{\ensuremath{a}}\fi}
\newunicodechar{ᵢ}{\ifmmode_{i}\else\textsubscript{\ensuremath{i}}\fi}
\newunicodechar{ᵣ}{\ifmmode_{r}\else\textsubscript{\ensuremath{r}}\fi}
\newunicodechar{ₖ}{\ifmmode_{k}\else\textsubscript{\ensuremath{k}}\fi}
\newunicodechar{ᵦ}{\ifmmode_{\beta}\else\textsubscript{\ensuremath{\beta}}\fi}
\newunicodechar{ₓ}{\ifmmode_{x}\else\textsubscript{\ensuremath{x}}\fi}
\newfontfamily\popdisplay{ArchivoBlack-Regular.ttf}[Path=../../../fonts/]
\newfontfamily\poplogo{SpaceGrotesk-Bold.ttf}[Path=../../../fonts/]
\newfontfamily\popsemi{PlusJakartaSans-SemiBold.ttf}[Path=../../../fonts/]
\newfontfamily\popxbold{PlusJakartaSans-ExtraBold.ttf}[Path=../../../fonts/]
\newfontfamily\popxboldls{PlusJakartaSans-ExtraBold.ttf}[Path=../../../fonts/,LetterSpace=3.0]

\setlist[enumerate]{leftmargin=1.7em,itemsep=5pt,topsep=4pt}
\setlist[itemize]{leftmargin=1.7em,itemsep=4pt,topsep=4pt,label={\color{poporange}\textbullet}}
\setlength{\parindent}{0pt}
\setlength{\parskip}{5pt}
\renewcommand{\arraystretch}{1.25}

% pgfplots : style Pop par défaut
\pgfplotsset{
  every axis/.append style={axis line style={popink,line width=1pt},tick style={popink},
    grid style={popline,line width=0.5pt},label style={font=\small},tick label style={font=\footnotesize}},
  popcurve/.style={poporange,line width=1.8pt,no markers,smooth},
}
% Même style disponible dans un \draw TikZ simple (hors pgfplots)
\tikzset{popcurve/.style={poporange,line width=1.8pt}}
\tikzset{popgrid/.style={popline,very thin}}
\colorlet{popcurve}{poporange} % le nom du style est parfois utilisé comme couleur

% ============================================================
% Pill / squiggle
% ============================================================
\newcommand{\poppill}[3]{%
  \tikz[baseline=-0.55ex]\node[draw=popink,line width=1.2pt,rounded corners=2.6pt,%
    fill=#2,inner xsep=6.5pt,inner ysep=3pt]{%
    \popxboldls\fontsize{7}{7}\selectfont\color{#3}\MakeUppercase{#1}};%
}
\newcommand{\popsquiggle}[1]{%
  \tikz[baseline=-0.4ex]{
    \draw[#1,line width=2.6pt,line cap=round]
      plot [smooth] coordinates {(0,0.09) (0.34,0.01) (0.68,0.21) (1.02,0.01) (1.36,0.10)};
  }%
}
% Mot-clé surligné (marqueur orange sous le texte)
\newcommand{\cle}[1]{{\popxbold\color{popdark}#1}}

% ============================================================
% En-tête / pied
% ============================================================
\pagestyle{fancy}
\fancyhf{}
\renewcommand{\headrulewidth}{0pt}
\renewcommand{\footrulewidth}{0pt}
\lhead{\raisebox{-0.15ex}{\poplogo\fontsize{13}{13}\selectfont\color{popink}OnBuch\color{poporange}+}}
\rhead{\poppill{\DOCMATIERE\ \textbullet\ \DOCNIVEAU\ \textbullet\ Leçon \DOCLECON}{white}{popink}}
\lfoot{\popxbold\fontsize{7.6}{7.6}\selectfont\color{popmuted}\MakeUppercase{Cours signé OnBuch+}\ \textbullet\ {\popsemi\DOCTITRE}}
\rfoot{\tikz[baseline=-0.6ex]\node[rounded corners=8.5pt,fill=popink,minimum height=17pt,minimum width=17pt,inner xsep=5pt,inner sep=0]{\popxbold\fontsize{8}{8}\selectfont\color{popcream}\thepage\,/\,\pageref*{LastPage}};}

% ============================================================
% Titres
% ============================================================
\newcommand{\chaptitre}[2]{%
  \begin{tcolorbox}[enhanced,colback=poporange,colframe=popink,boxrule=2.6pt,arc=11pt,
      drop shadow=popink,left=15pt,right=15pt,top=12pt,bottom=11pt,before skip=3pt,after skip=18pt]
    \poppill{#2}{popink}{popcream}\par\vspace{9pt}
    {\raggedright\hyphenpenalty=10000\exhyphenpenalty=10000\popdisplay\fontsize{22}{24}\selectfont\color{popcream}\MakeUppercase{#1}\par}
  \end{tcolorbox}
}
% Section numérotée : pastille orange avec le numéro + titre Archivo
\newcounter{popsec}
\newcommand{\coursec}[1]{%
  \stepcounter{popsec}\setcounter{popsub}{0}%
  \par\vspace{16pt}\noindent
  \begin{minipage}{\linewidth}
  \tikz[baseline=-0.7ex]\node[draw=popink,line width=1.6pt,fill=poporange,rounded corners=4pt,minimum size=22pt,inner sep=0]{\popdisplay\fontsize{12}{12}\selectfont\color{popcream}\thepopsec};%
  \hspace{8pt}{\popdisplay\fontsize{15}{17}\selectfont\color{popink}\MakeUppercase{#1}}\par
  \vspace{4pt}\noindent{\color{poporange}\rule{36pt}{3.6pt}}
  \end{minipage}\par\nopagebreak\vspace{8pt}
}
\newcounter{popsub}
\newcommand{\courssub}[1]{%
  \stepcounter{popsub}%
  \par\vspace{9pt}\noindent{\popxbold\fontsize{11.5}{13}\selectfont\color{popink}{\color{poporange}\thepopsec.\thepopsub}\hspace{6pt}#1}\par\nopagebreak\vspace{4pt}
}

% ============================================================
% Boîtes pédagogiques — même recette : fond blanc, bordure épaisse colorée,
% ombre dure encre, badge plein. Titre optionnel [#1].
% ============================================================
\tcbset{popbase/.style={breakable,enhanced,colback=white,boxrule=2.3pt,arc=9pt,
  shadow={3pt}{-3pt}{0pt}{popink},left=14pt,right=14pt,top=11pt,bottom=11pt,before skip=11pt,after skip=15pt,
  fontupper=\popsemi\fontsize{10.6}{14.5}\selectfont\color{popink},
  fontlower=\popsemi\fontsize{10.6}{14.5}\selectfont\color{popink}}}
% Coche dessinée (le glyphe ✓ n'existe pas dans les polices du document)
\newcommand{\popcheck}[1]{\tikz[baseline=-0.5ex]\draw[#1,line width=1.6pt,line cap=round,line join=round](-0.13,0.0)--(-0.04,-0.09)--(0.13,0.1);}
\providecommand{\coche}{\popcheck{popgreen}}
\providecommand{\resultat}[1]{\par\smallskip\noindent\fcolorbox{popgreen}{popgreenL}{\parbox{\dimexpr\linewidth-2\fboxsep-2\fboxrule\relax}{\textbf{Résultat.} #1}}\par\smallskip}
\newcommand{\popboxhead}[3]{\poppill{#1}{#2}{white}\ifx\relax#3\relax\else\hspace{6pt}{\popxbold\fontsize{10.6}{12}\selectfont\color{popink}#3}\fi\par\vspace{7pt}}

\newtcolorbox{aretenir}[1][]{popbase,colframe=popgreen,before upper={\popboxhead{À retenir}{popgreen}{#1}}}
\newtcolorbox{exemplebox}[1][]{popbase,colframe=poporange,before upper={\popboxhead{Exemple}{poporange}{#1}}}
\newtcolorbox{definition}[1][]{popbase,colframe=popblue,before upper={\popboxhead{Définition}{popblue}{#1}}}
\newtcolorbox{propriete}[1][]{popbase,colframe=poppurple,before upper={\popboxhead{Propriété}{poppurple}{#1}}}
\newtcolorbox{methode}[1][]{popbase,colframe=popgold,before upper={\popboxhead{Méthode}{popgold}{#1}}}
\newtcolorbox{attention}[1][]{popbase,colframe=poppink,before upper={\popboxhead{Attention}{poppink}{#1}}}
\newtcolorbox{experience}[1][]{popbase,colframe=popblue,colback=popblueL,before upper={\popboxhead{Expérience}{popblue}{#1}}}
% « Le savais-tu ? » : carte violette pleine (contexte, culture, Cameroun)
\newtcolorbox{savaistu}[1][]{popbase,colback=poppurple,colframe=popink,
  fontupper=\popsemi\fontsize{10.4}{14.2}\selectfont\color{white},
  before upper={\poppill{Le savais-tu\,?}{white}{poppurple}\ifx\relax#1\relax\else\hspace{6pt}{\popxbold\color{white}#1}\fi\par\vspace{7pt}}}
% Exercice résolu : énoncé en haut, \tcblower puis la solution sur fond crème
\newcounter{popexo}\newcounter{popexr}
\newtcolorbox{exoresolu}[1][]{popbase,colframe=poporange,colbacklower=popcream,
  segmentation style={popink,line width=1.2pt,dashed},
  before upper={\stepcounter{popexr}\popboxhead{Exercice résolu \thepopexr}{poporange}{#1}},
  before lower={\poppill{Solution}{popink}{popcream}\par\vspace{6pt}}}
% Exercice (non résolu en place) — la correction est dans la partie Corrigés
\newtcolorbox{exercice}[1][]{popbase,colframe=popink,
  before upper={\stepcounter{popexo}\popboxhead{Exercice \thepopexo}{popink}{#1}}}
% Corrigé
\newtcolorbox{corrige}[1][]{popbase,colframe=popgreen,colback=popgreenL,
  before upper={\popboxhead{Corrigé}{popgreen}{#1}}}
% Objectifs / prérequis / bilan
\newtcolorbox{objectifs}{popbase,colframe=popink,colback=poporangeL,
  before upper={\popboxhead{Objectifs de la leçon}{popink}{}}}
\newtcolorbox{prerequis}{popbase,colframe=popink,colback=popblueL,
  before upper={\popboxhead{Prérequis}{popblue}{}}}
\newtcolorbox{bilan}{popbase,colframe=popink,colback=white,boxrule=2.8pt,
  before upper={\popboxhead{Fiche bilan}{poporange}{L'essentiel à connaître pour l'examen}}}
% Figure encadrée avec légende
\newtcolorbox{popfigure}[1][]{enhanced,breakable=false,colback=white,colframe=popline,boxrule=1.4pt,arc=8pt,
  left=10pt,right=10pt,top=10pt,bottom=8pt,before skip=10pt,after skip=12pt,halign=center,
  lower separated=false,halign lower=center,
  fontlower=\popsemi\fontsize{9}{11.5}\selectfont\color{popmuted},
  #1}
\newcommand{\legende}[1]{\par\vspace{6pt}{\popsemi\fontsize{9}{11.5}\selectfont\color{popmuted}#1}}

% ============================================================
% Signature OnBuch+
% ============================================================
\newcommand{\poplogoinline}[1]{{\poplogo\fontsize{#1}{#1}\selectfont\color{popink}OnBuch\color{poporange}+}}
\newcommand{\signatureonbuch}{%
  \begin{tcolorbox}[enhanced,colback=popink,colframe=popink,boxrule=2.2pt,arc=10pt,
      drop shadow=poporange,left=14pt,right=14pt,top=10pt,bottom=10pt,before skip=0pt,after skip=0pt]
    \tikz[baseline=-0.7ex]\node[circle,fill=popgreen,draw=popcream,line width=1.3pt,minimum size=22pt,inner sep=0]{\popcheck{white}};%
    \hspace{9pt}\begin{minipage}[c]{0.62\linewidth}
      {\popxbold\fontsize{12}{13}\selectfont\color{popcream}Signé\ \textbullet\ {\poplogo OnBuch\color{poporange}+}}\par\vspace{2pt}
      {\raggedright\popsemi\fontsize{8}{10.5}\selectfont\color{popline}\MakeUppercase{Équipe pédagogique OnBuch+}\\\MakeUppercase{Conforme au programme officiel MINESEC}\par}
    \end{minipage}\hfill
    {\poplogo\fontsize{22}{22}\selectfont\color{popcream}OnBuch\color{poporange}+}
  \end{tcolorbox}%
}

% ============================================================
% Page de garde
% #1 titre · #2 matière · #3 niveau · #4 module · #5 n° leçon · #6 durée ·
% #7 description / objectifs (texte ou itemize)
% ============================================================
\newcommand{\pagegarde}[7]{%
  \thispagestyle{empty}
  \newgeometry{a4paper,margin=1.7cm,top=1.6cm,bottom=1.4cm}%
  \begin{tikzpicture}[remember picture,overlay]
    \fill[poporange!18] ([xshift=-3.2cm,yshift=-3.6cm]current page.north east) circle (4.6cm);
    \fill[poppurple!14] ([xshift=2.4cm,yshift=7.6cm]current page.south west) circle (3.8cm);
  \end{tikzpicture}%
  {\poplogo\fontsize{30}{30}\selectfont\color{popink}OnBuch\color{poporange}+}\hfill
  \raisebox{0.6ex}{\poppill{Cours signé \textbullet\ Premium}{popink}{popcream}}\par
  \vspace{1pt}\popsquiggle{poppurple}\par
  \vspace{8pt}
  \begin{tcolorbox}[enhanced,colback=poporange,colframe=popink,boxrule=2.8pt,arc=13pt,
      drop shadow=popink,left=18pt,right=18pt,top=12pt,bottom=13pt,after skip=10pt]
    \poppill{#2 \textbullet\ #3}{popink}{popcream}\hspace{5pt}\poppill{#4}{white}{popink}\par\vspace{8pt}
    {\popdisplay\fontsize{14}{14}\selectfont\color{popink}LEÇON #5}\par\vspace{4pt}
    {\raggedright\hyphenpenalty=10000\exhyphenpenalty=10000\popdisplay\fontsize{25}{27}\selectfont\color{popcream}\MakeUppercase{#1}\par}
    \vspace{8pt}
    {\popxbold\fontsize{9.5}{9.5}\selectfont\color{popink}\MakeUppercase{Durée conseillée : #6}}
  \end{tcolorbox}
  \begin{tcolorbox}[enhanced,colback=white,colframe=popink,boxrule=2.2pt,arc=10pt,
      drop shadow=popink,left=16pt,right=16pt,top=10pt,bottom=10pt,after skip=8pt]
    \poppill{Dans cette leçon}{poppurple}{white}\par\vspace{5pt}
    \popsemi\fontsize{9.3}{12.6}\selectfont\color{popink}#7
  \end{tcolorbox}
  \noindent\begin{minipage}[t]{0.487\linewidth}
    \begin{tcolorbox}[enhanced,colback=popgreen,colframe=popink,boxrule=2.2pt,arc=10pt,
        drop shadow=popink,left=11pt,right=11pt,top=10pt,bottom=10pt,height=3.1cm,valign=center]
      \tcbox[colback=white,colframe=popink,boxrule=1.3pt,arc=5pt,boxsep=0pt,nobeforeafter,
        left=3pt,right=3pt,top=3pt,bottom=3pt,valign=center]{\qrcode[height=1.9cm]{https://play.google.com/store/apps/details?id=com.luvvix.onbuch&hl=fr}}%
      \hspace{8pt}\begin{minipage}[c]{0.5\linewidth}
        {\popdisplay\fontsize{11}{12.5}\selectfont\color{white}\MakeUppercase{Télécharge OnBuch}}\par\vspace{4pt}
        {\raggedright\popsemi\fontsize{8.4}{11}\selectfont\color{white}Gratuit \textbullet\ Google Play\\Tuteur IA, annales, concours, résultats\par}
      \end{minipage}
    \end{tcolorbox}
  \end{minipage}\hfill
  \begin{minipage}[t]{0.487\linewidth}
    \begin{tcolorbox}[enhanced,colback=poppurple,colframe=popink,boxrule=2.2pt,arc=10pt,
        drop shadow=popink,left=11pt,right=11pt,top=10pt,bottom=10pt,height=3.1cm,valign=center]
      \tikz[baseline=-0.7ex]\node[circle,fill=white,draw=popink,line width=1.3pt,minimum size=1.6cm,inner sep=0]{\popdisplay\fontsize{24}{24}\selectfont\color{poppurple}+};%
      \hspace{8pt}\begin{minipage}[c]{0.6\linewidth}
        {\popdisplay\fontsize{11}{12.5}\selectfont\color{white}\MakeUppercase{Passe à OnBuch+}}\par\vspace{4pt}
        {\raggedright\popsemi\fontsize{8.4}{11}\selectfont\color{white}Tous les cours signés, Tuteur IA illimité, fiches PDF — onglet OnBuch+ de l'app\par}
      \end{minipage}
    \end{tcolorbox}
  \end{minipage}
  \vspace{8pt}
  \popcontenu
  \vfill
  \signatureonbuch
  \restoregeometry
  \newpage
}

% Rangée « Ce cours contient » (page de garde)
\newcommand{\poptile}[3]{% #1 couleur #2 gros texte #3 légende
  \begin{tcolorbox}[enhanced,colback=#1,colframe=popink,boxrule=2pt,arc=9pt,shadow={2.5pt}{-2.5pt}{0pt}{popink},
      left=6pt,right=6pt,top=9pt,bottom=9pt,halign=center,valign=center,height=1.75cm,width=\linewidth,nobeforeafter]
    {\popdisplay\fontsize{12.5}{13}\selectfont\color{white}\MakeUppercase{#2}}\par\vspace{3pt}
    {\centering\popsemi\fontsize{7.8}{9.5}\selectfont\color{white}#3\par}
  \end{tcolorbox}}
\newcommand{\popcontenu}{%
  \par\noindent{\popxboldls\fontsize{7.5}{7.5}\selectfont\color{popmuted}CE COURS SIGNÉ CONTIENT}\par\vspace{7pt}
  \noindent\begin{minipage}[t]{0.235\linewidth}\poptile{popblue}{Cours}{complet, illustré, pas à pas}\end{minipage}\hfill
  \begin{minipage}[t]{0.235\linewidth}\poptile{poporange}{Méthodes}{et exercices résolus}\end{minipage}\hfill
  \begin{minipage}[t]{0.235\linewidth}\poptile{poppink}{Exercices}{type examen + corrigés}\end{minipage}\hfill
  \begin{minipage}[t]{0.235\linewidth}\poptile{popgreen}{Fiche bilan}{l'essentiel à retenir}\end{minipage}\par}

% Sommaire
\newtcolorbox{sommaire}{enhanced,colback=white,colframe=popink,boxrule=2.2pt,arc=10pt,
  drop shadow=popink,left=18pt,right=18pt,top=13pt,bottom=13pt,before skip=6pt,after skip=20pt,
  before upper={\poppill{Sommaire}{poppurple}{white}\par\vspace{9pt}\popsemi\fontsize{10.6}{14.5}\selectfont\color{popink}}}

% Signature de fin de cours — poussée tout en bas de la dernière page
\newcommand{\findecours}{%
  \par\vfill\nopagebreak
  \begin{center}\popsquiggle{poporange}\end{center}\vspace{4pt}
  \signatureonbuch
}
\AtBeginDocument{\def\llbracket{\mathopen{[\mkern-3.5mu[}}\def\rrbracket{\mathclose{]\mkern-3.5mu]}}} % intervalles d'entiers (glyphe ⟦ absent de la police)
% Glyphes absents de Fira Math : redéfinis à partir de glyphes disponibles
\AtBeginDocument{%
  \def\setminus{\mathbin{\backslash}}\let\smallsetminus\setminus
  \def\vdots{\mathord{\vbox{\baselineskip3.2pt\lineskiplimit0pt\kern4pt\hbox{.}\hbox{.}\hbox{.}}}}%
  \def\ddots{\mathinner{\mkern1mu\raise6.5pt\hbox{.}\mkern2mu\raise3.6pt\hbox{.}\mkern2mu\raise0.7pt\hbox{.}\mkern1mu}}%
  \def\triangle{\mathord{\scalebox{0.9}{$\symup{\Delta}$}}}%
  \def\nabla{\mathord{\raisebox{\depth}{\scalebox{1}[-1]{$\symup{\Delta}$}}}}%
  \def\checkmark{\popcheck{popdark}}%
  \def\star{\mathbin{\ast}}\def\bigstar{\mathord{\ast}}%
  \def\diamond{\mathbin{\scalebox{0.6}{$\Diamond$}}}%
  \def\bigcup{\mathop{\mathchoice{\raisebox{-0.3em}{\scalebox{1.7}{$\cup$}}}{\scalebox{1.25}{$\cup$}}{\cup}{\cup}}\displaylimits}%
  \def\bigcap{\mathop{\mathchoice{\raisebox{-0.3em}{\scalebox{1.7}{$\cap$}}}{\scalebox{1.25}{$\cap$}}{\cap}{\cap}}\displaylimits}%
  \def\lesseqqgtr{\mathrel{\lessgtr}}\def\bigoplus{\mathop{\oplus}\displaylimits}%
}
\providecommand{\ssi}{si et seulement si} % abréviation fréquente
\AtBeginDocument{\providecommand{\wideparen}{\overparen}} % arc de cercle

%% FIGFIX-BEGIN v4 — correctifs globaux des figures (docs/onbuch-generation/tools/figfix)
\usepackage{adjustbox}\usepackage{etoolbox}
% toute figure TikZ/pgfplots plus large que la ligne est réduite à la largeur disponible
\BeforeBeginEnvironment{tikzpicture}{\begin{adjustbox}{max width=\linewidth}}
\AfterEndEnvironment{tikzpicture}{\end{adjustbox}}
% légendes pgfplots lisibles par-dessus les courbes
\pgfplotsset{every axis/.append style={legend style={fill=white,fill opacity=0.92,text opacity=1,draw=black!15,inner sep=3pt}}}
% étiquettes d'arêtes (syntaxe edge["..."]) avec fond blanc
\tikzset{every edge quotes/.append style={fill=white,inner sep=1.2pt}}
%% FIGFIX-END
\begin{document}

\pagegarde{VI : rédiger un CV}{Français}{Tle Tech}{Français – classe de Terminale technique}{N18}{9 séances (fiche de progression harmonisée)}{Cette leçon apprend à lire, analyser et rédiger un curriculum vitae (CV), texte fonctionnel indispensable à l'insertion professionnelle. Elle mobilise la maîtrise des structures externe et interne du CV, des valeurs du subjonctif et du conditionnel, et aboutit à la production et à la présentation d'un CV adapté à une offre d'embauche réelle.
\vspace{4pt}
\begin{multicols}{2}\small
\begin{itemize}
\item À la fin de cette leçon, je sais identifier un CV parmi d'autres textes fonctionnels et décrire sa fonction sociale.
\item À la fin de cette leçon, je sais décrire la structure externe d'un CV (mise en page, rubriques, photographie).
\item À la fin de cette leçon, je sais analyser la structure interne d'un CV (état civil, formation, expérience, compétences, centres d'intérêt).
\item À la fin de cette leçon, je sais interpréter les éléments iconiques d'un CV (photo, couleurs, typographie, pictogrammes).
\item À la fin de cette leçon, je sais employer correctement le subjonctif et le conditionnel dans les formulations d'un CV et d'une candidature.
\item À la fin de cette leçon, je sais rédiger un CV complet, clair et sans faute.
\end{itemize}\end{multicols}}

\begin{sommaire}
\begin{multicols}{2}
\begin{enumerate}
\item Le CV : un texte fonctionnel
\item La structure externe du CV
\item La structure interne du CV
\item Langue : les valeurs du subjonctif et du conditionnel
\item Rédiger et adapter un CV
\item Intégration, évaluation et remédiation
\item Activité d'intégration
\item Méthodes et exercices résolus
\item Exercices
\item Corrigés des exercices
\item Fiche bilan
\end{enumerate}
\end{multicols}
\end{sommaire}

\begin{objectifs}
\begin{itemize}
\item À la fin de cette leçon, je sais identifier un CV parmi d'autres textes fonctionnels et décrire sa fonction sociale.
\item À la fin de cette leçon, je sais décrire la structure externe d'un CV (mise en page, rubriques, photographie).
\item À la fin de cette leçon, je sais analyser la structure interne d'un CV (état civil, formation, expérience, compétences, centres d'intérêt).
\item À la fin de cette leçon, je sais interpréter les éléments iconiques d'un CV (photo, couleurs, typographie, pictogrammes).
\item À la fin de cette leçon, je sais employer correctement le subjonctif et le conditionnel dans les formulations d'un CV et d'une candidature.
\item À la fin de cette leçon, je sais rédiger un CV complet, clair et sans faute.
\item À la fin de cette leçon, je sais adapter un CV à une situation d'embauche donnée (offre d'emploi, stage, concours).
\item À la fin de cette leçon, je sais présenter oralement mon CV et justifier mes choix de rédaction.
\end{itemize}
\end{objectifs}

\begin{prerequis}
\begin{itemize}
\item Les textes fonctionnels étudiés en classe de Première (lettre administrative, lettre de motivation).
\item Les modes verbaux : indicatif, impératif, infinitif (notions de base du subjonctif et du conditionnel vues en 1ère).
\item La mise en page d'un document : titre, paragraphes, listes.
\item La lecture de documents iconiques (photographies, schémas) en réception des textes.
\item L'expression orale de présentation personnelle (se présenter en classe).
\end{itemize}
\end{prerequis}

\newpage

\coursec{Le CV : un texte fonctionnel}

\courssub{Situation d'accroche et problématique}

À Douala, une entreprise de BTP publie une offre d'emploi de technicien en génie civil. En une semaine, 350 candidatures arrivent sur le bureau du recruteur, qui ne retient que 10 CV pour les entretiens. Parmi les candidats écartés figurent pourtant de bons techniciens dont les CV étaient mal structurés, truffés de fautes ou inadaptés au poste. Comment rédiger un CV clair, convaincant et adapté à une situation d'embauche précise ? C'est tout l'enjeu de cette leçon.

Cette situation n'a rien d'exceptionnel. Chaque jour, à Yaoundé, à Bafoussam, à Garoua ou à Buea, des jeunes techniciens diplômés voient leur candidature écartée non pas parce qu'ils sont incompétents, mais parce que leur CV ne dit pas clairement qui ils sont et ce qu'ils savent faire. Le CV est donc un \cle{texte à enjeu} : de sa qualité dépend souvent l'accès à un entretien, à un stage ou à une formation. Avant d'apprendre à le rédiger, il faut comprendre sa nature : c'est un \cle{texte fonctionnel}.

\courssub{Qu'est-ce qu'un texte fonctionnel ?}

Commençons par l'intuition. Dans la vie quotidienne, nous ne lisons pas seulement des romans ou des poèmes. Nous remplissons des formulaires à la mairie, nous lisons des factures d'électricité, nous répondons à des petites annonces, nous rédigeons des demandes d'emploi. Ces textes ont un point commun : ils ne cherchent ni à émouvoir ni à raconter une histoire, mais à \emph{faire faire quelque chose} ou à \emph{transmettre une information utile} de manière efficace.

\begin{definition}[Texte fonctionnel]
Un \cle{texte fonctionnel} (ou texte utilitaire) est un texte écrit dans un but pratique précis : informer, demander, attester, vendre, recruter. Il obéit à des conventions de présentation fixes (rubriques, en-têtes, formules codifiées), privilégie la clarté et la concision, et vise une efficacité immédiate auprès d'un destinataire identifié.
\end{definition}

On distingue les textes fonctionnels des textes littéraires par plusieurs critères :

\begin{center}
\begin{tabularx}{\linewidth}{|l|X|X|}
\hline
\rowcolor{popblueL} \textbf{Critère} & \textbf{Texte fonctionnel} & \textbf{Texte littéraire} \\
\hline
But visé & Agir, informer, obtenir & Plaire, émouvoir, faire réfléchir \\
\hline
Forme & Codifiée, rubriques fixes & Libre, créative \\
\hline
Langue & Claire, directe, précise & Imagée, riche, ambiguïté possible \\
\hline
Destinataire & Identifié et précis & Lecteur indéterminé \\
\hline
Réussite & Efficacité (le but est atteint) & Valeur esthétique \\
\hline
\end{tabularx}
\end{center}

\courssub{Définition du curriculum vitae}

\begin{definition}[Curriculum vitae (CV)]
Le \cle{curriculum vitae} (expression latine signifiant littéralement « déroulement de la vie »), abrégé en \cle{CV}, est un document écrit qui récapitule, de manière structurée et synthétique, le parcours d'une personne : son identité, sa formation, ses diplômes, son expérience professionnelle, ses compétences et ses centres d'intérêt. Il est rédigé en vue d'obtenir un emploi, un stage ou une admission dans une formation.
\end{definition}

Le CV possède trois caractéristiques essentielles qu'il faut retenir :

\begin{itemize}
\item C'est un texte \cle{récapitulatif} : il résume un parcours, il ne le raconte pas. On n'y trouve pas de phrases longues ni de récit, mais des rubriques et des informations brèves.
\item C'est un texte \cle{structuré} : il obéit à une organisation en rubriques que le lecteur reconnaît immédiatement.
\item C'est un texte \cle{finalisé} : il est toujours écrit pour un destinataire précis et dans un but précis (un poste, un stage, une formation).
\end{itemize}

\courssub{La fonction sociale du CV}

Pourquoi rédige-t-on un CV ? Trois fonctions sociales se combinent.

\textbf{1. Se présenter.} Le CV est la carte d'identité professionnelle du candidat. Avant même de rencontrer la personne, le recruteur se fait une image d'elle à travers ce document : son âge, son niveau d'études, son expérience, ses compétences. Le CV parle donc \emph{à la place} du candidat, en son absence.

\textbf{2. Se valoriser.} Le CV est aussi un outil de mise en valeur. Il ne s'agit pas de mentir, mais de choisir et d'organiser les informations de manière à faire apparaître ses points forts : un stage réussi, une certification, une langue maîtrisée, un logiciel connu. Un même parcours peut paraître pauvre ou riche selon la manière dont il est présenté.

\textbf{3. Obtenir un entretien.} C'est la fonction ultime : le CV ne donne pas le poste, il donne \emph{accès à l'entretien}. Le recruteur qui lit un CV se pose une seule question : « Ce candidat mérite-t-il que je le rencontre ? » Tout le travail de rédaction vise donc à provoquer une réponse positive.

\begin{savaistu}[Moins d'une minute pour convaincre]
Des études menées auprès de recruteurs montrent qu'un CV est lu en moyenne \textbf{moins d'une minute} lors de la première sélection, parfois seulement 30 à 40 secondes. Le recruteur balaie le document du regard : identité, diplôme le plus récent, dernière expérience. Si ces informations ne sont pas immédiatement visibles, le CV est écarté. C'est pourquoi la clarté de la présentation et la hiérarchie des rubriques sont décisives : un CV n'est pas seulement lu, il est d'abord \emph{scanné}.
\end{savaistu}

\courssub{Le CV parmi les textes fonctionnels}

Le CV appartient à une grande famille de textes utilitaires. Il est important de le distinguer de ses voisins, avec lesquels on le confond parfois.

\begin{center}
\begin{tabularx}{\linewidth}{|l|X|X|}
\hline
\rowcolor{popblueL} \textbf{Texte} & \textbf{Fonction} & \textbf{Traits distinctifs} \\
\hline
CV & Résumer un parcours pour obtenir un emploi, un stage ou une formation & Rubriques, style télégraphique, pas de formules de politesse développées \\
\hline
Lettre de motivation & Argumenter, expliquer pourquoi on postule & Forme épistolaire, phrases rédigées, formules de politesse \\
\hline
Facture & Attester un achat et réclamer un paiement & Montants, références, mentions légales \\
\hline
Formulaire & Recueillir des informations normalisées & Cases à remplir, questions fermées \\
\hline
Annonce (offre d'emploi) & Proposer un bien, un service ou un poste & Brève, accrocheuse, coordonnées de contact \\
\hline
\end{tabularx}
\end{center}

Le critère de distinction fondamental est la \cle{fonction} : chaque texte fonctionnel répond à un besoin social précis, et sa forme découle de cette fonction. La facture chiffre, le formulaire interroge, l'annonce propose, le CV résume un parcours.

\begin{attention}[Ne pas confondre le CV et la lettre de motivation]
Erreur fréquente à l'examen comme dans la vie : rédiger le CV comme une lettre, ou la lettre comme un CV. Retiens la formule : \textbf{le CV décrit, la lettre argumente}. Le CV présente des faits (diplômes, dates, postes occupés) sous forme de rubriques et de courtes mentions, sans phrases de politesse. La lettre de motivation, elle, est un texte rédigé en phrases complètes qui explique \emph{pourquoi} le candidat postule et \emph{pourquoi} il est fait pour le poste. Les deux documents sont complémentaires : on envoie presque toujours les deux ensemble, mais ils ne se remplacent jamais l'un l'autre.
\end{attention}

\courssub{Les destinataires du CV}

Un texte fonctionnel s'écrit toujours pour quelqu'un. Le CV a quatre grandes catégories de destinataires :

\begin{itemize}
\item \textbf{Le recruteur d'entreprise} : chef d'entreprise, responsable du personnel, qui cherche un profil précis pour un poste précis.
\item \textbf{Le directeur des ressources humaines (DRH)} : dans les grandes entreprises et les administrations, il gère les candidatures et organise les sélections.
\item \textbf{Le maître de stage} : l'élève de Terminale technique qui cherche un stage académique adresse son CV au responsable de l'entreprise ou du service qui l'accueillera.
\item \textbf{Le jury de concours ou de formation} : certaines écoles et certains concours exigent un CV pour évaluer le parcours du candidat avant les épreuves.
\end{itemize}

Chaque destinataire attend des informations légèrement différentes : le maître de stage s'intéresse surtout à la formation et aux compétences techniques ; le recruteur d'entreprise regarde d'abord l'expérience. C'est pourquoi un bon CV est toujours \emph{adapté} à son destinataire, comme nous le verrons dans la suite de la leçon.

\courssub{Les qualités attendues d'un CV}

Quatre qualités font la valeur d'un CV.

\textbf{1. La clarté.} Le recruteur doit trouver chaque information en quelques secondes. Cela suppose des rubriques visibles, une mise en page aérée, une police lisible, un alignement soigné.

\textbf{2. La concision.} Un CV tient sur \cle{une à deux pages} maximum. Un élève de Terminale ou un jeune diplômé rédige un CV d'\emph{une seule page}. Tout ce qui n'est pas utile au poste visé doit être supprimé.

\textbf{3. L'exactitude.} Dates exactes, intitulés exacts des diplômes et des postes, coordonnées exactes (un numéro de téléphone faux rend le CV inutile). L'exactitude inclut l'orthographe : une faute dans un CV est souvent éliminatoire, car elle suggère la négligence.

\textbf{4. L'honnêteté.} On ne ment jamais dans un CV : un diplôme inventé ou une expérience exagérée sera découvert à l'entretien ou à l'essai, et détruira la crédibilité du candidat. Se valoriser, oui ; tricher, non.

\begin{exemplebox}[Une offre d'emploi et ce qu'elle exige du CV]
Lisons cette offre publiée à Yaoundé : « \emph{Entreprise CAMCONSULT recrute un technicien comptable. Profil : BTS ou DUT en comptabilité ; maîtrise de Sage Saari et d'Excel ; au moins un stage en cabinet comptable ; rigueur et sens de l'organisation. Envoyer CV et lettre de motivation avant le 30 novembre.} »

Cette annonce implique des attentes précises pour le CV : la rubrique « Formation » doit faire apparaître le BTS ou le DUT ; la rubrique « Compétences » doit mentionner Sage Saari et Excel ; la rubrique « Expérience » doit signaler le stage en cabinet. Un CV qui ne montre pas immédiatement ces trois éléments sera écarté, même si le candidat possède le profil. L'offre d'emploi est donc le \emph{mode d'emploi} du CV : on rédige son CV en répondant point par point à ce que l'annonce demande.
\end{exemplebox}

\courssub{Observation guidée : deux CV camerounais comparés}

Observons maintenant deux CV rédigés pour répondre à une offre de technicien en maintenance industrielle à Douala. Le premier, celui d'Arnaud M., est réussi ; le second, celui de Jean-Paul N., est défaillant.

Le CV d'Arnaud tient sur une page. Ses rubriques sont clairement titrées : État civil, Formation, Expérience professionnelle, Compétences, Langues, Centres d'intérêt. Chaque diplôme est accompagné de sa date et de son établissement (« BTS Maintenance industrielle, 2024, Lycée technique de Douala-Bassa »). Aucune faute d'orthographe. Le numéro de téléphone et l'adresse électronique sont exacts.

Le CV de Jean-Paul fait trois pages. Les rubriques se mélangent : le stage est mentionné au milieu de la formation, les dates sont absentes ou incohérentes (« 2020-2019 »). On relève huit fautes (« diplome », « maintennance », « j'ai travaillais »). La photo est floue, l'adresse électronique est fantaisiste et le paragraphe d'introduction raconte sa vie en phrases longues, comme une lettre. Le recruteur ne trouve pas le diplôme en moins d'une minute : le CV est écarté.

\begin{popfigure}
\begin{center}
\begin{tikzpicture}
\node[draw=popgreen, fill=popgreenL, rounded corners, thick, text width=0.44\linewidth, inner sep=8pt] (ok) at (0,0) {
\textbf{\textcolor{popdark}{CV réussi (Arnaud M.)}}\\[4pt]
\textbullet\ Lisibilité : mise en page aérée, rubriques visibles\\
\textbullet\ Rubriques : complètes et bien ordonnées\\
\textbullet\ Fautes : aucune\\
\textbullet\ Longueur : 1 page\\
\textbullet\ Dates : exactes et cohérentes\\
\textbullet\ Adaptation : compétences en lien avec le poste
};
\node[draw=poppink, fill=poppinkL, rounded corners, thick, text width=0.44\linewidth, inner sep=8pt, right=0.6cm of ok] (ko) {
\textbf{\textcolor{popdark}{CV défaillant (Jean-Paul N.)}}\\[4pt]
\textbullet\ Lisibilité : texte dense, paragraphes confus\\
\textbullet\ Rubriques : incomplètes, mélangées\\
\textbullet\ Fautes : 8 fautes d'orthographe\\
\textbullet\ Longueur : 3 pages\\
\textbullet\ Dates : absentes ou incohérentes\\
\textbullet\ Adaptation : informations hors sujet
};
\node[above=4pt of ok.north, font=\bfseries\small, text=popgreen] {Retenu pour l'entretien};
\node[above=4pt of ko.north, font=\bfseries\small, text=poppink] {Écarté dès la première lecture};
\end{tikzpicture}
\end{center}
\legende{Tableau comparatif des deux CV observés : à gauche, le CV retenu ; à droite, le CV écarté. Les mêmes critères (lisibilité, rubriques, fautes, longueur, dates, adaptation) expliquent le sort différent des deux candidatures.}
\end{popfigure}

Cette observation nous permet de dégager une conclusion générale : le sort d'un CV se joue sur des critères objectifs et contrôlables. Ce n'est pas la chance qui fait retenir un CV, c'est le respect des qualités attendues : clarté, concision, exactitude, honnêteté, adaptation au poste.

\begin{exoresolu}[Identifier la fonction d'un texte]
Énoncé. Voici trois extraits : (a) « BTS Génie civil, 2023, Lycée technique de Yaoundé ; stage de 3 mois, entreprise SABC, 2022 » ; (b) « Madame, Monsieur, vivement intéressé par le poste de technicien que vous proposez, je me permets de vous adresser ma candidature... » ; (c) « Montant dû : 45 000 FCFA, à régler avant le 15 du mois ». Identifie la nature de chaque texte et justifie ta réponse.
\tcblower
Solution détaillée. (a) Il s'agit d'un extrait de \textbf{CV} : les informations sont présentées sous forme de rubriques (formation, expérience), en style télégraphique, sans phrases rédigées ni formules de politesse ; le texte décrit un parcours. (b) Il s'agit d'une \textbf{lettre de motivation} : forme épistolaire (formule d'appel « Madame, Monsieur »), phrases complètes, démarche argumentative (« vivement intéressé... je me permets ») ; le texte argumente au lieu de décrire. (c) Il s'agit d'une \textbf{facture} : mention d'un montant et d'une échéance de paiement. Le critère décisif est donc la fonction : décrire un parcours (CV), argumenter (lettre), réclamer un paiement (facture).
\end{exoresolu}

\begin{exercice}[À toi de jouer]
Classe les textes suivants dans la catégorie qui convient (CV, lettre de motivation, annonce, formulaire, facture) et justifie chaque choix en une phrase : (1) « Nom : \ldots\ldots\ ; Date de naissance : \ldots\ldots\ ; Cocher la case correspondant à votre situation » ; (2) « Recherche secrétaire bilingue, expérience exigée, contacter le 6 99 00 00 00 » ; (3) « Expérience : technicien stagiaire, SONEL Douala, juillet-septembre 2023 » ; (4) « Total TTC : 120 000 FCFA » ; (5) « Je vous serais reconnaissant de bien vouloir examiner ma candidature ».
\end{exercice}

\begin{corrige}[Exercice : À toi de jouer]
(1) \textbf{Formulaire} : cases à remplir et informations normalisées. (2) \textbf{Annonce} : texte bref qui propose un poste avec un contact. (3) \textbf{CV} : rubrique « Expérience » en style télégraphique, avec dates et lieu. (4) \textbf{Facture} : mention d'un montant à payer. (5) \textbf{Lettre de motivation} : phrase rédigée, formule de politesse, démarche argumentative. Le critère de classement est toujours la fonction du texte.
\end{corrige}

\begin{aretenir}[L'essentiel de la section]
\begin{itemize}
\item Le CV est un \cle{texte fonctionnel} : il récapitule un parcours en vue d'un emploi, d'un stage ou d'une formation.
\item Ses trois fonctions sociales : se présenter, se valoriser, obtenir un entretien.
\item Le CV \emph{décrit} ; la lettre de motivation \emph{argumente} : ne jamais les confondre.
\item Ses destinataires : recruteurs, DRH, maîtres de stage, jurys.
\item Quatre qualités exigées : clarté, concision (1 à 2 pages), exactitude, honnêteté.
\item Un CV se rédige toujours en réponse à une offre ou à une situation d'embauche précise.
\end{itemize}
\end{aretenir}

\coursec{La structure externe du CV}

Après avoir compris \emph{ce qu'est} un CV — un texte fonctionnel, normé, destiné à convaincre en quelques secondes —, nous devons maintenant maîtriser \emph{comment il se présente}. La structure externe est l'écrin du contenu : avant même de lire un mot, le recruteur juge la rigueur, le sens de l'organisation et le professionnalisme du candidat à travers la mise en page. Une forme soignée ne garantit pas l'embauche, mais une forme négligée provoque souvent le rejet immédiat. Observons, analysons et codifions les règles de cette architecture visuelle.

\courssub{La mise en page générale : le cadre de la lisibilité}

\begin{definition}[Format, marges et aération]
La structure externe d'un CV repose sur un canevas technique strict : format \textbf{A4} (210 × 297 mm), marges symétriques de \textbf{2 à 2,5 cm} de chaque côté, interligne simple (1,15) à l'intérieur des blocs et espace vertical d'au moins \textbf{6 à 12 pt} entre deux rubriques. L'objectif est de créer un \emph{équilibre des masses} : ni page saturée (effet « mur de texte »), ni page vide (manque d'expérience). Le regard doit circuler naturellement de haut en bas, guidé par des \emph{blocs} visuels distincts.
\end{definition}

\begin{methode}[Organiser visuellement un CV en rubriques claires]
\begin{enumerate}
    \item \textbf{Diviser la page en zones logiques :} en-tête (identité), colonne de gauche (rubriques statiques : coordonnées, compétences, langues) et colonne de droite (rubriques dynamiques : formation, expériences, centres d'intérêt) \emph{ou} flux unique centré (une colonne).
    \item \textbf{Poser une grille invisible :} aligner tous les débuts de ligne sur un même axe vertical (marge gauche du corps de texte). Les dates, les titres de poste, les noms d'établissement doivent démarrer au même abscisse.
    \item \textbf{Aérer par le blanc :} réserver un espace vertical constant (ex. 8 pt) après chaque titre de rubrique et après chaque entrée (un poste, un diplôme). Le blanc n'est pas du vide, c'est un \emph{signe de ponctuation visuelle}.
    \item \textbf{Vérifier l'alignement :} utiliser les tabulations ou un tableau masqué (bordures nulles) pour garantir que les secondes lignes d'une même entrée repartent exactement sous la première.
\end{enumerate}
\end{methode}

\begin{popfigure}
\begin{tikzpicture}[scale=0.85, every node/.style={font=\small}]
    % Page A4 frame
    \draw[popdark, thick] (0,0) rectangle (18,25.4);
    % Marges
    \draw[popmuted, dashed] (1.5,1.5) rectangle (16.5,23.9);
    \node[popmuted, anchor=north west] at (1.6,23.8) {Marges 2 cm};
    % En-tête
    \fill[popblue!10] (1.5,21.5) rectangle (16.5,23.9);
    \draw[popblue, thick] (1.5,21.5) -- (16.5,21.5);
    \node[popblue, align=center, font=\bfseries] at (9,22.7) {EN-TÊTE : Nom, Prénom, Titre visé};
    \node[popblue, align=center] at (9,22.0) {Téléphone | Email | Ville | LinkedIn};
    \node[popblue, align=center] at (9,21.7) {\textit{(Photo : coin sup. droit ou gauche)}};
    % Colonne gauche (sidebar)
    \fill[popgreen!5] (1.5,1.5) rectangle (6.5,21.5);
    \draw[popgreen, thick] (6.5,1.5) -- (6.5,21.5);
    \node[popgreen, rotate=90, anchor=south] at (6.7,13) {COLONNE GAUCHE (Sidebar)};
    \node[popgreen, align=center, font=\bfseries] at (4,20) {COMPÉTENCES};
    \node[popgreen, align=center] at (4,18) {• Technique\\• Logiciels\\• Langues};
    \node[popgreen, align=center, font=\bfseries] at (4,15) {QUALITÉS};
    \node[popgreen, align=center] at (4,13) {• Rigueur\\• Esprit d'équipe};
    \node[popgreen, align=center, font=\bfseries] at (4,10) {CENTRES D'INTÉRÊT};
    \node[popgreen, align=center] at (4,8) {• Sport\\• Bénévolat};
    \node[popgreen, align=center, font=\bfseries] at (4,5) {RÉFÉRENCES};
    \node[popgreen, align=center] at (4,3) {Disponibles sur demande};
    % Colonne droite (corps principal)
    \node[poporange, align=center, font=\bfseries] at (11.5,20.5) {CORPS PRINCIPAL (Flux chronologique inverse)};
    \draw[poporange, thick] (6.5,20.2) -- (16.5,20.2);
    \node[poporange, font=\bfseries] at (7.5,19.5) {FORMATION};
    \node[poporange, anchor=west] at (7.5,19.0) {2023--2024 \hfill Baccalauréat Technique Génie Civil};
    \node[poporange, anchor=west] at (7.5,18.6) {2021--2023 \hfill Probatoire Série F4};
    \node[poporange, font=\bfseries] at (7.5,17.6) {EXPÉRIENCES / STAGES};
    \node[poporange, anchor=west] at (7.5,17.1) {Juil--Août 2023 \hfill Stagiaire BTP -- Entreprise X};
    \node[poporange, anchor=west] at (7.5,16.7) {\textit{Suivi chantier, DAO, métrés}};
    \node[poporange, anchor=west] at (7.5,16.2) {Juil--Août 2022 \hfill Ouvrier qualifié -- Chantier Y};
    \node[poporange, font=\bfseries] at (7.5,15.2) {PROJETS SCOLAIRES};
    \node[poporange, anchor=west] at (7.5,14.7) {2024 \hfill Maquette pont en béton armé};
    \node[poporange, anchor=west] at (7.5,14.3) {2023 \hfill Étude de sol pour bâtiment R+2};
    % Flèches d'annotation
    \draw[poppink, <->, thick] (0.5,21.5) -- (0.5,23.9) node[midway, left, poppink, font=\bfseries] {Zone En-tête};
    \draw[poppink, <->, thick] (0.5,1.5) -- (0.5,21.5) node[midway, left, poppink, font=\bfseries] {Zone Corps};
    \draw[poppink, <->, thick] (1.5,0.8) -- (6.5,0.8) node[midway, below, poppink, font=\bfseries] {Sidebar (≈30\%)};
    \draw[poppink, <->, thick] (6.5,0.8) -- (16.5,0.8) node[midway, below, poppink, font=\bfseries] {Corps principal (≈70\%)};
\end{tikzpicture}
\legende{Schéma annoté d'une structure externe type « deux colonnes » (sidebar + corps principal). Les proportions 30/70 sont indicatives ; l'essentiel est l'alignement vertical rigoureux des blocs.}
\end{popfigure}

\courssub{L'en-tête : l'identité professionnelle au premier coup d'œil}

L'en-tête occupe le tiers supérieur de la page. Il ne s'agit pas d'un simple « nom + téléphone », mais d'une \emph{signature visuelle}.

\begin{propriete}[Composition de l'en-tête optimal]
\begin{enumerate}
    \item \textbf{Nom et Prénom :} \textsc{NOM} en majuscules petites capitales (ou gras taille 18--22 pt), Prénom en titre (taille 14--16 pt). Exemple : \textsc{MBARGA} Étienne.
    \item \textbf{Titre visé / Accroche :} une ligne sous le nom : \emph{« Technicien supérieur en Génie Civil — Spécialité Bâtiment »} ou \emph{« Candidat au BTS Travaux Publics »}. Cela ancre immédiatement la candidature.
    \item \textbf{Coordonnées essentielles :} téléphone (format international +237 6XX XX XX XX), e-mail professionnel (\texttt{prenom.nom@domaine.com}), ville de résidence (ex. « Douala, Akwa »), lien LinkedIn/portfolio si à jour.
    \item \textbf{Exclure :} âge, état civil, nationalité, numéro CNI, photo d'identité non professionnelle (voir infra).
\end{enumerate}
\end{propriete}

\begin{attention}[Erreur fréquente : l'en-tête « passeport »]
Ne pas écrire : « Je m'appelle Étienne Mbarga, je suis né le 12/03/2005 à Yaoundé, je suis célibataire, mon numéro est... ». Le CV n'est pas une fiche d'état civil. L'en-tête doit être \textbf{compact, aligné, scannable en 3 secondes}. Utilisez des séparateurs discrets (\textbullet\ |\ \textbullet) ou des pictogrammes (☎ ✉ 📍) pour structurer l'information sur une ou deux lignes maximum.
\end{attention}

\courssub{Les rubriques : titres visibles, ordre logique, séparateurs}

\begin{definition}[Rubrique standardisée]
Une rubrique est un conteneur sémantique identifié par un \textbf{titre court, en majuscules, gras, taille 12--14 pt}, éventuellement précédé d'un pictogramme (🎓 Formation, 💼 Expériences, 🛠 Compétences, 🌐 Langues). L'ordre suit la \emph{pertinence pour le poste visé} :
\begin{itemize}
    \item Profil junior (étudiant, jeune diplômé) : Formation → Expériences/Stages → Projets → Compétences → Langues → Centres d'intérêt.
    \item Profil expérimenté : Expériences → Formation → Compétences → Certifications → Langues.
\end{itemize}
Un séparateur visuel (filet fin 0,5 pt, couleur popblue ou popgray) marque la transition entre rubriques.
\end{definition}

\begin{exemplebox}[Hiérarchie des titres dans le CV]
\begin{flushleft}
\texttt{FORMATION}                          (Niveau 1 : Rubrique principale)\\
\texttt{\hspace{1em}2023--2024  Baccalauréat F4}        (Niveau 2 : Entrée chronologique)\\
\texttt{\hspace{2em}Lycée Technique de Douala}        (Niveau 3 : Détail / Lieu)\\
\texttt{\hspace{3em}Spécialité Génie Civil}         (Niveau 4 : Précision option)\\
\texttt{\hspace{4em}Mention Assez Bien}           (Niveau 5 : Résultat)\\
\texttt{EXPÉRIENCES PROFESSIONNELLES}\\
\texttt{\hspace{1em}2023  Stage ouvrier, Entreprise X}\\
\texttt{\hspace{2em}- Coffrage, ferraillage, coulage}
\end{flushleft}
L'indentation (retrait) signale la subordination. La cohérence de cette hiérarchie sur toute la page prouve la rigueur du candidat.
\end{exemplebox}

\courssub{La photographie d'identité : texte iconique intégré}

\begin{savaistu}[La photo au Cameroun : usage et attentes]
Au Cameroun, comme dans la zone CEMAC et une grande partie de l'Afrique francophone, la photo d'identité \textbf{n'est pas obligatoire légalement} (lutte contre la discrimination), mais elle reste \textbf{très fortement attendue} par les recruteurs locaux (administrations, grandes entreprises, PME). Son absence peut être perçue comme un oubli ou un manque de transparence. En revanche, pour des candidatures internationales (France, Canada, ONG internationales), la tendance est à l'absence de photo. \textbf{Règle pratique :} si l'offre ne précise pas, mettez une photo professionnelle pour le marché camerounais ; retirez-la pour l'international.
\end{savaistu}

\begin{propriete}[Standards de la photo professionnelle]
\begin{itemize}
    \item \textbf{Format :} 3,5 × 4,5 cm (standard passeport), résolution 300 dpi.
    \item \textbf{Cadre :} fond uni clair (blanc, gris très clair, bleu pâle), cadrage visage/épaules, regard vers l'objectif.
    \item \textbf{Tenue :} adaptée au secteur (chemise/cravate pour tertiaire, tenue de chantier propre avec EPI pour BTP, blouse pour laboratoire).
    \item \textbf{Position :} coin supérieur droit (classique) ou gauche (moderne, aligné sur la colonne latérale). Jamais au centre de la page.
    \item \textbf{Qualité :} pas de selfie, pas de photo de groupe recadrée, pas de filtre Snapchat/Instagram. Faites-la faire par un photographe ou un camarade avec un bon appareil et un fond tendu.
\end{itemize}
\end{propriete}

\courssub{La typographie : police unique, tailles hiérarchisées, gras modéré}

\begin{definition}[Charte typographique du CV]
Le CV n'utilise \textbf{qu'une seule famille de police} (police à empattement \emph{serif} comme Garamond, Times New Roman, Georgia pour le classicisme ; ou sans empattement \emph{sans-serif} comme Calibri, Arial, Helvetica, Roboto, Open Sans pour la modernité technique). La hiérarchie s'obtient par :
\begin{itemize}
    \item \textbf{Taille :} Nom (18--22 pt) > Titres rubriques (12--14 pt, majuscules) > Corps de texte (10--11 pt) > Détails (9--10 pt, gris foncé).
    \item \textbf{Gras :} réservé aux titres de rubriques, noms d'entreprises/établissements, intitulés de poste/diplôme. \textbf{Jamais} sur le corps de texte courant.
    \item \textbf{Italique :} pour les précisions (spécialité, mention, lieu, dates) ou les verbes d'action en début de puce.
\end{itemize}
\end{definition}

\begin{attention}[Piège mortel : le « carnaval typographique »]
Multiplier les polices (Arial pour le nom, Times pour le corps, Comic Sans pour les centres d'intérêt...), varier les tailles au hasard, abuser du gras, du soulignement, des majuscules dans le corps de texte. \textbf{Conséquence immédiate :} le CV paraît amateur, brouillon, difficile à lire. Le recruteur fatigue en 10 secondes et passe au suivant. \textbf{Solution :} choisissez \emph{une} police, définissez \emph{trois} tailles (titre, rubrique, corps), \emph{deux} graisses (normal, gras), \emph{une} couleur de texte (noir ou gris anthracite \#333333). Verrouillez ce style.
\end{attention}

\courssub{Couleurs et pictogrammes : sobriété professionnelle}

\begin{methode}[Appliquer une identité visuelle cohérente]
\begin{enumerate}
    \item \textbf{Choisir une couleur d'accent unique :} bleu marine (\#1B3A5C), vert sapin (\#1A5C2E), bordeaux (\#7B1E3A), gris anthracite. Cette couleur sert pour : filets de séparation, titres de rubriques, pictogrammes, puces.
    \item \textbf{Limiter la palette :} fond blanc, texte noir/gris, \textbf{une} couleur d'accent. Pas de dégradés, pas d'ombres portées, pas de fonds colorés derrière le texte (sauf sidebar très claire, ex. \#F0F4F8).
    \item \textbf{Pictogrammes :} utiliser une bibliothèque cohérente (ex. FontAwesome, Material Icons) : 🎓 Formation, 💼 Expérience, 🛠 Compétences, 🌐 Langues, 📍 Lieu, ☎ Téléphone, ✉ Email. Taille 10--11 pt, alignés sur la ligne de base du texte.
    \item \textbf{Test d'impression :} imprimez en noir et blanc. La hiérarchie doit rester lisible (couleur d'accent → gris foncé). Si le contraste disparaît, changez la teinte.
\end{enumerate}
\end{methode}

\begin{popfigure}
\begin{tikzpicture}[scale=0.9, every node/.style={font=\small}]
    % Maquette 1 colonne
    \draw[popdark, thick] (0,0) rectangle (8,22);
    \draw[popmuted, dashed] (0.5,0.5) rectangle (7.5,21.5);
    \fill[popblue!10] (0.5,19) rectangle (7.5,21.5);
    \node[popblue, font=\bfseries, align=center] at (4,20.5) {EN-TÊTE};
    \node[popblue, align=center] at (4,19.5) {Nom | Titre | Coordonnées | Photo};
    \node[poporange, font=\bfseries] at (4,18.2) {FORMATION};
    \draw[poporange] (0.5,17.8) -- (7.5,17.8);
    \node[anchor=west] at (1,17.2) {2024 BAC F4 Génie Civil};
    \node[anchor=west] at (1,16.7) {2022 Probatoire F4};
    \node[poporange, font=\bfseries] at (4,15.5) {EXPÉRIENCES};
    \draw[poporange] (0.5,15.1) -- (7.5,15.1);
    \node[anchor=west] at (1,14.5) {2023 Stage Entreprise X};
    \node[anchor=west] at (1,14.0) {2022 Ouvrier Chantier Y};
    \node[poporange, font=\bfseries] at (4,12.8) {COMPÉTENCES};
    \draw[poporange] (0.5,12.4) -- (7.5,12.4);
    \node[anchor=west] at (1,11.8) {DAO (AutoCAD), Béton, Métrés};
    \node[anchor=west] at (1,11.3) {Français C1, Anglais B1};
    \node[poporange, font=\bfseries] at (4,10.1) {PROJETS / INTÉRÊTS};
    \draw[poporange] (0.5,9.7) -- (7.5,9.7);
    \node[anchor=west] at (1,9.1) {Maquette pont, Bénévolat Croix-Rouge};
    \node[popdark, font=\bfseries, align=center] at (4,-0.8) {MAQUETTE 1 COLONNE};
    \node[popmuted, align=center] at (4,-1.5) {Flux linéaire, idéal pour profils juniors,\\ATS-friendly, impression simple.};
    % Maquette 2 colonnes
    \begin{scope}[xshift=10cm]
        \draw[popdark, thick] (0,0) rectangle (8,22);
        \draw[popmuted, dashed] (0.5,0.5) rectangle (7.5,21.5);
        % Sidebar gauche
        \fill[popgreen!5] (0.5,0.5) rectangle (2.8,21.5);
        \draw[popgreen, thick] (2.8,0.5) -- (2.8,21.5);
        \node[popgreen, rotate=90, anchor=south] at (3.0,11) {SIDEBAR};
        \node[popgreen, font=\bfseries, align=center] at (1.65,20) {CONTACT};
        \node[popgreen, align=center] at (1.65,19) {☎ +237 6XX...};
        \node[popgreen, align=center] at (1.65,18.5) {✉ email...};
        \node[popgreen, align=center] at (1.65,18) {📍 Douala};
        \node[popgreen, font=\bfseries, align=center] at (1.65,16.5) {COMPÉTENCES};
        \node[popgreen, align=center] at (1.65,15.5) {AutoCAD, Revit};
        \node[popgreen, align=center] at (1.65,15) {Béton armé};
        \node[popgreen, font=\bfseries, align=center] at (1.65,13.5) {LANGUES};
        \node[popgreen, align=center] at (1.65,12.5) {Français ★★★★★};
        \node[popgreen, align=center] at (1.65,12) {Anglais ★★★☆☆};
        \node[popgreen, font=\bfseries, align=center] at (1.65,10.5) {INTÉRÊTS};
        \node[popgreen, align=center] at (1.65,9.5) {Football, Lecture};
        % Corps droit
        \fill[popblue!10] (3.0,19) rectangle (7.5,21.5);
        \node[popblue, font=\bfseries, align=center] at (5.25,20.5) {EN-TÊTE};
        \node[popblue, align=center] at (5.25,19.5) {Nom Prénom | Titre visé};
        \node[poporange, font=\bfseries] at (5.25,18.2) {FORMATION};
        \draw[poporange] (3.0,17.8) -- (7.5,17.8);
        \node[anchor=west] at (3.3,17.2) {2024 BAC F4 Génie Civil};
        \node[anchor=west] at (3.3,16.7) {Mention AB - Lycée Tech Douala};
        \node[poporange, font=\bfseries] at (5.25,15.5) {EXPÉRIENCES};
        \draw[poporange] (3.0,15.1) -- (7.5,15.1);
        \node[anchor=west] at (3.3,14.5) {2023 Stage Entreprise X};
        \node[anchor=west, font=\itshape] at (3.3,14.1) {Suivi chantier, DAO, Métrés};
        \node[anchor=west] at (3.3,13.5) {2022 Ouvrier Chantier Y};
        \node[anchor=west, font=\itshape] at (3.3,13.1) {Coffrage, Ferraillage};
        \node[poporange, font=\bfseries] at (5.25,11.8) {PROJETS SCOLAIRES};
        \draw[poporange] (3.0,11.4) -- (7.5,11.4);
        \node[anchor=west] at (3.3,10.8) {2024 Maquette pont béton};
        \node[anchor=west] at (3.3,10.3) {2023 Étude de sol R+2};
        \node[popdark, font=\bfseries, align=center] at (4,-0.8) {MAQUETTE 2 COLONNES};
        \node[popmuted, align=center] at (4,-1.5) {Densité info max, lecture en Z,\\idéal si beaucoup de compétences techniques.};
    \end{scope}
\end{tikzpicture}
\legende{Comparaison de deux mises en page pour un même profil (Élève Tle Génie Civil). La version 1 colonne (gauche) privilégie la lecture linéaire et la compatibilité ATS. La version 2 colonnes (droite) densifie l'information et met en avant les compétences techniques dès le premier regard.}
\end{popfigure}

\courssub{Lecture iconique : ce que la forme révèle du candidat}

Au-delà des règles techniques, la structure externe est un \emph{texte iconique} que le recruteur « lit » en quelques secondes. Cette lecture pré-attentive déclenche des inférences automatiques sur la personnalité professionnelle.

\begin{exoresolu}[Analyse comparative de deux mises en page]
\textbf{Contexte :} Deux élèves de Terminale Génie Civil, même parcours (BAC F4, stages identiques, même projet de fin d'études), postulent au même stage de fin d'études dans une entreprise de BTP à Douala.

\textbf{Candidat A (CV « Propre ») :}
\begin{itemize}
    \item Police unique (Calibri 11), titres rubriques Calibri 13 Gras Bleu Marine.
    \item Marge 2 cm, alignement parfait à gauche (tabulations), inter-blocs 8 pt constants.
    \item Photo pro fond bleu clair, coin sup. droit, 3,5×4,5 cm.
    \item Sidebar gauche : Compétences techniques (AutoCAD, Robot, Métrés) + Langues + Permis B.
    \item Corps droit : Formation, Expériences (verbes d'action alignés), Projet PFE.
    \item Imprimé sur papier 90g blanc, PDF généré depuis Word/LaTeX (pas d'image scannée).
\end{itemize}

\textbf{Candidat B (CV « Bricolé ») :}
\begin{itemize}
    \item Police : Nom en Comic Sans 24 rose, rubriques en Times 14 souligné, corps en Arial 10,5.
    \item Marges irrégulières (1,2 cm à gauche, 3 cm à droite), dates parfois à gauche, parfois centrées.
    \item Photo : selfie anniversaire recadré, fond cuisine, 5×5 cm, au milieu de l'en-tête.
    \item Pas de rubriques distinctes : tout en vrac avec des tirets « - ».
    \item Fautes de frappe : « Geni Civil », « AutoCad », « mettres ».
    \item Fond de page gris foncé (gaspillage encre), texte blanc par endroits, illisible en N\&B.
\end{itemize}

\textbf{Lecture iconique du recruteur (30 secondes) :}
\begin{center}
\begin{tabularx}{\linewidth}{|l|X|X|}\hline
\textbf{Critère visuel} & \textbf{Candidat A} & \textbf{Candidat B} \\ \hline
Rigueur / Soin & \textbf{Élevée} : alignements, cohérence & \textbf{Faible} : désalignements, incohérences \\ \hline
Compétence technique & \textbf{Supposée} : DAO mis en avant, structure claire & \textbf{Doutée} : fautes sur noms logiciels, présentation chaotique \\ \hline
Adaptabilité entreprise & \textbf{Forte} : codes pro respectés (photo, sobre) & \textbf{Faible} : codes ignorés (selfie, couleurs) \\ \hline
Respect procédure & \textbf{Oui} : PDF propre, ATS-ready & \textbf{Non} : format instable, fond coloré \\ \hline
Décision probable & \textbf{Appelé en entretien} & \textbf{Rejeté (poubelle)} \\ \hline
\end{tabularx}
\end{center}

\textbf{Conclusion :} La structure externe n'est pas du « décor ». Elle \emph{est} le premier message de compétence. Pour un technicien, la précision du tracé, le respect des normes de présentation, la lisibilité du plan sont des preuves \emph{in situ} de la qualité de son futur travail sur chantier ou en bureau d'études.
\end{exoresolu}

\begin{aretenir}[Check-list finale de la structure externe (à valider avant envoi)]
\begin{itemize}
    \item [ ] Format A4, marges 2 cm, police unique (Calibri/Roboto/Garamond), 3 tailles max.
    \item [ ] En-tête complet : NOM Prénom, Titre visé, Tél, Email pro, Ville, LinkedIn (option).
    \item [ ] Photo pro (3,5×4,5 cm, fond uni, tenue adaptée) OU absence assumée (candidature internationale).
    \item [ ] Rubriques standardisées (FORMATION, EXPÉRIENCES, COMPÉTENCES, LANGUES, PROJETS, INTÉRÊTS).
    \item [ ] Ordre chronologique inverse dans chaque rubrique (le plus récent en haut).
    \item [ ] Alignement vertical parfait (tabulations), indentation cohérente pour sous-niveaux.
    \item [ ] Une seule couleur d'accent (bleu marine/vert/bordeaux) pour titres, filets, icônes.
    \item [ ] Pictogrammes uniformes (même librairie, même taille).
    \item [ ] Test impression N\&B réussi (contraste conservé, pas de fond noir).
    \item [ ] Fichier final : \textbf{PDF} (nommé \texttt{Nom\_Prenom\_CV\_Poste.pdf}), pas de .docx, pas de .jpg.
\end{itemize}
\end{aretenir}

La structure externe est désormais maîtrisée. Elle pose le cadre rigoureux dans lequel viendra s'inscrire le contenu — la structure interne — que nous détaillerons à la section suivante. Un CV bien architecturé est déjà à moitié convaincu.

\coursec{La structure interne du CV}

Dans la section précédente, nous avons étudié la \cle{structure externe} du CV : sa présentation matérielle, sa mise en page, sa longueur, sa typographie. Mais une belle enveloppe ne suffit pas si le contenu est vide ou mal organisé. Le recruteur qui ouvre un CV cherche des informations précises, rangées dans des \cle{rubriques} claires. C'est ce que nous appelons la \cle{structure interne} du CV : l'ensemble des rubriques et des informations qui composent le document. Un CV efficace, c'est un CV dont chaque rubrique répond à une question que l'employeur se pose : Qui es-tu ? Que sais-tu faire ? Qu'as-tu déjà fait ? Peut-on te faire confiance ?

\courssub{Les rubriques du CV : vue d'ensemble}

\begin{definition}[Structure interne du CV]
La \cle{structure interne} d'un CV est l'organisation du contenu du document en \cle{rubriques} ordonnées : l'état civil, le titre ou l'objectif professionnel, la formation, l'expérience professionnelle, les compétences, les centres d'intérêt et, éventuellement, les références. Chaque rubrique regroupe des informations de même nature, présentées de manière brève et hiérarchisée.
\end{definition}

\begin{popfigure}
\begin{center}
\begin{tikzpicture}[
  grow=right,
  level 1/.style={sibling distance=1.55cm, level distance=3.6cm},
  level 2/.style={sibling distance=0.62cm, level distance=3.4cm},
  every node/.style={font=\small},
  edge from parent/.style={draw=popmuted, thick}
]
\node[draw=popblue, fill=popblueL, rounded corners, font=\bfseries] {CV}
  child { node[draw=popgreen, fill=popgreenL, rounded corners] {\textbf{État civil}}
    child { node {contacts} }
    child { node {nom, naissance} }
  }
  child { node[draw=poppink, fill=poppinkL, rounded corners] {\textbf{Objectif}}
    child { node {phrase ciblée} }
  }
  child { node[draw=popblue, fill=popblueL, rounded corners] {\textbf{Formation}}
    child { node {ordre antichronologique} }
    child { node {diplômes, années} }
  }
  child { node[draw=popgold, fill=popgoldL, rounded corners] {\textbf{Expérience}}
    child { node {missions} }
    child { node {poste, structure, dates} }
  }
  child { node[draw=poppurple, fill=poppurpleL, rounded corners] {\textbf{Compétences}}
    child { node {linguistiques} }
    child { node {techniques} }
  }
  child { node[draw=poporange, fill=poporangeL, rounded corners] {\textbf{Centres d'intérêt}}
    child { node {engagement} }
    child { node {sport, lecture} }
  }
  child { node[draw=popgreen, fill=popgreenL, rounded corners] {\textbf{Références}}
    child { node {personnes à contacter} }
  };
\end{tikzpicture}
\end{center}
\legende{Arbre de la structure interne d'un CV : sept rubriques principales dans l'ordre logique de lecture (haut $\rightarrow$ bas) et leurs sous-éléments essentiels.}
\end{popfigure}

\courssub{L'état civil ou informations personnelles}

C'est la première rubrique, placée en haut du document. Elle permet au recruteur d'\textbf{identifier} le candidat et de le \textbf{contacter} rapidement. Elle comprend :

\begin{itemize}
\item le \cle{nom} et les \cle{prénoms}, en toutes lettres (éventuellement en majuscules pour le nom) ;
\item la \cle{date} et le \cle{lieu de naissance} ;
\item la \cle{nationalité} ;
\item les \cle{contacts} : adresse postale, numéro de téléphone (9 chiffres au Cameroun), adresse électronique professionnelle ;
\item la \cle{situation matrimoniale} (célibataire, marié(e)) : elle est \textbf{facultative} et ne doit être mentionnée que si elle est demandée ou utile pour le poste.
\end{itemize}

\begin{exemplebox}[Rubrique « État civil »]
\textbf{MBALLA ESSOMBA Jean-Claude}\par
Né le 12 mars 2006 à Bafoussam\par
Nationalité camerounaise\par
B.P. 452 Dschang — Tél. : 6 90 00 00 00 — Courriel : jcmballa@mail.com
\end{exemplebox}

\begin{attention}[Piège : les informations inutiles ou dangereuses]
N'indique jamais dans l'état civil des informations sans rapport avec le poste : religion, ethnie, opinions politiques, numéro de CNI, groupe sanguin. Un numéro de téléphone faux ou une adresse électronique fantaisiste (du type « beaugossedouala@... ») discrédite immédiatement la candidature. Choisis une adresse électronique sobre, formée de ton nom et de ton prénom (ex. : prenom.nom@domaine.com).
\end{attention}

\courssub{Le titre ou l'objectif professionnel}

Sous l'état civil, le CV moderne comporte souvent un \cle{titre} (ou objectif professionnel) : une phrase courte qui annonce le poste recherché ou le projet professionnel. Cette phrase doit être \textbf{ciblée}, c'est-à-dire adaptée à l'offre d'emploi ou au secteur visé. Elle remplace l'ancienne formule « Madame, Monsieur, je vous adresse ma candidature... » qui figure dans la lettre de motivation, pas dans le CV.

\begin{exemplebox}[Exemples d'objectifs professionnels]
\begin{itemize}
\item « Jeune titulaire d'un Baccalauréat technique F3 (Construction mécanique), à la recherche d'un poste d'aide-comptable / technicien de maintenance. »
\item « Technicien en froid et climatisation (Bac F2), disponible immédiatement pour un CDD ou un stage. »
\item « À la recherche d'un stage académique de fin d'études en maintenance des engins de chantier (BTP). »
\end{itemize}
\end{exemplebox}

\courssub{La formation ou cursus scolaire et académique}

Cette rubrique présente les \cle{diplômes} obtenus, les \cle{établissements} fréquentés, les \cle{séries} ou \cle{spécialités} et les \cle{années} correspondantes. Elle obéit à une règle d'organisation fondamentale : l'\cle{ordre antichronologique}.

\begin{definition}[Ordre antichronologique]
L'\cle{ordre antichronologique} est une présentation qui va \textbf{du plus récent au plus ancien} : on cite d'abord le diplôme ou l'expérience le plus récent, puis on remonte le temps. C'est l'ordre exigé dans un CV, car le recruteur s'intéresse d'abord à ce que le candidat a fait dernièrement (son niveau actuel).
\end{definition}

\begin{exemplebox}[Rubrique « Formation » — exemple pour un élève de Terminale en 2026-2027]
\textbf{FORMATION}\par
\begin{itemize}
\item \textbf{2027 (prévu)} : Baccalauréat technique série F3 (Construction mécanique), Lycée technique de Douala-Bassa ;
\item \textbf{2026} : Probatoire technique F3, Lycée technique de Douala-Bassa ;
\item \textbf{2023} : BEPC, Collège bilingue de Bonabéri.
\end{itemize}
\end{exemplebox}

\begin{popfigure}
\begin{center}
\begin{tikzpicture}[scale=1]
\draw[->, very thick, popblue] (0,0) -- (11.5,0) node[right, font=\small\itshape] {temps chronologique};
\foreach \x/\annee in {1.2/2023, 4.6/2026, 8/2027} {
  \draw[thick, popdark] (\x,0.15) -- (\x,-0.15);
  \node[below, font=\small\bfseries] at (\x,-0.2) {\annee};
}
\node[draw=popgreen, fill=popgreenL, rounded corners, align=center, font=\small] at (1.2,1) {BEPC};
\node[draw=popgold, fill=popgoldL, rounded corners, align=center, font=\small] at (4.6,1.6) {Probatoire F3};
\node[draw=poporange, fill=poporangeL, rounded corners, align=center, font=\small] at (8,2.2) {Baccalauréat F3};
\draw[->, very thick, poppink] (9.6,3) -- (2.4,3) node[midway, above, font=\small\bfseries, poppink] {Ordre de lecture du CV : du plus récent (2027) au plus ancien (2023)};
\end{tikzpicture}
\end{center}
\legende{Frise chronologique : les diplômes s'acquièrent de gauche à droite (2023 $\rightarrow$ 2027), mais se présentent dans le CV dans l'ordre inverse (antichronologique).}
\end{popfigure}

\courssub{L'expérience professionnelle}

Cette rubrique recense les \cle{stages} (académiques, de perfectionnement), les \cle{emplois} occupés (jobs d'été, CDD, CDI) et les \cle{travaux pratiques} significatifs réalisés en atelier ou en entreprise. Pour chaque expérience, on précise quatre éléments : le \textbf{poste} occupé, la \textbf{structure} d'accueil (nom de l'entreprise, service), les \textbf{dates} (période ou durée) et les \textbf{missions} réalisées. Là encore, on respecte l'ordre antichronologique.

\begin{methode}[Rédiger une rubrique « Expérience » avec des verbes d'action]
\begin{enumerate}
\item Indique la période : « Juillet--septembre 2024 » ou « 3 mois (juil.--sept. 2024) ».
\item Nomme le poste et la structure : « Stagiaire en maintenance mécanique, SOCATERM, Douala ».
\item Décris deux ou trois missions principales à l'aide de \cle{verbes d'action} précis, à l'infinitif (ou au participe passé accordé) : \emph{concevoir, gérer, installer, contrôler, réparer, saisir, classer, accueillir, diagnostiquer, maintenir}.
\item Choisis uniquement les missions en rapport avec le poste visé (sélectivité).
\item Relis : aucune faute, aucune abréviation obscure (écrire « Société des Eaux du Cameroun » plutôt que « SNEC » si l'acronyme n'est pas connu de tous).
\end{enumerate}
\end{methode}

\begin{exemplebox}[Rubrique « Expérience » rédigée]
\textbf{EXPÉRIENCE PROFESSIONNELLE}\par
\begin{itemize}
\item \textbf{Juillet--septembre 2024} : Stagiaire en maintenance mécanique, SOCATERM, Douala. Installation de moteurs électriques, contrôle des soudures, rédaction de rapports de panne.
\item \textbf{Juillet 2023} : Aide-électricien, Entreprise KAMGA \& Fils, Bafoussam. Pose de câblage apparent, vérification des installations domestiques, respect des normes de sécurité.
\end{itemize}
\end{exemplebox}

\begin{attention}[Erreur fréquente : mentir sur un diplôme ou une expérience]
Certains candidats gonflent leur CV : diplôme inventé, stage jamais effectué, compétence exagérée (ex. « Maîtrise parfaite d'AutoCAD » alors qu'on a seulement vu l'interface). C'est une faute grave. L'employeur peut \textbf{vérifier} (appel à l'établissement, demande d'attestation de réussite, épreuve pratique à l'embauche, test de langue). Le mensonge découvert entraîne le \textbf{discrédit immédiat} du candidat, voire un licenciement pour faute lourde s'il a été embauché. Un CV honnête mais modeste vaut mieux qu'un CV mensonger.
\end{attention}

\courssub{Les compétences}

Cette rubrique distingue deux catégories majeures, très regardées par les recruteurs techniques :

\begin{itemize}
\item les \cle{compétences techniques} (savoir-faire métier) : informatique (Word, Excel, PowerPoint, logiciels de CAO/DAO type AutoCAD/SolidWorks, programmation), conduite d'engins (permis C, E), soudure (TIG, MIG, à l'arc), comptabilité (logiciels Sage, Ohada), secrétariat (prise de notes, archivage), électricité (câblage, lecture de schémas), maintenance préventive/curative ;
\item les \cle{compétences linguistiques} : français (lu, écrit, parlé — langue officielle), anglais (préciser le niveau : « scolaire », « lu/écrit/parlé », « niveau B1 », « technique »), langues nationales (ewondo, duala, fulfuldé, bassa... — préciser « parlé couramment » ou « notions »).
\end{itemize}

\begin{exemplebox}[Rubrique « Compétences »]
\textbf{COMPÉTENCES}\par
\begin{itemize}
\item Techniques : Maîtrise de la suite Microsoft Office (Word, Excel avancé -- TCD, fonctions SI) ; Lecture de plans mécaniques (vues en coupe, cotation ISO) ; Permis de conduire catégorie B (2025) ; Notions de soudure à l'arc (atelier lycée).
\item Linguistiques : Français (lu, écrit, parlé -- langue maternelle) ; Anglais (niveau B1 -- lu, écrit, parlé technique) ; Duala (parlé couramment).
\end{itemize}
\end{exemplebox}

\courssub{Les centres d'intérêt et les références}

Les \cle{centres d'intérêt} (ou activités extra-professionnelles) doivent rester \textbf{pertinents et valorisants} pour le poste : le sport collectif (football, handball $\rightarrow$ esprit d'équipe, endurance), la lecture (curiosité intellectuelle, veille technique), l'engagement associatif ou communautaire (bureau d'association, délégué de classe $\rightarrow$ sens des responsabilités, leadership), bénévolat (solidarité). Évite les mentions vagues et passives (« cinéma, musique, voyages, internet ») qui n'apportent aucune information sur tes soft skills.

Les \cle{références} (facultatives, souvent remplacées par « Références : sur demande ») sont des personnes que l'employeur peut contacter pour confirmer tes qualités professionnelles et humaines : ancien maître de stage, enseignant de matière technique, chef d'atelier. \textbf{Demande toujours leur accord} avant de citer leur nom, fonction, téléphone et courriel.

\courssub{La langue du CV : style et conventions}

Le CV obéit à un style télégraphique, normé, sans phrases complètes conjuguées :

\begin{itemize}
\item des \cle{phrases nominales} (sans verbe conjugué) : « Baccalauréat technique F3, 2027 » ; « Stage maintenance, SOCATERM, 2024 » ;
\item des \cle{verbes d'action} à l'infinitif (ou participe passé) pour les missions : « Installer, contrôler, gérer, diagnostiquer » ;
\item un \cle{vocabulaire précis} et professionnel (terminologie du métier), sans abréviations maison (sauf sigles très connus : BEPC, Bac, BTP, CAO, DAO, OHADA) ;
\item une \textbf{absence totale de fautes} d'orthographe, de grammaire, de syntaxe et de typographie : une seule erreur peut éliminer une candidature (manque de rigueur).
\end{itemize}

\begin{propriete}[Principe de sélectivité]
Un CV efficace est \cle{sélectif} : il ne retient que les informations \textbf{utiles au poste visé}. Tout ce qui n'éclaire pas le recruteur sur ton aptitude à occuper le poste doit être écarté. Le CV n'est pas une autobiographie exhaustive : c'est un outil de sélection et de marketing personnel. Adapte-le à chaque candidature.
\end{propriete}

\begin{exoresolu}[Construire la rubrique « Formation »]
Énoncé : Aïcha a obtenu le BEPC en 2023 au Collège de Maroua, le Probatoire technique F4 (Génie civil) en 2026 et passera le Baccalauréat technique F4 (Génie civil) en 2027 au Lycée technique de Maroua. Rédige sa rubrique « Formation » dans le bon ordre.
\tcblower
Solution : on applique l'ordre antichronologique, du plus récent (ou prévu) au plus ancien.\par
\textbf{FORMATION}\par
\begin{itemize}
\item \textbf{2027 (session de juin)} : Baccalauréat technique F4 (Génie civil), Lycée technique de Maroua (en cours) ;
\item \textbf{2026} : Probatoire technique F4 (Génie civil), Lycée technique de Maroua ;
\item \textbf{2023} : BEPC, Collège de Maroua.
\end{itemize}
\end{exoresolu}

\begin{exercice}[À toi de jouer]
Voici des informations en désordre concernant Alain, élève de Terminale F2 (Électrotechnique) : stage d'un mois à la SABC (Société Anonyme des Brasseries du Cameroun) en juillet 2024 (contrôle qualité boissons) ; Baccalauréat F2 prévu juin 2027 ; BEPC obtenu en 2023 au Lycée de Ngaoundéré ; job d'été vendeur en juillet-août 2022 (supermarché U) ; Probatoire F2 obtenu en 2026. Rédige les rubriques « Formation » et « Expérience » de son CV, dans l'ordre correct et avec des verbes d'action adaptés.
\end{exercice}

\begin{corrige}[Exercice : rubriques d'Alain]
\textbf{FORMATION}\par
\begin{itemize}
\item \textbf{2027 (prévu)} : Baccalauréat technique F2 (Électrotechnique), Lycée technique de Ngaoundéré ;
\item \textbf{2026} : Probatoire technique F2, Lycée technique de Ngaoundéré ;
\item \textbf{2023} : BEPC, Lycée de Ngaoundéré.
\end{itemize}

\textbf{EXPÉRIENCE PROFESSIONNELLE}\par
\begin{itemize}
\item \textbf{Juillet 2024} : Stagiaire contrôle qualité, SABC, Douala. Contrôler la conformité des produits (densité, taux d'alcool, étiquetage), rédiger les fiches de non-conformité, appliquer les normes HACCP.
\item \textbf{Juillet--août 2022} : Vendeur rayon électroménager, Supermarché U, Ngaoundéré. Accueillir et conseiller la clientèle, encaisser les paiements (espèces, mobile money), gérer le réapprovisionnement linéaire.
\end{itemize}

\emph{Remarque :} L'année 2026 comporte le Probatoire (rubrique Formation) ; l'année 2024 comporte le stage (rubrique Expérience). Chaque information est rangée dans la rubrique qui lui correspond par nature. L'ordre antichronologique est respecté dans chaque rubrique.
\end{corrige}

\begin{aretenir}[L'essentiel de la structure interne]
Un CV complet et standard comporte sept rubriques principales, dans cet ordre : \textbf{1. État civil} (identité, contacts), \textbf{2. Objectif professionnel} (poste visé), \textbf{3. Formation} (diplômes, antichronologique), \textbf{4. Expérience professionnelle} (stages, emplois, antichronologique, verbes d'action), \textbf{5. Compétences} (techniques + linguistiques), \textbf{6. Centres d'intérêt} (pertinents, soft skills), \textbf{7. Références} (sur demande ou nommées avec accord). Le style est télégraphique : phrases nominales, verbes d'action, vocabulaire précis, \textbf{zéro faute}. Le CV est \textbf{sélectif} (utile au poste) et \textbf{honnête} (vérifiable).
\end{aretenir}

\coursec{Langue : les valeurs du subjonctif et du conditionnel}

Après avoir maîtrisé la structure interne du CV — rubriques, mise en page, sélection des informations pertinentes —, il nous reste à soigner la \textbf{forme linguistique} de votre candidature. Un CV impeccable sur le plan visuel mais truffé de maladresses syntaxiques ou de tons inappropriés perd toute sa crédibilité. Cette section vise à vous donner les outils grammaticaux pour exprimer avec justesse votre volonté, votre politesse, vos hypothèses et votre prudence, en mobilisant deux modes aux valeurs riches : le \cle{subjonctif} et le \cle{conditionnel}.

\courssub{Rappel de formation : subjonctif présent et conditionnel présent}

Avant d'explorer les valeurs, assurons-nous que la mécanique verbale est huilée. En Terminale technique, la conjugaison n'est plus un apprentissage mais un \textbf{outil de précision}.

\begin{definition}[Formation du subjonctif présent]
Pour la grande majorité des verbes, le radical du subjonctif présent se construit sur la \textbf{3\textsuperscript{e} personne du pluriel de l'indicatif présent} (ils/elles), auquel on ajoute les terminaisons : \textbf{-e, -es, -e, -ions, -iez, -ent}.
\begin{itemize}
    \item \emph{Parler} (ils parl\textbf{ent}) $\rightarrow$ que je parl\textbf{e}, que tu parl\textbf{es}, qu'il parl\textbf{e}, que nous parl\textbf{ions}, que vous parl\textbf{iez}, qu'ils parl\textbf{ent}.
    \item \emph{Finir} (ils fin\textbf{issent}) $\rightarrow$ que je fin\textbf{isse}, que tu fin\textbf{isses}, qu'il fin\textbf{isse}, que nous fin\textbf{issions}, que vous fin\textbf{issiez}, qu'ils fin\textbf{issent}.
    \item \emph{Prendre} (ils pren\textbf{nent}) $\rightarrow$ que je pren\textbf{ne}, que tu pren\textbf{nes}, qu'il pren\textbf{ne}, que nous pren\textbf{nions}, que vous pren\textbf{niez}, qu'ils pren\textbf{nent}.
\end{itemize}
\textbf{Verbes irréguliers fréquents dans le CV :} \emph{être} (que je sois), \emph{avoir} (que j'aie), \emph{faire} (que je fasse), \emph{aller} (que j'aille), \emph{savoir} (que je sache), \emph{pouvoir} (que je puisse), \emph{vouloir} (que je veuille), \emph{devoir} (que je doive), \emph{falloir} (qu'il faille).
\end{definition}

\begin{definition}[Formation du conditionnel présent]
Le conditionnel présent utilise le \textbf{radical du futur simple} (souvent l'infinitif pour les verbes du 1\textsuperscript{er} et 2\textsuperscript{e} groupe, radical irrégulier pour le 3\textsuperscript{e} groupe) + les \textbf{terminaisons de l'imparfait de l'indicatif} : \textbf{-ais, -ais, -ait, -ions, -iez, -aient}.
\begin{itemize}
    \item \emph{Parler} $\rightarrow$ je parl\textbf{erais}.
    \item \emph{Finir} $\rightarrow$ je fin\textbf{irais}.
    \item \emph{Prendre} $\rightarrow$ je prend\textbf{rais}.
    \item \emph{Être} $\rightarrow$ je s\textbf{erais} ; \emph{Avoir} $\rightarrow$ j'\textbf{aurais} ; \emph{Faire} $\rightarrow$ je f\textbf{erais} ; \emph{Aller} $\rightarrow$ j'\textbf{irais} ; \emph{Vouloir} $\rightarrow$ je v\textbf{oudrais} ; \emph{Pouvoir} $\rightarrow$ je p\textbf{ourrais} ; \emph{Devoir} $\rightarrow$ je d\textbf{evrais}.
\end{itemize}
\end{definition}

\begin{attention}[Piège graphique : \textbf{-rais} vs \textbf{-rai}]
Ne confondez pas \textbf{conditionnel} (\emph{je voudrais, je pourrais}) et \textbf{futur simple} (\emph{je voudrai, je pourrai}).
\begin{itemize}
    \item \emph{Je \textbf{voudrais} un entretien} (souhait poli, conditionnel).
    \item \emph{Je \textbf{voudrai} un entretien} (certitude future, futur simple — inapproprié dans une lettre de motivation).
\end{itemize}
À l'oral, la différence s'entend souvent : [vudr\textbf{$\varepsilon$}] (conditionnel) vs [vudr\textbf{e}] (futur). À l'écrit, la vigilance est de mise.
\end{attention}

\courssub{Le subjonctif : le mode du subjectif, du souhaité et du nécessaire}

Le subjonctif ne décrit pas un fait réel (c'est le rôle de l'indicatif), il exprime une \textbf{réalité virtuelle, souhaitée, doutée, nécessaire ou ressentie}. Dans un CV ou une lettre, il structure votre rapport à l'employeur : ce que vous \emph{voulez} qu'il arrive, ce que vous \emph{jugez nécessaire}, ce que vous \emph{ressentez}.

\begin{propriete}[Valeurs principales du subjonctif en contexte professionnel]
\begin{itemize}
    \item \textbf{Volonté / Souhait / Désir :} Verbes \emph{vouloir, souhaiter, désirer, préférer, aimer mieux, exiger} + \textbf{que} + subjonctif.
    \item \textbf{Obligation / Nécessité / Impératif :} Expressions \emph{il faut que, il est nécessaire que, il est indispensable que, il est impératif que} + subjonctif.
    \item \textbf{Doute / Incertitude / Probabilité :} Verbes \emph{douter, nier, ne pas croire, ne pas penser, il est douteux que, il est possible que, il se peut que} + subjonctif.
    \item \textbf{Sentiment / Émotion / Jugement subjectif :} Adjectifs/verbes \emph{être content que, regretter que, craindre que, apprécier que, trouver normal que, il est important que} + subjonctif.
\end{itemize}
\end{propriete}

\begin{exemplebox}[Le subjonctif de la volonté et du souhait dans la candidature]
\begin{itemize}
    \item \emph{Je \textbf{souhaite} que ma candidature \textbf{retienne} votre attention.} (Volonté tournée vers l'autre).
    \item \emph{Nous \textbf{voulons} que vous \textbf{puissiez} apprécier mes compétences.} (Volonté de l'émetteur sur l'action du destinataire).
    \item \emph{Je \textbf{préfère} que l'entretien \textbf{ait} lieu la semaine prochaine.} (Choix exprimé poliment).
\end{itemize}
\textbf{Observation :} Le sujet de la principale (Je) est \textbf{différent} du sujet de la subordonnée (ma candidature, vous, l'entretien). C'est la condition syntaxique majeure du subjonctif (sauf cas de \emph{il faut que} où le sujet est impersonnel).
\end{exemplebox}

\begin{exemplebox}[Le subjonctif de l'obligation et de la nécessité]
\begin{itemize}
    \item \emph{Il \textbf{faut} que je \textbf{sois} disponible immédiatement.} (Contrainte objective, impersonnelle).
    \item \emph{Il est \textbf{indispensable} que vous \textbf{receviez} mes références avant vendredi.} (Nécessité professionnelle).
    \item \emph{Il est \textbf{nécessaire} que le contrat \textbf{prenne} effet au 1\textsuperscript{er} juillet.} (Impératif calendrier).
\end{itemize}
\textbf{Note :} \emph{Il faut que} + subjonctif est la structure standard de l'obligation. Évitez \emph{Il faut que je dois} (pléonasme fautif).
\end{exemplebox}

\begin{exemplebox}[Le subjonctif du doute et du sentiment : nuances de candidature]
\begin{itemize}
    \item \emph{Je \textbf{crains} que mon expérience ne \textbf{soit} pas tout à fait adaptée.} (Crainte + \textbf{ne} explétif, pas de négation).
    \item \emph{Je \textbf{doute} que ce poste \textbf{corresponde} exactement à mon profil.} (Doute réel).
    \item \emph{Il est \textbf{possible} que je \textbf{doive} déménager.} (Probabilité, incertitude).
    \item \emph{Je \textbf{suis ravi} que vous \textbf{ayez} pris le temps de lire mon CV.} (Sentiment positif sur un fait accompli).
\end{itemize}
\end{exemplebox}

\begin{methode}[Identifier et justifier le subjonctif dans un texte professionnel]
\begin{enumerate}
    \item \textbf{Repérer le déclencheur :} Verbe de volonté (\emph{vouloir, souhaiter}), d'obligation (\emph{il faut que}), de sentiment (\emph{craindre, regretter}), de doute (\emph{douter, nier}) ou expression impersonnelle (\emph{il est important que}).
    \item \textbf{Vérifier la structure :} \textbf{Que} (ou \emph{qu'}) introduit une proposition subordonnée dont le sujet est \textbf{différent} de celui de la principale (sauf impersonnel \emph{il faut que}).
    \item \textbf{Analyser la valeur :} Est-ce un souhait ? Une obligation ? Un doute ? Une émotion ? Cela confirme le choix du mode.
    \item \textbf{Conjuguer correctement :} Radical de la 3\textsuperscript{e} pers. pl. indicatif présent + terminaisons \emph{-e, -es, -e, -ions, -iez, -ent}.
\end{enumerate}
\end{methode}

\courssub{Le conditionnel : le mode de l'hypothèse, de la politesse et de l'atténuation}

Si le subjonctif exprime ce qui \emph{n'est pas encore réel mais souhaité/nécessaire}, le conditionnel exprime ce qui \emph{serait réel sous condition} (hypothèse) ou ce que l'on \emph{présente comme non imposé} (politesse, atténuation). C'est \textbf{le mode roi de la lettre de motivation}.

\begin{propriete}[Valeurs principales du conditionnel en contexte professionnel]
\begin{itemize}
    \item \textbf{Politesse / Courtoisie / Atténuation de la demande :} \emph{Je \textbf{voudrais}, j'\textbf{aimerais}, je \textbf{souhaiterais}, je \textbf{serais reconnaissant de}} + infinitif. Adoucit l'injonction directe (\emph{Je veux, Donnez-moi}).
    \item \textbf{Hypothèse / Condition (si + imparfait $\rightarrow$ conditionnel présent) :} \emph{Si j'\textbf{avais} le poste, je \textbf{m'investirais} pleinement.} Projection dans un avenir conditionnel.
    \item \textbf{Atténuation journalistique / Prudence (Conditionnel de distanciation) :} \emph{Mon stage \textbf{aurait permis} de développer mes compétences.} (Prudence sur le résultat, modestie). \emph{Cette expérience \textbf{serait} un atout.} (Hypothèse modérée, pas certitude arrogante).
    \item \textbf{Futur dans le passé (Récit de démarche) :} \emph{Il m'a dit qu'il \textbf{me rappellerait} la semaine prochaine.}
\end{itemize}
\end{propriete}

\begin{exemplebox}[Le conditionnel de politesse : formules types de la lettre de motivation]
\begin{itemize}
    \item \emph{Je \textbf{voudrais} vous soumettre ma candidature au poste de technicien de maintenance.}
    \item \emph{J'\textbf{aimerais} vous rencontrer afin de vous exposer mes motivations.}
    \item \emph{Je \textbf{souhaiterais} obtenir un entretien à votre convenance.}
    \item \emph{Je \textbf{serais} honoré de pouvoir échanger avec vous sur mon parcours.}
    \item \emph{\textbf{Auriez-vous} la disponibilité pour me recevoir la semaine prochaine ?} (Inversion sujet-verbe, très formel).
\end{itemize}
\textbf{Pourquoi pas l'indicatif ?} \emph{Je veux ce poste} = exigence, manque de recul. \emph{Je voudrais ce poste} = demande respectueuse, conscience de la sélection.
\end{exemplebox}

\begin{exemplebox}[Le conditionnel hypothétique : projection professionnelle]
Structure type : \textbf{Si + Imparfait (condition) $\rightarrow$ Conditionnel Présent (conséquence)}.
\begin{itemize}
    \item \emph{Si je \textbf{devais} rejoindre votre entreprise, je \textbf{mettrais} mes compétences en automatisme à votre service.}
    \item \emph{Si l'entretien \textbf{avait} lieu mardi, je \textbf{pourrais} vous présenter mon portfolio.}
    \item \emph{Si vous \textbf{me donniez} votre accord, je \textbf{serais} opérationnel dès le 1\textsuperscript{er} septembre.}
\end{itemize}
\textbf{Erreur fatale (voir \textbf{Attention} ci-dessous) :} \emph{Si j'\textbf{aurais} le temps...} $\rightarrow$ \textbf{Jamais de conditionnel après \textbf{si} !}
\end{exemplebox}

\begin{exemplebox}[Le conditionnel journalistique / atténué : prudence et modestie]
Dans le CV (rubrique \emph{Compétences}, \emph{Réalisations}) ou la lettre, le conditionnel permet d'éviter l'affirmation péremptoire ou l'autosatisfaction.
\begin{itemize}
    \item \emph{Cette formation \textbf{m'aurait permis} d'acquérir une rigueur méthodologique.} (Plutôt que \emph{a permis} : je ne juge pas moi-même le résultat final).
    \item \emph{Mon stage chez X \textbf{serait} un atout pour ce poste.} (Hypothèse modérée, laisse l'employeur juge).
    \item \emph{Je \textbf{serais} en mesure de gérer une équipe de 5 personnes.} (Disponibilité potentielle, pas certitude actuelle).
\end{itemize}
\end{exemplebox}

\begin{attention}[Erreurs fréquentes : les 3 pièges mortels au Bac / Probatoire]
\begin{enumerate}
    \item \textbf{« Si j'aurais su, j'aurais pas venu »} $\rightarrow$ \textbf{Jamais de conditionnel après \textbf{si} (condition).}
        \begin{itemize}
            \item Correct : \emph{Si j'\textbf{avais} su, je \textbf{n'aurais pas venu} (ou \textbf{ne serais pas venu}).}
            \item Règle : \textbf{Si + Imparfait (ou Plus-que-parfait) $\rightarrow$ Conditionnel Présent (ou Passé).}
        \end{itemize}
    \item \textbf{Subjonctif après « après que »} $\rightarrow$ \textbf{Faux.}
        \begin{itemize}
            \item \emph{Après qu'il \textbf{a reçu} mon CV, il m'a appelé.} (Indicatif : antériorité réelle, fait accompli).
            \item \emph{Avant qu'il \textbf{reçoive} mon CV...} (Subjonctif : antériorité virtuelle, non encore réalisée).
        \end{itemize}
    \item \textbf{Concordance des temps (Si passé $\rightarrow$ Conditionnel passé) :}
        \begin{itemize}
            \item \emph{Si j'\textbf{avais eu} le bac l'an dernier, j'\textbf{aurais intégré} cette école.} (Plus-que-parfait $\rightarrow$ Conditionnel Passé).
            \item Ne pas écrire : \emph{Si j'avais eu, j'intègre} (mélange incohérent).
        \end{itemize}
\end{enumerate}
\end{attention}

\courssub{Synthèse visuelle : valeurs, structures et exemples CV}

Le tableau ci-dessous récapitule les articulations entre déclencheurs, modes, valeurs et formulations types pour votre dossier de candidature.

\begin{popfigure}
\centering
\begin{tikzpicture}[
    cell/.style={rectangle, draw=popline, minimum height=1.1cm, align=center, text width=3.2cm, font=\small},
    head/.style={cell, fill=popblue, text=white, font=\small\bfseries},
    rowalt/.style={fill=popblueL!30},
    rowstd/.style={fill=white},
    valcell/.style={rectangle, draw=popline, minimum height=1.1cm, align=left, text width=4.5cm, font=\small\itshape}
]
% En-têtes
\node[head] (m1) at (0,0) {Déclencheur / Contexte};
\node[head] (m2) at (3.4,0) {Mode};
\node[head] (m3) at (6.8,0) {Valeur principale};
\node[head] (m4) at (10.2,0) {Exemple CV / Lettre};
% Ligne 1
\node[rowalt, cell, align=center] (r1c1) at (0,-1.1){Verbe de volonté :\\ vouloir, souhaiter, préférer};
\node[rowalt, cell] (r1c2) at (3.4,-1.1) {Subjonctif};
\node[rowalt, cell] (r1c3) at (6.8,-1.1) {Souhait, désir};
\node[rowalt, valcell] (r1c4) at (10.2,-1.1) {Je \textbf{souhaite} que ma candidature \textbf{retienne} votre attention.};
% Ligne 2
\node[rowstd, cell] (r2c1) at (0,-2.2) {Il faut que, il est nécessaire / indispensable / impératif que};
\node[rowstd, cell] (r2c2) at (3.4,-2.2) {Subjonctif};
\node[rowstd, cell] (r2c3) at (6.8,-2.2) {Obligation, nécessité};
\node[rowstd, valcell] (r2c4) at (10.2,-2.2) {Il \textbf{faut} que je \textbf{sois} disponible dès lundi.};
% Ligne 3
\node[rowalt, cell, align=center] (r3c1) at (0,-3.3){Verbe de sentiment :\\ craindre, regretter, apprécier, être content};
\node[rowalt, cell] (r3c2) at (3.4,-3.3) {Subjonctif};
\node[rowalt, cell] (r3c3) at (6.8,-3.3) {Émotion, jugement subjectif};
\node[rowalt, valcell] (r3c4) at (10.2,-3.3) {Je \textbf{crains} que mon profil ne \textbf{soit} trop généraliste.};
% Ligne 4
\node[rowstd, cell, align=center] (r4c1) at (0,-4.4){Verbe de doute / négation :\\ douter, nier, ne pas croire, il est possible que};
\node[rowstd, cell] (r4c2) at (3.4,-4.4) {Subjonctif};
\node[rowstd, cell] (r4c3) at (6.8,-4.4) {Incertitude, probabilité};
\node[rowstd, valcell] (r4c4) at (10.2,-4.4) {Il est \textbf{possible} que je \textbf{doive} passer un test technique.};
% Ligne 5
\node[rowalt, cell, align=center] (r5c1) at (0,-5.5){Demande polie, courtoisie :\\ voudrais, aimerais, souhaiterais, serais reconnaissant};
\node[rowalt, cell] (r5c2) at (3.4,-5.5) {Conditionnel};
\node[rowalt, cell] (r5c3) at (6.8,-5.5) {Politesse, atténuation};
\node[rowalt, valcell] (r5c4) at (10.2,-5.5) {J'\textbf{aimerais} vous rencontrer pour un entretien.};
% Ligne 6
\node[rowstd, cell, align=center] (r6c1) at (0,-6.6){Hypothèse conditionnelle :\\ Si + Imparfait};
\node[rowstd, cell] (r6c2) at (3.4,-6.6) {Conditionnel};
\node[rowstd, cell] (r6c3) at (6.8,-6.6) {Projection, conséquence hypothétique};
\node[rowstd, valcell] (r6c4) at (10.2,-6.6) {Si j'\textbf{intégrais} l'équipe, je \textbf{participerais} aux projets en cours.};
% Ligne 7
\node[rowalt, cell, align=center] (r7c1) at (0,-7.7){Prudence, modestie, distanciation :\\ description de compétences / résultats};
\node[rowalt, cell] (r7c2) at (3.4,-7.7) {Conditionnel};
\node[rowalt, cell] (r7c3) at (6.8,-7.7) {Atténuation, conditionnel journalistique};
\node[rowalt, valcell] (r7c4) at (10.2,-7.7) {Cette expérience \textbf{aurait renforcé} mon sens des responsabilités.};
\end{tikzpicture}
\legende{Tableau de synthèse : valeurs du subjonctif et du conditionnel en contexte de candidature (CV / Lettre de motivation).}
\end{popfigure}

\courssub{Entraînement guidé : transformer et corriger}

La maîtrise s'acquiert par la manipulation. Voici deux exercices types d'examen (Probatoire / Baccalauréat technique) : transformation de registre et correction d'erreurs.

\begin{exemplebox}[Transformation : du direct au courtois (Indicatif/Impératif $\rightarrow$ Conditionnel/Subjonctif)]
\textbf{Consigne :} Réécrivez les phrases suivantes en adoptant le ton formel et courtois d'une lettre de motivation.
\begin{enumerate}
    \item \emph{Je veux ce poste.} $\rightarrow$ \textbf{Je \textbf{voudrais} / \textbf{souhaiterais} obtenir ce poste.}
    \item \emph{Donnez-moi un entretien.} $\rightarrow$ \textbf{Je \textbf{serais reconnaissant} de vous \textbf{accorder} un entretien.} / \textbf{J'\textbf{aimerais} que vous \textbf{m'accordiez} un entretien.} (Subjonctif après verbe de volonté).
    \item \emph{Il faut que je parte tôt.} $\rightarrow$ \textbf{Il \textbf{serait nécessaire} que je \textbf{puisse} partir plus tôt.} / \textbf{J'\textbf{aimerais} pouvoir partir plus tôt.}
    \item \emph{Je suis sûr que je vais réussir.} $\rightarrow$ \textbf{Je \textbf{suis convaincu} que je \textbf{réussirais} (conditionnel d'atténuation) / que je \textbf{saurais} relever le défi.}
    \item \emph{Mon stage m'a appris la rigueur.} $\rightarrow$ \textbf{Mon stage \textbf{m'aurait permis} d'acquérir de la rigueur.} (Conditionnel atténué).
\end{enumerate}
\end{exemplebox}

\begin{exoresolu}[Exemple chiffré : 6 phrases de candidature à corriger]
\textbf{Énoncé :} Les phrases ci-dessous contiennent chacune une erreur de mode (subjonctif/conditionnel/indicatif) ou de concordance. Identifiez l'erreur, expliquez-la et réécrivez la phrase correcte.

\begin{enumerate}
    \item \emph{« Je voudrais que vous \textbf{recevrez} mon CV avant vendredi. »}
    \item \emph{« Si j'\textbf{aurais} su que le poste était libre, j'\textbf{aurais postulé} plus tôt. »}
    \item \emph{« Après que j'\textbf{aie} envoyé la lettre, j'ai reçu un accusé de réception. »}
    \item \emph{« Il faut que je \textbf{dois} mettre à jour mes compétences informatiques. »}
    \item \emph{« Je doute que vos horaires \textbf{sont} flexibles. »}
    \item \emph{« Cette formation me \textbf{permettra} (futur) d'être opérationnel, je l'espère. »} (Contexte : lettre d'accompagnement, ton souhaité modéré).
\end{enumerate}

\tcblower
\textbf{Solutions détaillées :}

\begin{enumerate}
    \item \textbf{Erreur :} Indicatif futur (\emph{recevrez}) après verbe de volonté (\emph{vouloir que}). \textbf{Règle :} Volonté $\rightarrow$ Subjonctif. \textbf{Correction :} \emph{Je voudrais que vous \textbf{receviez} mon CV avant vendredi.}
    \item \textbf{Erreur :} Conditionnel (\emph{aurais}) après \textbf{si} (condition). \textbf{Règle :} \emph{Si} + Imparfait (ou PQP) $\rightarrow$ Conditionnel. \textbf{Correction :} \emph{Si j'\textbf{avais} su que le poste était libre, j'\textbf{aurais postulé} plus tôt.}
    \item \textbf{Erreur :} Subjonctif (\emph{aie}) après \emph{après que}. \textbf{Règle :} \emph{Après que} + Indicatif (fait accompli) ; \emph{Avant que} + Subjonctif. \textbf{Correction :} \emph{Après que j'\textbf{ai envoyé} la lettre, j'ai reçu un accusé de réception.}
    \item \textbf{Erreur :} Pléonasme / Doublon modal (\emph{falloir} + \emph{devoir}). \textbf{Règle :} \emph{Il faut que} + Subjonctif suffit à exprimer l'obligation. \textbf{Correction :} \emph{Il faut que je \textbf{mette} à jour mes compétences...} OU \emph{Je \textbf{dois} mettre à jour...}
    \item \textbf{Erreur :} Indicatif (\emph{sont}) après \emph{douter que}. \textbf{Règle :} Doute / Négation $\rightarrow$ Subjonctif. \textbf{Correction :} \emph{Je doute que vos horaires \textbf{soient} flexibles.}
    \item \textbf{Erreur de registre :} Futur simple (\emph{permettra}) trop assertif pour une lettre de motivation où l'on projette une compétence future non acquise. \textbf{Correction privilégiée :} \emph{Cette formation me \textbf{permettrait} d'être opérationnel...} (Conditionnel d'atténuation/hypothèse) ou \emph{Cette formation \textbf{me permettra}...} si certitude totale assumée. Au Bac, le conditionnel est attendu pour la modestie.
\end{enumerate}
\end{exoresolu}

\courssub{Application intégrée : rédaction d'un paragraphe « Motivation / Disponibilité »}

Pour clore cette section langue, appliquons les deux modes dans un paragraphe cohérent, typique de la fin d'une lettre de motivation ou de la zone « Objectif professionnel » d'un CV moderne.

\begin{experience}[Rédaction guidée : paragraphe de clôture de lettre de motivation]
\textbf{Consigne :} Rédigez un paragraphe (5-6 lignes) en utilisant \textbf{au moins} : 1 subjonctif de volonté, 1 subjonctif de nécessité, 1 conditionnel de politesse, 1 conditionnel hypothétique, 1 conditionnel atténué.

\textbf{Proposition de rédaction :}
\emph{« \textbf{Je souhaiterais} (cond. politesse) vous rencontrer afin de vous présenter mon parcours. \textbf{Il est indispensable} (subj. nécessité) \textbf{que cet entretien me permette} (subj. après expression d'impératif) de valoriser mon expérience en maintenance préventive. \textbf{Si vous me donniez} (hypothèse : si + imparfait) \textbf{cette opportunité, je m'investirais} (cond. conséquence) \textbf{pleinement dans les projets de l'équipe.} Mon stage chez SONEL \textbf{m'aurait permis} (cond. atténué) de développer une rigueur que je \textbf{mettrais} (cond. hypothétique/futur dans le passé projeté) au service de votre structure. »}

\textbf{Analyse grammaticale :}
\begin{itemize}
    \item \emph{Souhaiterais} : Conditionnel présent, politesse.
    \item \emph{Permette} : Subjonctif présent (après \emph{il est indispensable que}), nécessité.
    \item \emph{Donniez / m'investirais} : Si + Imparfait $\rightarrow$ Conditionnel, hypothèse.
    \item \emph{Aurait permis} : Conditionnel passé, atténuation / modestie sur le passé.
    \item \emph{Mettrais} : Conditionnel, projection future conditionnelle.
\end{itemize}
\end{experience}

\begin{aretenir}[Check-list « Modes \& Candidature » pour le Bac / Probatoire]
\begin{itemize}
    \item \textbf{Subjonctif :} Après \emph{que} (volonté, nécessité, sentiment, doute) \textbf{si sujets différents}. Formation : radical 3\textsuperscript{e} pl. indicatif + \emph{-e, -es, -e, -ions, -iez, -ent}.
    \item \textbf{Conditionnel :} Politesse (\emph{voudrais, aimerais}), Hypothèse (\emph{Si} + imparfait $\rightarrow$ cond.), Atténuation (\emph{serait, aurait permis}).
    \item \textbf{3 Interdits majeurs :} 1) \emph{Si} + Conditionnel. 2) \emph{Après que} + Subjonctif. 3) \emph{Il faut que} + \emph{devoir} (pléonasme).
    \item \textbf{Concordance :} \emph{Si} + PQP $\rightarrow$ Conditionnel Passé (\emph{aurais fait}).
    \item \textbf{Orthographe :} \emph{-rais} (cond.) $\neq$ \emph{-rai} (futur). \emph{Je pourrais} vs \emph{Je pourrai}.
\end{itemize}
\end{aretenir}

\begin{savaistu}[Le conditionnel, marqueur de compétence « soft » au Cameroun]
Dans le monde professionnel camerounais (administration, entreprises parapubliques, PME), la maîtrise des formules de politesse au conditionnel (\emph{Je voudrais solliciter..., Je saurais gré...}) est perçue comme un indicateur fort de \textbf{professionnalisme} et de \textbf{maîtrise des codes de l'écrit administratif}. Un candidat qui écrit « Je veux un poste » ou « Donnez-moi du travail » dans une lettre formelle est souvent écarté avant même la lecture du CV, perçu comme manquant de « savoir-être ». La grammaire est ici un levier d'employabilité.
\end{savaistu}

\begin{exercice}[Entraînement autonome : Modes et candidature]
\textbf{1. Conjuguez le verbe entre parenthèses au mode et au temps appropriés (Subjonctif présent ou Conditionnel présent).}
\begin{enumerate}
    \item Il est urgent que le technicien \_\_\_\_\_\_\_\_\_\_ (intervenir) sur la panne.
    \item Je \_\_\_\_\_\_\_\_\_\_ (aimer) que vous me confirmiez la date de l'entretien.
    \item Si le directeur \_\_\_\_\_\_\_\_\_\_ (accepter) ma candidature, je \_\_\_\_\_\_\_\_\_\_ (être) ravi.
    \item Je \_\_\_\_\_\_\_\_\_\_ (souhaiter) que mon dossier \_\_\_\_\_\_\_\_\_\_ (être) complet.
    \item \_\_\_\_\_\_\_\_\_\_ (Avoir - vous) l'amabilité de m'indiquer la procédure ?
\end{enumerate}

\textbf{2. Corrigez les erreurs dans les phrases suivantes (une erreur par phrase).}
\begin{enumerate}
    \item Si j'aurais le temps, je viendrais vous voir.
    \item Après qu'il ait signé le contrat, il est parti.
    \item Il faut que je dois réviser mon CV.
    \item Je doute qu'il vient demain.
    \item Je veux que vous me donnez une réponse.
\end{enumerate}

\textbf{3. Transformez ces phrases directes en formulations courtoises (Conditionnel de politesse ou Subjonctif de volonté).}
\begin{enumerate}
    \item Je veux le poste de commercial.
    \item Envoyez-moi la fiche de poste.
    \item Je peux commencer lundi.
    \item Il faut que tu me répondes vite.
\end{enumerate}
\end{exercice}

\begin{corrige}[Exercice n°18 - Modes et candidature]
\textbf{1. Conjugaison :}
\begin{enumerate}
    \item \textbf{intervienne} (Subj. présent — nécessité, \emph{il est urgent que}).
    \item \textbf{aimerais} (Cond. présent — politesse) / \textbf{confirmiez} (Subj. présent — volonté, \emph{que vous}).
    \item \textbf{acceptait} (Imparfait — \emph{si} condition) / \textbf{serais} (Cond. présent — conséquence).
    \item \textbf{souhaite} (Indicatif — fait réel) / \textbf{soit} (Subj. présent — volonté, \emph{que mon dossier}).
    \item \textbf{Auriez-vous} (Cond. présent — politesse, inversion).
\end{enumerate}

\textbf{2. Corrections :}
\begin{enumerate}
    \item \emph{Si j'\textbf{avais} le temps, je viendrais vous voir.} (\emph{Si} + Imparfait).
    \item \emph{Après qu'il \textbf{a signé} le contrat, il est parti.} (\emph{Après que} + Indicatif).
    \item \emph{Il faut que je \textbf{revoie} / \textbf{revisite} mon CV.} (Subjonctif, pas \emph{dois}).
    \item \emph{Je doute qu'il \textbf{vienne} demain.} (\emph{Douter que} + Subjonctif).
    \item \emph{Je \textbf{voudrais} / \textbf{souhaiterais} que vous me \textbf{donniez} une réponse.} (Conditionnel de politesse + Subjonctif de volonté).
\end{enumerate}

\textbf{3. Transformations (propositions) :}
\begin{enumerate}
    \item \emph{Je \textbf{voudrais} / \textbf{souhaiterais} obtenir le poste de commercial.}
    \item \emph{Je \textbf{serais reconnaissant} de recevoir la fiche de poste.} / \emph{J'\textbf{aimerais} que vous me \textbf{envoyiez} la fiche de poste.}
    \item \emph{Je \textbf{pourrais} commencer lundi.} / \emph{Je \textbf{serais disponible} pour commencer lundi.}
    \item \emph{Il \textbf{serait nécessaire} que vous me \textbf{répondiez} dans les meilleurs délais.} / \emph{J'\textbf{aimerais} que vous me \textbf{répondiez} rapidement.}
\end{enumerate}
\end{corrige}

\coursec{Rédiger et adapter un CV}

Dans la section précédente, nous avons étudié les valeurs du \cle{subjonctif} et du \cle{conditionnel}, ces modes qui permettent d'exprimer le souhait, l'hypothèse ou la politesse : « Je souhaiterais intégrer votre entreprise », « Il faudrait que je sois disponible dès le mois de juin ». Ces formes verbales trouvent leur emploi naturel dans la lettre de motivation et dans l'entretien d'embauche. Mais avant de parler, il faut un document qui parle pour nous : le \cle{CV}. Nous avons analysé sa structure externe (rubriques, mise en page) et sa structure interne (contenu de chaque rubrique). Il est temps maintenant de passer à l'acte : \textbf{rédiger} un CV, puis l'\textbf{adapter} à une situation d'embauche précise, enfin le \textbf{présenter} à l'oral. C'est l'aboutissement de toute l'unité, et c'est exactement ce que le Baccalauréat technique peut vous demander de faire.

\courssub{Les cinq étapes de la rédaction d'un CV}

Rédiger un CV ne s'improvise pas. Comme toute production de texte au programme de Terminale technique, elle obéit à une démarche en étapes. Observons d'abord une situation concrète : deux élèves de Terminale F4 répondent à la même offre d'emploi. Le premier envoie en une heure un CV rédigé « au feeling » ; le second consacre une journée à analyser l'offre, à trier ses informations et à peaufiner sa mise en page. Lequel sera convoqué à l'entretien ? L'expérience des recruteurs est sans appel : un CV se construit méthodiquement, en cinq étapes.

\begin{popfigure}
\begin{center}
\begin{tikzpicture}[
  node distance=0.55cm,
  etape/.style={rectangle, rounded corners=4pt, draw=popblue, fill=popblueL, thick, text width=9.6cm, align=left, inner sep=6pt, font=\small},
  num/.style={circle, draw=popdark, fill=poporange, text=white, font=\bfseries\small, inner sep=2.2pt, minimum size=0.62cm},
  fleche/.style={-{Stealth[length=3mm]}, thick, popdark}
]
\node[etape] (e1) {\textbf{Analyser l'offre} : poste proposé, profil recherché, compétences exigées, mots-clés.};
\node[num, left=0.25cm of e1] {1};
\node[etape, below=of e1] (e2) {\textbf{Inventorier ses informations} : formation, stages, expériences, compétences, langues, centres d'intérêt.};
\node[num, left=0.25cm of e2] {2};
\node[etape, below=of e2] (e3) {\textbf{Choisir et ordonner les rubriques} : la rubrique la plus pertinente pour le poste vient en premier.};
\node[num, left=0.25cm of e3] {3};
\node[etape, below=of e3] (e4) {\textbf{Rédiger} : style télégraphique maîtrisé, verbes d'action, orthographe irréprochable.};
\node[num, left=0.25cm of e4] {4};
\node[etape, below=of e4] (e5) {\textbf{Mettre en page et relire} : présentation aérée, relecture personnelle puis relecture par un tiers.};
\node[num, left=0.25cm of e5] {5};
\draw[fleche] (e1) -- (e2);
\draw[fleche] (e2) -- (e3);
\draw[fleche] (e3) -- (e4);
\draw[fleche] (e4) -- (e5);
\end{tikzpicture}
\end{center}
\legende{Organigramme des cinq étapes de rédaction d'un CV : chaque étape prépare la suivante ; aucune ne peut être sautée.}
\end{popfigure}

\textbf{Étape 1 : analyser l'offre d'emploi ou la situation d'embauche.} Tout commence par une lecture attentive. Une offre d'emploi contient toujours trois informations essentielles : le \cle{poste} proposé (intitulé, lieu, missions), le \cle{profil recherché} (diplôme exigé, expérience, qualités humaines) et les \cle{compétences} attendues (savoir-faire techniques, langues, outils). Soulignez les mots-clés : ce sont eux que le recruteur cherchera dans votre CV. À l'examen, la « situation donnée » joue le même rôle que l'offre : lisez-la deux fois et repérez ce qu'elle exige.

\textbf{Étape 2 : inventorier ses informations personnelles.} Sur une feuille de brouillon, listez exhaustivement : votre formation (diplômes, établissements, années), vos stages et expériences (même informels : un stage en atelier, un travail saisonnier), vos compétences (techniques, informatiques, linguistiques), vos centres d'intérêt. À ce stade, on ne trie pas encore : on rassemble la matière première.

\textbf{Étape 3 : choisir et ordonner les rubriques.} Toutes les informations ne se valent pas face à un poste donné. La règle d'or : \textbf{la rubrique la plus pertinente pour le poste vient en premier}. Pour un \cle{jeune diplômé} sans longue expérience, c'est la rubrique « Formation » qui ouvre le CV ; pour un technicien expérimenté, ce sera « Expérience professionnelle ». On ordonne ensuite les rubriques par ordre d'intérêt décroissant pour le recruteur.

\textbf{Étape 4 : rédiger.} Le CV exige un \cle{style télégraphique} : des phrases courtes, sans article ni pronom personnel (« Installation et maintenance de groupes froids » et non « J'ai installé et j'ai maintenu des groupes froids »). Chaque ligne commence idéalement par un \cle{verbe d'action} à l'infinitif ou au participe passé : installer, réparer, gérer, concevoir, accueillir, saisir. L'orthographe et les accords doivent être irréprochables : une faute dans un CV est souvent éliminatoire.

\textbf{Étape 5 : mettre en page, relire, faire relire.} La mise en page doit être aérée : une page maximum, des rubriques clairement séparées, une seule police de caractères, des alignements rigoureux. Puis vient la relecture : d'abord à voix haute (l'oreille détecte ce que l'œil ne voit plus), ensuite par une autre personne — un camarade, un enseignant — qui repérera les coquilles et les ambiguïtés.

\begin{methode}[Adapter un CV à une offre d'emploi en 4 gestes]
\begin{enumerate}
\item \textbf{Lire l'offre} attentivement, en entier, deux fois : identifier le poste, les missions et le profil recherché.
\item \textbf{Repérer les mots-clés} : les compétences nommées (« maintenance », « comptabilité », « accueil client »), les diplômes exigés, les qualités demandées.
\item \textbf{Sélectionner} dans son inventaire personnel les informations qui répondent à ces mots-clés, et écarter ou réduire le reste.
\item \textbf{Reformuler} : reprendre dans le CV le vocabulaire de l'offre, adapter l'objectif professionnel et faire passer en premier la rubrique décisive.
\end{enumerate}
\end{methode}

\begin{attention}[Erreur fréquente : le CV « passe-partout »]
L'erreur la plus courante — et la plus sanctionnée, dans la vie comme à l'examen — consiste à envoyer \textbf{le même CV à toutes les offres}, sans aucune adaptation. Un recruteur reconnaît immédiatement un CV générique : l'objectif professionnel est vague (« recherche un emploi »), les compétences mises en avant ne correspondent pas au poste, les centres d'intérêt sont hors sujet. Un CV non adapté donne l'impression que le candidat ne s'est pas intéressé à l'entreprise. \textbf{Une offre = un CV adapté.}
\end{attention}

\courssub{Adapter son CV : un même candidat, deux offres différentes}

\begin{propriete}[Les trois qualités d'un bon CV]
Un bon CV est à la fois \textbf{complet} (il contient toutes les rubriques attendues : état civil, objectif, formation, expérience, compétences, centres d'intérêt), \textbf{sélectif} (il ne retient que les informations pertinentes pour le poste visé) et \textbf{ciblé} (il est adapté à une offre précise, dans son vocabulaire, son ordre de rubriques et son objectif professionnel). Un CV complet mais non ciblé est un mauvais CV.
\end{propriete}

Prenons un cas concret. Etoundi Marc, élève de Terminale F4 (Génie civil) au Lycée technique de Douala, dispose aussi de bonnes notions de comptabilité acquises lors d'un stage chez un commerçant. Il repère deux offres : \textbf{offre A} — « Entreprise BATIPRO recrute un technicien BTP, chantier à Bonabéri » ; \textbf{offre B} — « Cabinet FIDEX recrute un assistant comptable, Douala-Akwa ». Le même jeune homme ne présentera pas le même document.

\begin{popfigure}
\begin{center}
\begin{tikzpicture}[
  cell/.style={rectangle, draw=popline, fill=white, text width=4.55cm, align=left, inner sep=5pt, font=\scriptsize, minimum height=1.9cm},
  head/.style={rectangle, draw=popdark, fill=popblue, text=white, text width=4.55cm, align=center, inner sep=5pt, font=\small\bfseries, minimum height=0.8cm},
  lab/.style={rectangle, draw=popdark, fill=poporangeL, text width=2.6cm, align=center, inner sep=5pt, font=\scriptsize\bfseries, minimum height=1.9cm}
]
\node[lab] (l0) at (0,0) {\phantom{x}};
\node[head] (h1) at (4.05,0) {Offre A : technicien BTP (BATIPRO)};
\node[head] (h2) at (8.85,0) {Offre B : assistant comptable (FIDEX)};
\node[lab] (l1) at (0,-1.45) {Rubrique mise en premier};
\node[cell] (c11) at (4.05,-1.45) {Formation : Tle F4 Génie civil, Lycée technique de Douala};
\node[cell] (c12) at (8.85,-1.45) {Formation : Tle F4 + notions de comptabilité (stage)};
\node[lab] (l2) at (0,-3.4) {Expérience valorisée};
\node[cell] (c21) at (4.05,-3.4) {Stage de 2 mois sur chantier : coffrage, ferraillage, lecture de plans};
\node[cell] (c22) at (8.85,-3.4) {Stage chez un commerçant : tenue du journal de caisse, facturation};
\node[lab] (l3) at (0,-5.35) {Compétences mises en avant};
\node[cell] (c31) at (4.05,-5.35) {Lecture de plans, métré, sécurité sur chantier, maçonnerie};
\node[cell] (c32) at (8.85,-5.35) {Saisie comptable, tableur, classement de pièces, rigueur};
\node[lab] (l4) at (0,-7.3) {Objectif professionnel};
\node[cell] (c41) at (4.05,-7.3) {« Mettre mes compétences de technicien BTP au service de vos chantiers »};
\node[cell] (c42) at (8.85,-7.3) {« Contribuer à la tenue comptable de votre cabinet avec rigueur »};
\node[lab] (l5) at (0,-9.25) {Centres d'intérêt ajustés};
\node[cell] (c51) at (4.05,-9.25) {Bricolage, dessin technique, football (esprit d'équipe)};
\node[cell] (c52) at (8.85,-9.25) {Lecture, jeux de logique, calcul mental};
\end{tikzpicture}
\end{center}
\legende{Tableau d'adaptation : le même candidat (Etoundi Marc, Tle F4) face à deux offres différentes. Chaque rubrique est ajustée ; seul le fond d'informations réelles reste identique.}
\end{popfigure}

Ce tableau montre un point capital : adapter un CV, ce n'est \textbf{pas mentir}. Etoundi n'invente rien ; il \emph{sélectionne}, \emph{ordonne} et \emph{reformule} des informations vraies en fonction de ce que chaque recruteur attend. C'est exactement la propriété énoncée plus haut : complet, sélectif, ciblé.

\begin{exemplebox}[Offre « Technicien en froid et climatisation », Douala — CV adapté d'un élève de Tle F4]
\textbf{Offre} : « La société FRIGOCAM recrute un technicien en froid et climatisation à Douala. Profil : Baccalauréat technique F4 ou équivalent ; connaissance des installations frigorifiques ; rigueur, ponctualité ; permis B apprécié. »

\textbf{CV adapté (extrait)} :

\textbf{MVONDO Alain} — Douala, quartier Deïdo — Tél. : 6 90 00 00 00 — mvondo.alain@mail.com

\emph{Objectif professionnel} : technicien en froid et climatisation, disponible immédiatement.

\emph{Formation} : 2024-2025 : Baccalauréat technique F4 (en cours), Lycée technique de Douala-Bassa. 2022 : Probatoire technique F4.

\emph{Expérience} : juillet-septembre 2024 : stage d'entretien de chambres froides, marché central de Douala — dépannage, recharge en gaz, nettoyage des condenseurs.

\emph{Compétences} : installation et maintenance de groupes froids ; lecture de schémas électriques ; soudure ; conduite (permis B, 2024).

\emph{Langues} : français (courant), anglais (scolaire).

\emph{Centres d'intérêt} : mécanique, électronique, handball.
\end{exemplebox}

Observez comment chaque ligne du CV « répond » à un élément de l'offre : le diplôme F4 exigé figure en premier, le stage porte précisément sur le froid, le permis B « apprécié » est signalé, et même les centres d'intérêt (mécanique, électronique) restent cohérents avec le poste. Rien n'est laissé au hasard.

\courssub{Présenter son CV à l'oral}

Le programme prévoit que l'élève sache non seulement rédiger son CV, mais aussi le \textbf{présenter}. Cette compétence est évaluée en classe et devient décisive lors d'un entretien d'embauche réel. Trois exercices sont à maîtriser.

\textbf{Se présenter en deux minutes.} Le fameux « Parlez-moi de vous » du recruteur. La réponse suit un plan simple : qui je suis (nom, formation), ce que j'ai fait (stages, expériences marquantes), ce que je sais faire (deux ou trois compétences clés), ce que je veux (l'objectif, en lien avec le poste). Deux minutes, c'est court : il faut répéter, chronomètre en main, pour éviter les hésitations et les détails inutiles.

\textbf{Justifier ses choix.} Le recruteur — ou l'examinateur — peut demander : « Pourquoi avez-vous placé la formation avant l'expérience ? », « Pourquoi ce centre d'intérêt ? ». Il faut pouvoir expliquer sa démarche : « Je suis jeune diplômé, ma formation est mon atout principal ; je l'ai donc placée en premier. » C'est ici que le conditionnel étudié dans la section précédente devient utile : « J'ai pensé qu'il serait plus pertinent de mettre ce stage en avant. »

\textbf{Répondre aux questions.} Rester calme, écouter la question jusqu'au bout, répondre de manière concise et honnête. Une question piège classique : « Vous n'avez pas d'expérience professionnelle ? » — Réponse possible : « Je n'ai pas encore de longue expérience, mais mon stage chez FRIGOCAM m'a permis de me confronter au terrain, et j'apprends vite. »

\begin{exoresolu}[Rédiger un CV à partir d'une situation donnée]
\textbf{Énoncé} : « Vous êtes NGONO Sylvie, élève de Terminale technique au Lycée technique de Yaoundé. Vous répondez à l'annonce suivante : "La société AGROPLUS recrute un(e) secrétaire-comptable à Mbalmayo. Profil : Baccalauréat technique, maîtrise de l'outil informatique, sens de l'organisation." Rédigez votre CV. »
\tcblower
\textbf{Solution détaillée} : on applique les cinq étapes.

\emph{Étape 1 — Analyse de l'offre} : poste = secrétaire-comptable ; lieu = Mbalmayo ; mots-clés = Baccalauréat technique, informatique, organisation.

\emph{Étape 2 — Inventaire} : formation (Baccalauréat technique en préparation, Probatoire 2023) ; stage d'un mois à la mairie de Mbalmayo (accueil, classement) ; compétences (Word, Excel, dactylographie, français courant, anglais scolaire) ; qualités (ponctualité, discrétion).

\emph{Étape 3 — Ordre des rubriques} : jeune diplômée $\Rightarrow$ la Formation passe avant l'Expérience.

\emph{Étape 4 — Rédaction} (style télégraphique, verbes d'action) :

\textbf{NGONO Sylvie} — Yaoundé, quartier Nkolbisson — Tél. : 6 99 00 00 00 — ngono.sylvie@mail.com

\emph{Objectif professionnel} : secrétaire-comptable, disponible immédiatement à Mbalmayo.

\emph{Formation} : 2024-2025 : Baccalauréat technique (en cours), Lycée technique de Yaoundé. 2023 : Probatoire technique.

\emph{Expérience} : août 2024 : stage de secrétariat, mairie de Mbalmayo — accueil du public, classement des dossiers, saisie de courriers.

\emph{Compétences} : traitement de texte et tableur (Word, Excel) ; dactylographie ; organisation et classement de documents.

\emph{Langues} : français (courant), anglais (scolaire).

\emph{Centres d'intérêt} : lecture, chant choral, volley-ball.

\emph{Étape 5 — Mise en page et relecture} : une page, rubriques alignées, aucune faute d'orthographe ; relecture par un camarade avant remise.

\emph{Vérification finale} : chaque mot-clé de l'offre (informatique, organisation, Baccalauréat technique) trouve une réponse dans le CV. Le document est complet, sélectif et ciblé.
\end{exoresolu}

\courssub{Complément utile au Bac : l'articulation CV / lettre de motivation}

\begin{savaistu}[Le CV au Baccalauréat technique]
Au Baccalauréat technique, l'épreuve de production de texte peut proposer une \textbf{situation d'embauche} et demander la rédaction d'un CV complet à partir des éléments fournis par le sujet. Le correcteur évalue alors : le respect de la structure externe (rubriques présentes et ordonnées), la qualité de la rédaction (style télégraphique, orthographe), et surtout \textbf{l'adéquation entre le CV et la situation} — autrement dit, votre capacité d'adaptation. Un CV générique, même bien présenté, perd des points.
\end{savaistu}

Signalons enfin un complément fréquemment exigé : le CV est rarement envoyé seul. Il s'accompagne d'une \cle{lettre de motivation}, et certains sujets de production demandent les deux documents ensemble. L'articulation est simple à comprendre : le \textbf{CV} présente les faits (formation, expériences, compétences) de manière brève et télégraphique ; la \textbf{lettre} les commente, les met en valeur et exprime la motivation, en phrases rédigées — c'est là que le conditionnel de politesse (« Je souhaiterais vous rencontrer », « Je vous serais reconnaissant de bien vouloir examiner ma candidature ») et le subjonctif après certaines tournures (« Il est important que je précise... ») trouvent toute leur utilité. Les deux documents doivent être \textbf{cohérents} : mêmes informations, même ciblage, même vocabulaire repris de l'offre. La lettre ne répète pas le CV : elle l'éclaire.

\begin{aretenir}[L'essentiel de la section]
\begin{itemize}
\item Rédiger un CV suit cinq étapes : analyser l'offre, inventorier ses informations, ordonner les rubriques, rédiger (style télégraphique, verbes d'action), mettre en page et relire.
\item Un bon CV est complet, sélectif et ciblé : on adapte chaque rubrique à l'offre, sans jamais mentir.
\item Pour un jeune diplômé, la rubrique « Formation » vient en premier.
\item Savoir présenter son CV à l'oral : se présenter en deux minutes, justifier ses choix, répondre aux questions.
\item Au Bac, le CV peut être demandé seul ou avec la lettre de motivation ; les deux documents doivent être cohérents et adaptés à la situation donnée.
\end{itemize}
\end{aretenir}

\coursec{Intégration, évaluation et remédiation}

Dans la section précédente, nous avons appris à rédiger un CV et à l'adapter à une offre d'emploi précise : choix des rubriques, mise en valeur des compétences, emploi correct du subjonctif et du conditionnel. Mais savoir rédiger un CV en s'exerçant n'est pas encore savoir le produire en situation réelle, sous la pression du temps et devant un jury. C'est tout l'objet de cette dernière section : \cle{intégrer} les acquis de l'unité dans des tâches complexes, \cle{évaluer} les productions selon des critères explicites, puis \cle{remédier} aux difficultés observées. C'est la démarche complète d'apprentissage exigée par l'approche par les compétences du programme MINESEC.

\courssub{Activité d'intégration 1 : rédiger un CV complet à partir d'une situation donnée}

\begin{definition}[Activité d'intégration]
Une \cle{activité d'intégration} est une tâche complexe, proche d'une situation réelle de la vie, qui oblige l'élève à mobiliser \emph{ensemble} plusieurs savoirs et savoir-faire de l'unité (structure du CV, adaptation à une offre, correction langagière) pour produire un texte fonctionnel abouti.
\end{definition}

\textbf{Situation d'intégration proposée.}\par

Votre établissement reçoit l'offre de stage suivante :

\begin{exemplebox}[Offre d'emploi / de stage (support de l'activité)]
\textbf{ENTREPRISE CAMTEC SARL — Douala, Bonabéri.} Recrute pour un stage pré-professionnel de deux mois : un(e) assistant(e) de gestion, titulaire d'un Baccalauréat technique (comptabilité-gestion) ou d'un Probatoire avec une forte expérience. Compétences exigées : maîtrise du tableur, sens de l'organisation, bonne expression écrite et orale en français. Candidatures (CV + lettre de motivation) à déposer au plus tard le 15 juin.
\end{exemplebox}

\textbf{Consigne.} En une heure, rédigez le CV complet d'un candidat fictif (ou le vôtre) répondant à cette offre. Le CV doit tenir sur une page, respecter la structure externe et interne étudiées, et être adapté à l'offre.

\begin{methode}[Conduire l'activité d'intégration en cinq étapes]
\begin{enumerate}
\item \textbf{Analyser la situation} : souligner dans l'offre le poste, le diplôme exigé, les compétences attendues, le délai.
\item \textbf{Sélectionner les informations} : ne retenir que les rubriques et les expériences utiles au poste (pas de rubrique « loisirs » si elle n'apporte rien).
\item \textbf{Rédiger} : état civil, objectif, formation, expériences, compétences, centres d'intérêt pertinents ; verbes d'action au passé composé ou à l'infinitif.
\item \textbf{Adapter} : reprendre dans l'objectif et les compétences les mots-clés de l'offre (« gestion », « tableur », « organisation »).
\item \textbf{Relire} : accords, orthographe, modes (subjonctif après « il faut que », conditionnel pour la formulation prudente des souhaits), présentation aérée.
\end{enumerate}
\end{methode}

\begin{exoresolu}[Extrait de CV attendu (corrigé-type)]
Énoncé : rédigez les rubriques « Objectif » et « Compétences » du CV répondant à l'offre CAMTEC.
\tcblower
\textbf{Solution détaillée.}\par
\textbf{Objectif :} « Obtenir un stage pré-professionnel d'assistant de gestion afin de mettre en pratique ma formation en comptabilité et de développer mon sens de l'organisation au sein d'une entreprise dynamique. »\par
\textbf{Compétences :} « Maîtrise du tableur Excel (tableaux croisés, facturation) ; tenue de caisse et classement des pièces comptables ; rédaction de courriers administratifs en français correct ; aisance à l'oral acquise lors des exposés de classe. »\par
Chaque compétence reprend un mot-clé de l'offre (tableur, organisation, expression écrite et orale) : c'est cela, \cle{adapter} un CV. La formulation au conditionnel serait aussi possible dans l'objectif : « Je souhaiterais effectuer un stage... », marque de politesse et de prudence.
\end{exoresolu}

\courssub{Activité d'intégration 2 : présenter oralement son CV devant un jury simulé}

Le CV écrit ne suffit pas : à l'entretien, le recruteur demande souvent « Présentez-vous ». La classe se transforme alors en jury : trois ou quatre élèves jouent les recruteurs de CAMTEC, les autres présentent tour à tour leur CV en deux minutes, puis répondent à une question du jury.

\begin{methode}[Présenter son CV à l'oral en 2 minutes]
\begin{enumerate}
\item \textbf{Accroche (15 secondes)} : nom, diplôme préparé, et une phrase qui capte l'attention : « Je m'appelle..., je prépare un Baccalauréat technique en comptabilité-gestion et je me présente pour le stage d'assistant de gestion. »
\item \textbf{Parcours (40 secondes)} : formation et expériences, dans l'ordre chronologique inverse, sans tout détailler : retenir deux étapes fortes.
\item \textbf{Compétences (40 secondes)} : deux ou trois compétences directement liées à l'offre, illustrées par un fait concret (« J'ai tenu la caisse de la coopérative scolaire pendant un an »).
\item \textbf{Objectif et conclusion (25 secondes)} : ce que l'on apporterait à l'entreprise, formule de politesse : « Je vous remercie de votre attention et je reste à votre disposition. »
\end{enumerate}
\end{methode}

\begin{attention}[Erreur fréquente : réciter au lieu de valoriser]
Le piège classique est de \emph{réciter} le CV rubrique par rubrique, comme une liste : « Je suis né le..., j'ai fait la classe de... ». Le jury a le document sous les yeux ! Il faut au contraire \cle{commenter} et \cle{valoriser} : choisir deux ou trois éléments, les expliquer, montrer ce qu'ils prouvent. Autres pièges : regarder ses notes en permanence, parler trop vite, employer le conditionnel de prudence à l'excès (« je voudrais peut-être... ») au lieu d'affirmer calmement ses atouts.
\end{attention}

\courssub{Grille d'évaluation critériée}

Pour évaluer objectivement, on utilise une \cle{grille critériée} : chaque critère est coté, la somme donne une note sur 20. La grille ci-dessous est conforme aux habitudes du Baccalauréat technique.

\begin{exemplebox}[Barème type Bac technique (sur 20 points)]
\begin{center}
\begin{tabularx}{\linewidth}{|l|X|c|}
\hline
\rowcolor{popblueL} \textbf{Critère} & \textbf{Attendus} & \textbf{Points} \\
\hline
Structure externe & Titre, rubriques visibles, une page, mise en page aérée & 5 \\
\hline
Contenu et adaptation & Informations pertinentes, reprise des mots-clés de l'offre, objectif clair & 7 \\
\hline
Langue & Orthographe, accords, ponctuation, emploi correct du subjonctif et du conditionnel & 5 \\
\hline
Présentation orale & Posture, regard, voix, gestion du temps, réponses au jury & 3 \\
\hline
\textbf{Total} & & \textbf{20} \\
\hline
\end{tabularx}
\end{center}
\end{exemplebox}

\begin{popfigure}
\begin{tikzpicture}[x=1cm,y=0.62cm]
% axes
\draw[->,thick,popdark] (0,0) -- (11.5,0) node[below left,font=\small] {Note / 5};
\draw[->,thick,popdark] (0,0) -- (0,9.2);
\foreach \n in {0,1,2,3,4,5}{\draw[popdark] (2*\n,0.12) -- (2*\n,-0.12) node[below,font=\small] {\n};}
% barres (valeurs ramenées sur 5 pour le graphique)
\fill[popblue]   (0.3,7.6) rectangle (4.7,8.6);     % Structure : 4.4/5
\fill[poporange] (0.3,5.6) rectangle (5.23,6.6);   % Contenu : 6.9/7 = 4.93/5
\fill[popgreen]  (0.3,3.6) rectangle (4.6,4.6);     % Langue : 4.3/5
\fill[popgold]   (0.3,1.6) rectangle (3.97,2.6);    % Oral : 2.2/3 = 3.67/5
% labels critères
\node[right,font=\small\bfseries] at (4.9,8.1) {Structure externe : 4,4/5};
\node[right,font=\small\bfseries] at (5.43,6.1) {Contenu et adaptation : 6,9/7};
\node[right,font=\small\bfseries] at (4.8,4.1) {Langue : 4,3/5};
\node[right,font=\small\bfseries] at (4.17,2.1) {Présentation orale : 2,2/3};
% labels axe y
\node[left,font=\small] at (0,8.1) {Structure};
\node[left,font=\small] at (0,6.1) {Contenu};
\node[left,font=\small] at (0,4.1) {Langue};
\node[left,font=\small] at (0,2.1) {Oral};
\end{tikzpicture}
\legende{Grille critériée d'évaluation d'un CV : exemple de profil d'un candidat. Chaque critère est ramené à une échelle de 5 pour la lisibilité du graphique ; la note finale sur 20 est la somme des points réels (ici 4,4 + 6,9 + 4,3 + 2,2 = 17,8/20).}
\end{popfigure}

\courssub{Compte rendu collectif : points forts et difficultés observés}

Après les passages à l'oral, la classe, guidée par l'enseignant, dresse le bilan. Ce \cle{compte rendu collectif} est un moment d'apprentissage à part entière : on apprend aussi des erreurs des autres.

\begin{center}
\begin{tabularx}{\linewidth}{|X|X|}
\hline
\rowcolor{popgreenL} \textbf{Points forts observés} & \textbf{Difficultés fréquentes} \\
\hline
Rubriques claires et bien ordonnées ; objectifs adaptés à l'offre ; bonne gestion du temps à l'oral. & Confusion entre « expérience » et « formation » ; fautes d'accord du participe passé ; subjonctif oublié après « il faut que » ; récitation du CV au lieu de sa valorisation ; dates incohérentes. \\
\hline
\end{tabularx}
\end{center}

Le compte rendu doit être \emph{formulé de façon constructive} : on ne dit pas « ton CV est mauvais », mais « ta rubrique compétences gagnerait à reprendre les mots-clés de l'offre ».

\courssub{Remédiation ciblée}

La \cle{remédiation} consiste à retravailler précisément les points qui ont échoué, avec des exercices courts et ciblés.

\begin{exercice}[Remédiation 1 — Subjonctif et conditionnel]
Conjuguez les verbes entre parenthèses au mode qui convient :
\begin{enumerate}
\item Il faut que je (envoyer) mon CV avant le 15 juin.
\item Je souhaiterais que vous (étudier) ma candidature.
\item Si j'avais plus d'expérience, je (postuler) à un poste permanent.
\item Bien qu'il (être) jeune, ce candidat est compétent.
\end{enumerate}
\end{exercice}

\begin{corrige}[Exercice Remédiation 1]
\begin{enumerate}
\item Il faut que j'\textbf{envoie} (subjonctif présent, obligatoire après « il faut que »).
\item Je souhaiterais que vous \textbf{étudiiez} (subjonctif après un verbe de souhait ; « souhaiterais » est au conditionnel de politesse).
\item je \textbf{postulerais} (conditionnel présent dans la conséquence d'une hypothèse en « si + imparfait »).
\item Bien qu'il \textbf{soit} jeune (subjonctif après la conjonction de concession « bien que »).
\end{enumerate}
\end{corrige}

\begin{exercice}[Remédiation 2 — Correction orthographique et rubriques]
Corrigez les cinq erreurs de ce passage extrait d'un CV : « J'ai effectuer un stage de deux mois chez SONARA. Mes compétence en comptabilité et ma bonne maitrise du tableur serais un atout pour votre entreprise. »
\end{exercice}

\begin{corrige}[Exercice Remédiation 2]
« J'ai \textbf{effectué} (participe passé, non infinitif) un stage de deux mois chez SONARA. Mes \textbf{compétences} (pluriel) en comptabilité et ma bonne \textbf{maîtrise} (accent circonflexe) du tableur \textbf{seraient} (conditionnel pluriel, accord avec le sujet composé) un atout pour votre entreprise. » Remarque : « seraient un atout » est correct, mais à l'écrit comme à l'oral on préfère souvent l'affirmation : « sont un atout ».
\end{corrige}

Pour les rubriques défaillantes, la remédiation consiste à reprendre la fiche de structure interne (section 3) et à faire réécrire uniquement la rubrique fautive, en binôme avec un élève ayant réussi.

\courssub{Réinvestissement : rédiger son CV personnel réel}

Dernière étape, la plus importante : le \cle{réinvestissement}. Chaque élève rédige maintenant \emph{son propre} CV, avec ses informations authentiques, pour une demande de stage réelle (entreprise, administration, ONG de sa localité). Ce document, corrigé en classe, sera effectivement utilisable : c'est la finalité concrète de l'unité.

\begin{savaistu}[CV et lettre de motivation : le passeport de l'emploi au Cameroun]
De nombreuses entreprises camerounaises — la SONARA à Limbé, ENEO, les grandes PME industrielles et commerciales de Douala — exigent un CV accompagné d'une lettre de motivation dès le premier contact, avant même tout entretien. Un CV mal présenté ou fautif est souvent écarté sans que le candidat soit reçu. Savoir rédiger ce document en Terminale technique, c'est donc préparer concrètement son insertion professionnelle.
\end{savaistu}

\begin{aretenir}[Bilan de la section]
\begin{itemize}
\item L'activité d'intégration mobilise \emph{ensemble} tous les savoirs de l'unité dans une situation proche du réel.
\item La présentation orale du CV se construit en quatre temps : accroche, parcours, compétences, objectif ; on \emph{commente} et on \emph{valorise}, on ne récite pas.
\item L'évaluation critériée (structure 5 pts, contenu 7 pts, langue 5 pts, présentation 3 pts) rend la notation objective.
\item Le compte rendu collectif identifie points forts et difficultés ; la remédiation cible les lacunes (subjonctif, conditionnel, orthographe, rubriques).
\item Le réinvestissement final produit un CV personnel réel, utilisable pour une demande de stage.
\end{itemize}
\end{aretenir}

\newpage

\coursec{Activité d'intégration}

\courssub{Objectifs de l'activité}

À l'issue de cette activité d'intégration, l'élève doit être capable de :
\begin{itemize}
\item analyser une offre d'emploi réelle et en dégager les compétences exigées ;
\item rédiger un \cle{CV} d'une page respectant la \cle{structure externe} (mise en page, rubriques, ordre logique) et la \cle{structure interne} (état civil, formation, expériences, compétences, centres d'intérêt) ;
\item employer correctement le \cle{conditionnel de politesse} et le \cle{subjonctif} dans les formulations de l'objectif professionnel ;
\item adapter son CV à une situation d'embauche précise ;
\item présenter oralement son CV en deux minutes devant un jury ;
\item évaluer un CV à l'aide d'une grille critériée et formuler des axes de remédiation.
\end{itemize}

\courssub{Situation}

\textbf{Offre d'emploi publiée dans la presse locale.}

\begin{center}
\begin{tabularx}{\linewidth}{|l|X|}
\hline
\rowcolor{popblueL} \textbf{Rubrique} & \textbf{Contenu de l'offre} \\
\hline
Employeur & Société Anonyme des Brasseries du Cameroun (SABC), site industriel de Douala-Bassa \\
\hline
Poste & Technicien de maintenance industrielle (1 poste) \\
\hline
Profil exigé & Baccalauréat technique F3 (Mécanique) ou F4 (Électrotechnique) ; stage en milieu industriel souhaité ; maîtrise du français et de l'anglais ; sens de l'organisation et du travail en équipe \\
\hline
Missions & Maintenance préventive et curative des lignes de production ; diagnostic des pannes ; rédaction de rapports techniques \\
\hline
Dossier à fournir & Un CV d'une page et une lettre de motivation, à déposer au plus tard le 30 juin \\
\hline
\end{tabularx}
\end{center}

Tu es élève en classe de Terminale technique F3 au Lycée technique de Douala. Tu as effectué l'année dernière un stage de quatre semaines à la Chocolaterie du Cameroun (CHOCCAM), où tu as participé à l'entretien des machines d'emballage. Tu as obtenu ton Probatoire technique avec la mention « Assez bien ». En anglais, tu as suivi une formation de niveau intermédiaire au centre de langues de ton quartier. Ton professeur de français organise en classe une \textbf{simulation de recrutement} : tu dois rédiger ton CV d'une page adapté à cette offre, puis le présenter en deux minutes devant un jury composé de trois élèves et de l'enseignant, qui jouent le rôle des recruteurs de la SABC.

Le jury évalue chaque candidat à l'aide de la grille critériée suivante (sur 20 points) :

\begin{center}
\begin{tabularx}{\linewidth}{|X|c|}
\hline
\rowcolor{popblueL} \textbf{Critère d'évaluation} & \textbf{Note} \\
\hline
Structure externe : mise en page claire, rubriques visibles, une seule page & /4 \\
\hline
Structure interne : rubriques complètes et ordonnées, informations pertinentes & /4 \\
\hline
Adéquation du CV avec l'offre (compétences exigées mises en valeur) & /4 \\
\hline
Correction de la langue : subjonctif et conditionnel bien employés, orthographe & /4 \\
\hline
Présentation orale : clarté, durée respectée (2 minutes), réponses au jury & /4 \\
\hline
\textbf{Total} & \textbf{/20} \\
\hline
\end{tabularx}
\end{center}

\courssub{Consignes}

\begin{enumerate}
\item \textbf{Analyse de l'offre.} Relis l'offre de la SABC et relève dans un tableau à deux colonnes : (a) les compétences exigées par l'employeur ; (b) les éléments de ton profil qui y répondent. Quel élément de ton profil reste à valoriser malgré son apparente faiblesse ?
\item \textbf{Plan du CV.} Propose l'ordre des rubriques de ton CV en justifiant ton choix : pourquoi placer la rubrique « Formation » avant « Expérience professionnelle » dans ton cas ?
\item \textbf{Rédaction.} Rédige ton CV complet sur une page : en-tête (état civil et coordonnées), objectif professionnel rédigé au conditionnel de politesse, formation, expérience professionnelle, compétences (linguistiques et techniques), centres d'intérêt.
\item \textbf{Langue.} Dans ton objectif professionnel, emploie au moins une phrase au conditionnel de politesse et une phrase contenant un subjonctif après une expression de souhait. Souligne-les.
\item \textbf{Présentation orale.} Prépare une prise de parole de deux minutes : présente-toi, annonce ton projet professionnel, mets en valeur deux compétences en lien avec l'offre, conclus par une formule de politesse.
\item \textbf{Évaluation croisée.} En tant que membre du jury, note deux de tes camarades à l'aide de la grille critériée et formule pour chacun un point fort et un axe de remédiation.
\item \textbf{Compte rendu collectif.} Participe à la synthèse de la classe : quelles erreurs reviennent le plus souvent ? Quelles stratégies ont fait la différence ?
\end{enumerate}

\courssub{Ressources à mobiliser}

\begin{aretenir}[Les outils de la leçon à réutiliser]
\begin{itemize}
\item \textbf{Structure externe du CV} : une page, titre « Curriculum Vitae », en-tête avec identité et coordonnées, rubriques annoncées par des titres visibles, alignements réguliers, aucune phrase inutile.
\item \textbf{Structure interne du CV} : état civil et coordonnées ; objectif professionnel ; formation (du diplôme le plus récent au plus ancien) ; expériences professionnelles et stages (de la plus récente à la plus ancienne) ; compétences ; centres d'intérêt.
\item \textbf{Le conditionnel de politesse} : il atténue une demande ou exprime un souhait avec déférence : \emph{« Je souhaiterais intégrer votre équipe de maintenance. »}
\item \textbf{Le subjonctif} : il suit les verbes et expressions de souhait, de volonté, d'obligation ou de sentiment : \emph{« Il est essentiel que je mette mes compétences au service de votre entreprise. »}
\item \textbf{Présentation orale} : se présenter, annoncer le plan de sa prise de parole, parler distinctement, regarder le jury, respecter le temps imparti, conclure.
\end{itemize}
\end{aretenir}

\begin{attention}[Pièges fréquents à éviter]
\begin{itemize}
\item Ne confonds pas le \textbf{futur} et le \textbf{conditionnel de politesse} : on écrit \emph{« je voudrais »} et non \emph{« je voudrai »} ; \emph{« je serais honoré »} et non \emph{« je serai honoré »}.
\item Après \emph{« il faut que »}, \emph{« il est important que »}, \emph{« je souhaite que »}, le verbe se met au \textbf{subjonctif}, jamais à l'indicatif : \emph{« il faut que je sois »} et non \emph{« il faut que je suis »}.
\item Un CV ne doit jamais dépasser \textbf{une page} et ne doit contenir aucune faute d'orthographe : relis-toi ou fais relire.
\item N'invente pas de diplôme ni d'expérience : valorise ce qui est vrai, même si cela semble modeste.
\end{itemize}
\end{attention}

\courssub{Démarche et corrigé commenté}

\begin{exoresolu}[Recrutement à la SABC : rédiger et présenter son CV]
\textbf{Énoncé.} À partir de l'offre de la SABC et du profil donné dans la situation, rédige un CV d'une page adapté, puis prépare sa présentation orale de deux minutes.

\tcblower
\textbf{Étape 1 : analyse de l'offre (consigne 1).}

\begin{center}
\begin{tabularx}{\linewidth}{|X|X|}
\hline
\rowcolor{popblueL} \textbf{Compétences exigées par la SABC} & \textbf{Éléments de mon profil} \\
\hline
Baccalauréat technique F3 ou F4 & Élève en Terminale F3, Probatoire technique obtenu (mention Assez bien) \\
\hline
Stage en milieu industriel souhaité & Stage de 4 semaines à la CHOCCAM (entretien des machines d'emballage) \\
\hline
Maîtrise du français et de l'anglais & Français : langue de scolarité ; anglais : niveau intermédiaire (centre de langues) \\
\hline
Organisation, travail en équipe & Travaux de maintenance en binôme pendant le stage \\
\hline
\end{tabularx}
\end{center}

L'élément à valoriser malgré son apparente faiblesse est le \textbf{stage} : quatre semaines seulement, mais en milieu industriel réel, exactement ce que l'offre « souhaite ». Il faut donc le décrire avec des verbes d'action précis (\emph{diagnostiquer, démonter, graisser, rédiger un rapport}) plutôt que de le présenter comme une simple visite.

\textbf{Étape 2 : ordre des rubriques (consigne 2).}

Pour un jeune candidat sans longue expérience, on place la \textbf{Formation} avant l'Expérience professionnelle : c'est le diplôme (ici le Baccalauréat F3 en préparation) qui répond d'abord au profil exigé. L'ordre retenu est donc : en-tête ; objectif professionnel ; formation ; expérience professionnelle ; compétences ; centres d'intérêt.

\textbf{Étape 3 : rédaction du CV (consignes 3 et 4).}

\begin{center}
\begin{tabularx}{\linewidth}{|X|}
\hline
\rowcolor{popblueL} \textbf{CURRICULUM VITAE} \\
\hline
\textbf{ÉTAT CIVIL ET COORDONNÉES} \par
NOM : MBALLA ESSOMBA Jean-Claude \par
Né le 12 mars 2007 à Douala — Nationalité camerounaise \par
Adresse : Quartier Ndogpassi, Douala — Tél. : 6 90 00 00 00 \par
Courriel : jc.mballa@exemple.cm \\
\hline
\textbf{OBJECTIF PROFESSIONNEL} \par
Candidat au Baccalauréat technique F3, \emph{je souhaiterais} mettre mes compétences en mécanique industrielle au service de la SABC. \emph{Il est important que je contribue} à la fiabilité de vos lignes de production. \\
\hline
\textbf{FORMATION} \par
2025-2026 : classe de Terminale F3 (Mécanique), Lycée technique de Douala \par
2024-2025 : Probatoire technique F3, mention Assez bien \\
\hline
\textbf{EXPÉRIENCE PROFESSIONNELLE} \par
Juillet 2025 : stage de 4 semaines à la CHOCCAM (Douala) — entretien des machines d'emballage : diagnostic de pannes, graissage, remplacement de pièces, rédaction d'un rapport de stage \\
\hline
\textbf{COMPÉTENCES} \par
Techniques : lecture de plans, maintenance préventive, utilisation des outils d'atelier \par
Linguistiques : français (courant), anglais (intermédiaire) \par
Qualités : ponctualité, esprit d'équipe, sens de l'organisation \\
\hline
\textbf{CENTRES D'INTÉRÊT} \par
Football (capitaine de l'équipe du quartier), lecture de revues techniques \\
\hline
\end{tabularx}
\end{center}

\textbf{Commentaires.} Le CV tient sur une page ; les rubriques sont annoncées par des titres en gras ; les diplômes et expériences sont classés du plus récent au plus ancien ; chaque information répond à une exigence de l'offre. L'objectif professionnel contient bien un \textbf{conditionnel de politesse} (\emph{je souhaiterais}) et un \textbf{subjonctif} après une expression d'importance (\emph{il est important que je contribue}).

\textbf{Étape 4 : présentation orale de deux minutes (consigne 5).}

Plan de prise de parole : (1) salutations et présentation (nom, classe, établissement) — 20 secondes ; (2) annonce du projet professionnel — 20 secondes ; (3) deux compétences en lien avec l'offre : le stage à la CHOCCAM et le niveau d'anglais — 60 secondes ; (4) conclusion au conditionnel de politesse : \emph{« Je serais honoré de pouvoir mettre mes compétences au service de votre entreprise. Je vous remercie de votre attention. »} — 20 secondes.

\textbf{Étape 5 : évaluation et compte rendu (consignes 6 et 7).}

Le jury attribue les notes critère par critère. Exemple de compte rendu collectif : les points forts les plus fréquents sont la clarté des rubriques et la précision du stage ; les erreurs les plus fréquentes sont le dépassement d'une page, l'oubli des coordonnées, l'emploi du futur (\emph{je voudrai}) à la place du conditionnel de politesse (\emph{je voudrais}), et l'indicatif à la place du subjonctif après « il faut que ». La remédiation porte donc sur la reprise des conjugaisons et sur la reformulation des objectifs professionnels.
\end{exoresolu}

\courssub{Bilan de l'activité}

Cette activité d'intégration a permis de mettre en évidence plusieurs acquis essentiels de la leçon :

\begin{itemize}
\item Un CV n'est pas une simple liste d'informations : c'est un \cle{texte fonctionnel} qui doit \textbf{répondre point par point à une offre}. L'analyse préalable de l'offre conditionne le choix et l'ordre des rubriques.
\item La \cle{structure externe} (une page, rubriques visibles, mise en page aérée) est aussi importante que le contenu : un recruteur lit d'abord la forme avant le fond.
\item La langue joue un rôle décisif : le \cle{conditionnel de politesse} marque le respect du destinataire, et le \cle{subjonctif} exprime correctement le souhait et la volonté. Ces deux modes, étudiés dans la partie langue, trouvent ici un emploi concret.
\item La présentation orale prolonge le CV écrit : savoir se présenter en deux minutes, avec clarté et courtoisie, est une compétence professionnelle à part entière.
\item L'évaluation croisée et le compte rendu collectif ont montré que les difficultés sont partagées (dépassement de page, confusion futur/conditionnel, indicatif après « il faut que ») : la remédiation porte sur la relecture méthodique et la reprise des conjugaisons.
\end{itemize}

L'élève est désormais capable, sur la base d'une situation d'embauche réelle, de rédiger un CV adapté, de le présenter oralement et de justifier ses choix de rédaction : c'est exactement la compétence attendue au Baccalauréat technique et dans la vie professionnelle.

\newpage

\coursec{Méthodes et exercices résolus}

\coursec{Leçon 18 : Rédiger un CV}

\courssub{Méthode 1 : Identifier et analyser un CV (réception de texte fonctionnel)}

\begin{methode}[Lire et analyser un CV]
Pour analyser un CV à l'écrit comme à l'oral, suivre une démarche en cinq étapes :
\begin{enumerate}
\item \cle{Identifier} la nature du texte : c'est un \cle{texte fonctionnel} (il sert à obtenir un emploi, un stage), non un texte littéraire.
\item Repérer la \cle{structure externe} : en-tête (état civil, coordonnées), titre du poste visé, blocs (formation, expérience, compétences, centres d'intérêt), mise en page (colonnes, puces, photo éventuelle).
\item Repérer la \cle{structure interne} : ordre des rubriques (chronologique ou antichronologique), verbes d'action, noms de diplômes, dates, mots-clés du métier.
\item Évaluer la \cle{cohérence} : les rubriques correspondent-elles au poste visé ? Les dates se suivent-elles sans trou injustifié ?
\item Formuler un \cle{jugement argumenté} : dire si le CV est efficace et justifier par des éléments précis du document.
\end{enumerate}
Réflexe : toujours citer le document pour appuyer chaque observation.
\end{methode}

\begin{exoresolu}[Analyse d'un extrait de CV]
Énoncé : Voici l'extrait du CV de Mme Etoundi : « Formation : 2023 — Baccalauréat F3 (Génie civil), Lycée technique de Douala ; 2020 — Probatoire F3. Expérience : juin-septembre 2024 — stagiaire ouvrière, chantier SABC, Douala : coffrage, lecture de plans. » Questions : 1) Quel type de texte est-ce ? 2) Quelles rubriques apparaissent ? 3) Dans quel ordre sont présentées les dates ? Justifiez.
\tcblower
Solution détaillée :
\begin{enumerate}
\item \textbf{Nature du texte.} Il s'agit d'un \cle{CV} (curriculum vitae), c'est-à-dire un \cle{texte fonctionnel} : sa fonction n'est pas de divertir ni d'émouvoir, mais d'informer un employeur pour obtenir un emploi ou un stage. On le reconnaît à ses rubriques, à ses dates et à l'absence de phrases rédigées en récit.
\item \textbf{Rubriques présentes.} Deux rubriques de la structure interne apparaissent : la rubrique \emph{Formation} (diplômes obtenus, avec établissements et années) et la rubrique \emph{Expérience professionnelle} (poste occupé, lieu, tâches réalisées : « coffrage, lecture de plans »). Il manque ici l'en-tête (état civil, coordonnées), les compétences et les centres d'intérêt, qui compléteraient le document.
\item \textbf{Ordre des dates.} Dans la rubrique Formation, les dates vont de 2023 à 2020, donc du plus récent au plus ancien : c'est l'ordre \cle{antichronologique}. Cet ordre est le plus courant dans un CV car l'employeur lit d'abord le diplôme le plus élevé et le plus récent, information la plus pertinente pour le poste.
\end{enumerate}
\textbf{Conclusion :} cet extrait est un CV fonctionnel, organisé en rubriques datées selon l'ordre antichronologique, conforme aux normes du genre.
\end{exoresolu}

\courssub{Méthode 2 : Construire la structure externe d'un CV}

\begin{methode}[Organiser la présentation matérielle du CV]
La structure externe concerne la forme visible du document. Procéder ainsi :
\begin{enumerate}
\item Rédiger l'\cle{en-tête} : nom et prénom (en évidence), date et lieu de naissance, adresse, téléphone, adresse électronique, situation matrimoniale si utile.
\item Indiquer le \cle{titre} : le poste ou l'objectif visé (ex. « Technicien en froid et climatisation »).
\item Disposer les \cle{rubriques} en blocs séparés et titrés : Formation, Expérience professionnelle, Compétences, Langues, Centres d'intérêt.
\item Soigner la \cle{mise en page} : une page au maximum, puces, alignement des dates, police lisible, marges régulières ; la photo est facultative mais soignée si présente.
\item Relire pour l'\cle{orthographe} et la cohérence graphique : une faute dans un CV est éliminatoire.
\end{enumerate}
Formule à retenir : un CV tient sur \textbf{une page}, claire, aérée, sans rature.
\end{methode}

\begin{exoresolu}[Construire l'en-tête et le plan d'un CV]
Énoncé : À partir des informations suivantes, présentez l'en-tête et le plan des rubriques du CV d'Alain Mbarga, né le 12 mars 2005 à Yaoundé, demeurant à Mvan, téléphone 690 00 00 00, courriel alain.mbarga@mail.com, titulaire d'un Baccalauréat F4, qui postule un emploi de technicien comptable.
\tcblower
Solution détaillée :
\begin{enumerate}
\item \textbf{En-tête} (état civil et coordonnées, placé en haut, nom en évidence) :
\begin{center}
\textbf{MBARGA Alain}\\
Né le 12 mars 2005 à Yaoundé\\
Adresse : Mvan, Yaoundé — Tél. : 690 00 00 00\\
Courriel : alain.mbarga@mail.com
\end{center}
\item \textbf{Titre du CV} (le poste visé, placé juste sous l'en-tête) : \emph{« Technicien comptable »}. Ce titre montre immédiatement à l'employeur l'objectif du candidat.
\item \textbf{Plan des rubriques} (blocs titrés, dans un ordre logique pour un jeune diplômé) :
\begin{itemize}
\item \textbf{Formation} : Baccalauréat F4 (année, établissement) — ordre antichronologique ;
\item \textbf{Expérience professionnelle} : stages en cabinet comptable, avec dates et tâches ;
\item \textbf{Compétences} : tenue des livres, logiciels de comptabilité, tableur ;
\item \textbf{Langues} : français, anglais (niveau précisé) ;
\item \textbf{Centres d'intérêt} : éléments valorisants et véridiques.
\end{itemize}
\item \textbf{Mise en page} : une seule page, dates alignées à gauche, puces pour les tâches, aucune faute d'orthographe.
\end{enumerate}
\textbf{Conclusion :} le CV obtenu respecte la structure externe attendue : en-tête complet, titre de poste, rubriques claires et présentation soignée.
\end{exoresolu}

\begin{center}
\begin{tabularx}{\linewidth}{|l|X|X|}
\hline
\rowcolor{popblueL}
\textbf{Élément de la structure externe} & \textbf{Contenu attendu} & \textbf{Critères de réussite} \\
\hline
En-tête & Nom, prénom, naissance, adresse, téléphone, courriel & Complet, lisible, nom en gras/majuscules \\
\hline
Titre & Intitulé exact du poste visé & Présent, précis, adapté à l'annonce \\
\hline
Mise en page & Une page, police standard (Arial, Calibri 11), marges, puces & Aéré, aligné, sans faute, photo pro si incluse \\
\hline
Rubriques & Formation, Expérience, Compétences, Langues, Centres d'intérêt & Titrées, séparées, ordre logique \\
\hline
\end{tabularx}
\end{center}
\legende{Tableau récapitulatif : les composantes de la structure externe du CV}

\courssub{Méthode 3 : Remplir la structure interne (contenu des rubriques)}

\begin{methode}[Rédiger le contenu des rubriques]
\begin{enumerate}
\item \textbf{Formation} : diplôme + spécialité + établissement + année, du plus récent au plus ancien (ordre \cle{antichronologique}).
\item \textbf{Expérience} : pour chaque poste ou stage, donner la période, l'entreprise, le lieu, puis les tâches avec des \cle{verbes d'action} à l'infinitif ou au participe passé : « accueillir la clientèle », « saisie des factures », « installation et maintenance des moteurs ».
\item \textbf{Compétences} : savoir-faire techniques et qualités vérifiables (maîtrise d'un logiciel, d'une machine, d'une langue).
\item \textbf{Langues} : langue + niveau (ex. « Français : courant », « Anglais : B2 », « Fulfuldé : langue maternelle »).
\item \textbf{Centres d'intérêt} : 2 ou 3 éléments précis et honnêtes, en lien si possible avec le poste (éviter « lecture, cinéma » trop vagues).
\item \textbf{Vérifier} : aucune information fausse, aucun « trou » de dates inexpliqué, vocabulaire du métier correctement employé.
\end{enumerate}
Réflexe : chaque ligne doit répondre à la question de l'employeur : « que sait faire ce candidat ? »
\end{methode}

\begin{exoresolu}[Rédiger la rubrique « Expérience professionnelle »]
Énoncé : Rédigez la rubrique « Expérience professionnelle » du CV de Fatou, qui a effectué un stage de deux mois (juillet-août 2024) à l'hôtel Méridien de Kribi, où elle accueillait les clients, gérait les réservations et tenait la caisse.
\tcblower
Solution détaillée :
\begin{enumerate}
\item \textbf{Éléments obligatoires} à faire figurer : la période, le poste, l'entreprise, le lieu, les tâches.
\item \textbf{Choix de la formulation} : on utilise des \cle{verbes d'action} (infinitif ou substantifs d'action), jamais de longues phrases rédigées à la première personne, car le CV est un texte fonctionnel qui va à l'essentiel.
\item \textbf{Rédaction} :
\begin{center}
\textbf{Expérience professionnelle}\\
\textbf{Juillet-août 2024 — Hôtesse d'accueil (stage), Hôtel Méridien, Kribi}\\
• Accueil et orientation de la clientèle ;\\
• Gestion des réservations (téléphone et courriel) ;\\
• Tenue de la caisse et encaissements.
\end{center}
\item \textbf{Justification des choix} : la période est placée en premier car l'employeur vérifie d'abord la durée et la date de l'expérience ; les puces rendent la lecture rapide ; les verbes « accueil », « gestion », « tenue » nomment des compétences concrètes et vérifiables lors d'un entretien.
\end{enumerate}
\textbf{Conclusion :} la rubrique est complète (période, poste, entreprise, lieu, tâches), formulée en verbes d'action et présentée en puces, conformément aux normes du CV.
\end{exoresolu}

\begin{center}
\begin{tabularx}{\linewidth}{|l|X|X|}
\hline
\rowcolor{popgreenL}
\textbf{Rubrique} & \textbf{Contenu type (exemple)} & \textbf{Formulation attendue} \\
\hline
Formation & 2024 : Baccalauréat G1, Lycée technique de Garoua & Antichronologique, diplôme + série + établissement + année \\
\hline
Expérience & Juillet-sept. 2024 : Secrétaire (stage), Mairie de Maroua & Période, poste, structure, lieu + puces verbes d'action \\
\hline
Compétences & Traitement de texte (Word), Tableur (Excel), Classement archivage & Noms de logiciels, techniques, savoir-faire vérifiables \\
\hline
Langues & Français (courant), Anglais (B1), Fulfuldé (maternelle) & Langue + niveau CECRL ou description fonctionnelle \\
\hline
Centres d'intérêt & Volley-ball (club), Lecture (roman africain), Bénévolat Croix-Rouge & Précis, valorisants, véridiques \\
\hline
\end{tabularx}
\end{center}
\legende{Tableau récapitulatif : contenu type et formulation de la structure interne}

\courssub{Méthode 4 : Employer les valeurs du subjonctif et du conditionnel}

\begin{methode}[Choisir le bon mode dans les écrits fonctionnels]
\begin{enumerate}
\item Le \cle{subjonctif} exprime le souhait, l'obligation, la nécessité, la possibilité, le doute, le sentiment, la concession. Il suit notamment : \emph{vouloir que, souhaiter que, il faut que, bien que, pour que, avant que, afin que}. Ex. : « Je souhaite que ma candidature retienne votre attention. »
\item Le \cle{conditionnel} exprime l'hypothèse, la politesse atténuée, le souhait modéré, l'information non confirmée. Ex. : « Je \emph{voudrais} postuler… », « Je \emph{serais} disponible pour un entretien. »
\item Dans une lettre de motivation ou un courriel accompagnant un CV, employer le \textbf{conditionnel de politesse} pour formuler une demande (je voudrais, je souhaiterais, pourriez-vous) et le \textbf{subjonctif} après les verbes de souhait/volonté et les locutions comme « il faut que », « bien que », « pour que ».
\item Vérifier la \cle{concordance des temps} : après « bien que » + subjonctif présent ; après « je voudrais que » + subjonctif présent ; si le principal est au passé, le subjonctif passe à l'imparfait (ex. « Je voulais que vous examinassiez »).
\end{enumerate}
Formule à retenir : demande polie $\rightarrow$ conditionnel ; souhait, obligation, concession, finalité $\rightarrow$ subjonctif.
\end{methode}

\begin{exoresolu}[Compléter une lettre d'accompagnement de CV]
Énoncé : Complétez avec le mode et le temps qui conviennent, puis justifiez : 1) « Je (vouloir) …… postuler au poste de secrétaire. » 2) « Je souhaite que vous (examiner) …… ma candidature. » 3) « Bien que je (être) …… jeune diplômée, j'ai déjà une expérience de stage. » 4) « (Pouvoir)-vous …… me recevoir en entretien ? »
\tcblower
Solution détaillée :
\begin{enumerate}
\item « Je \textbf{voudrais} postuler… » : \cle{conditionnel présent} de politesse. Une demande directe (« je veux ») serait brutale ; le conditionnel atténue et respecte le destinataire.
\item « Je souhaite que vous \textbf{examiniez}… » : \cle{subjonctif présent}. Le verbe « souhaiter que » exprime un souhait portant sur autrui : il impose le subjonctif dans la subordonnée.
\item « Bien que je \textbf{sois} jeune diplômée… » : \cle{subjonctif présent}. La conjonction de concession « bien que » est toujours suivie du subjonctif ; « sois » est le subjonctif présent du verbe \emph{être} (que je sois, que tu sois, qu'il soit…).
\item « \textbf{Pourriez}-vous me recevoir… ? » : \cle{conditionnel présent} de politesse dans une demande interrogative ; l'indicatif « pouvez-vous » serait moins courtois dans un écrit professionnel.
\end{enumerate}
\textbf{Conclusion :} dans les écrits fonctionnels de candidature, le conditionnel marque la politesse de la demande et le subjonctif suit les verbes de souhait et les conjonctions de concession : ces deux modes construisent le ton respectueux attendu par l'employeur.
\end{exoresolu}

\begin{center}
\begin{tabularx}{\linewidth}{|l|X|X|}
\hline
\rowcolor{poppurpleL}
\textbf{Contexte / Structure} & \textbf{Mode / Temps} & \textbf{Exemple en situation de candidature} \\
\hline
Demande polie (je + verbe) & Conditionnel présent & Je \textbf{voudrais} / \textbf{souhaiterais} vous rencontrer. \\
\hline
Question polie (inversion) & Conditionnel présent & \textbf{Pourriez}-vous / \textbf{Serait}-il possible de… ? \\
\hline
Verbe de souhait/volonté + que & Subjonctif présent & Je \textbf{souhaite que} vous \textbf{receviez} mon dossier. \\
\hline
Obligation / Nécessité (il faut que) & Subjonctif présent & Il \textbf{faut que} je \textbf{joigne} mes diplômes. \\
\hline
Concession (bien que / quoique) & Subjonctif présent & \textbf{Bien que} je \textbf{sois} débutant, je suis motivé. \\
\hline
Finalité (pour que / afin que) & Subjonctif présent & Je joins mon CV \textbf{afin que} vous \textbf{puissiez} juger mes compétences. \\
\hline
\end{tabularx}
\end{center}
\legende{Tableau récapitulatif : choix du mode (subjonctif vs conditionnel) dans la candidature}

\courssub{Méthode 5 : Adapter un CV à une situation d'embauche}

\begin{methode}[Cibler un CV sur une offre d'emploi]
\begin{enumerate}
\item \textbf{Lire l'offre} et souligner les exigences : diplôme, compétences techniques, expérience, qualités demandées, mots-clés.
\item \textbf{Sélectionner} dans son parcours uniquement ce qui répond à ces exigences : on ne ment pas, mais on met en avant ce qui est pertinent (principe de pertinence).
\item \textbf{Réécrire le titre} du CV avec l'intitulé exact du poste de l'annonce.
\item \textbf{Réordonner} les rubriques : un jeune diplômé place la Formation d'abord ; un candidat expérimenté place l'Expérience d'abord.
\item \textbf{Reprendre les mots-clés} de l'annonce dans les rubriques (ex. si l'offre demande « maintenance des installations électriques », écrire précisément cette compétence si on la possède).
\item \textbf{Supprimer ou réduire} ce qui est hors sujet (centres d'intérêt sans rapport, expériences trop éloignées) tout en évitant les trous de dates inexpliqués.
\end{enumerate}
Réflexe : un même candidat possède autant de versions de son CV que de postes visés.
\end{methode}

\begin{exoresolu}[Adapter un CV à une offre]
Énoncé : Offre : « Garage moderne de Bafoussam recrute un mécanicien auto : Bac F2 exigé, maîtrise du diagnostic électronique, sens du contact client. » Le CV de Serge mentionne : Bac F2 (2023) ; stage de mécanique (2024) ; emploi de vendeur de légumes (2022) ; passion pour le football. Indiquez trois adaptations à faire et justifiez-les.
\tcblower
Solution détaillée :
\begin{enumerate}
\item \textbf{Analyse de l'offre} : trois exigences — le diplôme (Bac F2), une compétence technique (diagnostic électronique), une qualité relationnelle (contact client).
\item \textbf{Adaptation 1 — le titre} : écrire en titre « Mécanicien automobile », intitulé exact de l'annonce, pour montrer que le CV répond à ce poste précis.
\item \textbf{Adaptation 2 — mise en avant de la compétence clé} : dans la rubrique Compétences, écrire explicitement « diagnostic électronique des véhicules » (acquis pendant le stage de 2024), en reprenant le mot-clé de l'offre ; détailler le stage de 2024 avec des verbes d'action (« révision des moteurs, diagnostic électronique, accueil des clients »), ce qui prouve aussi le contact client.
\item \textbf{Adaptation 3 — tri des informations} : supprimer ou réduire l'emploi de vendeur de légumes (hors sujet technique) ; on peut toutefois le garder en une ligne dans « Autres expériences » pour ne pas laisser de trou de dates (2022-2023). Le football peut être conservé en centre d'intérêt s'il suggère l'esprit d'équipe et la forme physique, sinon supprimé au profit d'un intérêt technique (ex. « mécanique de compétition »).
\end{enumerate}
\textbf{Conclusion :} le CV adapté reprend le titre du poste, met en évidence le diplôme et la compétence exigés avec les mots de l'annonce, et élimine les informations non pertinentes : il devient une réponse ciblée à l'offre.
\end{exoresolu}

\courssub{Méthode 6 : Rédiger et présenter un CV complet (intégration)}

\begin{methode}[Produire un CV complet à partir d'une situation donnée]
\begin{enumerate}
\item \textbf{Dégager la situation} : qui est le candidat ? Quel poste vise-t-il ? Quelles informations fournit le sujet ?
\item \textbf{Classer les informations} en rubriques : état civil, formation, expérience, compétences, langues, centres d'intérêt.
\item \textbf{Rédiger} : en-tête complet, titre du poste, rubriques en ordre antichronologique, verbes d'action, puces.
\item \textbf{Mettre en page} : une page, alignement des dates, titres de rubriques en évidence, aucune faute.
\item \textbf{Présenter à l'oral} : se présenter en 1 à 2 minutes (identité, diplôme principal, expérience clé, motivation), répondre aux questions en utilisant le conditionnel de politesse.
\item \textbf{Auto-évaluer} avec la grille : structure externe respectée ? structure interne complète ? adaptation au poste ? correction de la langue ?
\end{enumerate}
\end{methode}

\begin{exoresolu}[Sujet d'intégration type Bac]
Énoncé : Vous êtes Aïcha Bello, née le 5 mai 2005 à Garoua, domiciliée à Maroua (tél. 699 11 22 33, courriel aicha.bello@mail.com). Titulaire du Baccalauréat G1 (2024, Lycée technique de Garoua), vous avez fait un stage de secrétariat à la mairie de Maroua (juillet-septembre 2024 : frappe de courriers, classement, accueil). Vous maîtrisez le traitement de texte et parlez français et fulfuldé. La SODECOTON recrute une secrétaire. Rédigez votre CV.
\tcblower
Solution détaillée :
\begin{enumerate}
\item \textbf{Analyse de la situation} : candidate jeune diplômée $\rightarrow$ rubrique Formation avant Expérience ; poste visé : secrétaire $\rightarrow$ mettre en avant stage de secrétariat, frappe, accueil, traitement de texte.
\item \textbf{CV rédigé} :
\begin{center}
\textbf{BELLO Aïcha}\\
Née le 5 mai 2005 à Garoua\\
Maroua — Tél. : 699 11 22 33 — Courriel : aicha.bello@mail.com\\[4pt]
\textbf{Secrétaire}\\[4pt]
\textbf{Formation}\\
2024 : Baccalauréat G1 (Techniques administratives), Lycée technique de Garoua.\\[4pt]
\textbf{Expérience professionnelle}\\
Juillet-septembre 2024 — Secrétaire (stage), Mairie de Maroua :\\
• frappe et mise en forme des courriers ;\\
• classement et archivage des dossiers ;\\
• accueil du public.\\[4pt]
\textbf{Compétences} : traitement de texte (Word), tableur (Excel), classement, accueil téléphonique et physique.\\
\textbf{Langues} : français (courant), fulfuldé (langue maternelle).\\
\textbf{Centres d'intérêt} : lecture (littérature africaine), couture (création de tenues).
\end{center}
\item \textbf{Justification des choix} : ordre antichronologique respecté ; verbes d'action en puces (« frappe », « classement », « accueil ») ; les compétences reprennent les besoins d'un poste de secrétaire ; la langue fulfuldé est un atout local pour la SODECOTON (zone cotonnière nord) ; le tout tient sur une page, sans faute.
\end{enumerate}
\textbf{Conclusion :} ce CV complet respecte la structure externe (en-tête, titre, rubriques, mise en page) et interne (ordre antichronologique, verbes d'action), et il est adapté au poste de secrétaire visé : il répond aux trois critères de l'évaluation au Baccalauréat technique.
\end{exoresolu}

\begin{aretenir}[Le texte iconique : la photo et la mise en page visuelle]
Le CV est aussi un \cle{texte iconique} : sa lisibilité repose sur des choix graphiques.
\begin{itemize}
\item \textbf{La photo} : facultative au Cameroun, mais si elle est présente, elle doit être professionnelle (fond neutre, tenue adaptée, format passeport, coin haut droit ou gauche). Une photo de loisir ou de mauvaise qualité décrédibilise le candidat.
\item \textbf{La hiérarchie visuelle} : noms de rubriques en gras/majuscules, dates alignées en colonne, puces pour les tâches, espaces blancs entre les blocs.
\item \textbf{La police} : sans empattement (Arial, Calibri, Helvetica), corps 11 ou 12 pour le texte, 14-16 pour le nom.
\item \textbf{La couleur} : utilisée avec parcimonie (un filet, un titre de rubrique en bleu marine ou gris foncé) pour guider l'œil sans distraire.
\end{itemize}
Un CV bien « lu » visuellement en 30 secondes a plus de chances d'être lu en détail.
\end{aretenir}

\begin{attention}[Les 5 erreurs qui coûtent des points au Bac]
\begin{enumerate}
\item \textbf{Confondre CV et lettre de motivation} : le CV est un texte fonctionnel en rubriques et en puces, sans phrases rédigées à la première personne ; la lettre, elle, est rédigée. Mélanger les deux fait perdre des points sur la structure externe.
\item \textbf{Oublier l'ordre antichronologique} : présenter les diplômes ou les expériences du plus ancien au plus récent est une faute de structure interne ; l'employeur lit d'abord le plus récent.
\item \textbf{Employer l'indicatif au lieu du conditionnel de politesse} : « je veux », « pouvez-vous » au lieu de « je voudrais », « pourriez-vous » donne un ton brutal, sanctionné dans les écrits de candidature.
\item \textbf{Mettre l'indicatif après « bien que » ou « il faut que »} : « bien que je suis » est fautif ; il faut le subjonctif : « bien que je \emph{sois} », « il faut que je \emph{sois} ».
\item \textbf{Négliger la relecture et l'adaptation} : fautes d'orthographe dans l'en-tête, coordonnées incomplètes, titre de poste absent, informations hors sujet ou inventées : autant de critères de la grille d'évaluation perdus. Un CV se relit toujours et s'adapte à chaque offre.
\end{enumerate}
\end{attention}

\newpage

\coursec{Exercices}

\courssub{Je vérifie mes connaissances}

\begin{exercice}[QCM : Nature et fonction du CV]
Parmi les propositions suivantes, une seule définit exactement le CV. Laquelle ?
\begin{enumerate}
\item Un texte narratif relatant la vie personnelle d'un candidat.
\item Un texte fonctionnel synthétique présentant le parcours professionnel et les compétences d'un candidat pour une candidature.
\item Une lettre de motivation détaillée expliquant les raisons de la candidature.
\item Un formulaire administratif obligatoire pour toute embauche au Cameroun.
\end{enumerate}
\end{exercice}

\begin{exercice}[Vrai / Faux justifié : Structure du CV]
Indiquez si les affirmations suivantes sont vraies ou fausses. Justifiez brièvement chaque réponse.
\begin{enumerate}
\item La structure externe du CV concerne uniquement le choix de la police de caractères et la couleur du papier.
\item La rubrique « Centres d'intérêt » est obligatoire et doit figurer en premier dans la structure interne.
\item L'ordre antéchronologique (du plus récent au plus ancien) est la norme pour présenter les expériences professionnelles et les diplômes.
\item La photographie d'identité est interdite sur un CV au Cameroun.
\end{enumerate}
\end{exercice}

\begin{exercice}[Définitions et terminologie]
Donnez la définition précise des termes suivants dans le contexte de la rédaction d'un CV :
\begin{enumerate}
\item Structure externe
\item Structure interne
\item Rubrique
\item Mots-clés (ou \emph{keywords})
\item CV ciblé (ou adapté)
\end{enumerate}
\end{exercice}

\begin{exercice}[Langue : Modes subjonctif et conditionnel]
Conjuguez les verbes entre parenthèses au mode et au temps appropriés (subjonctif présent ou conditionnel présent) pour exprimer la politesse, le souhait ou l'hypothèse.
\begin{enumerate}
\item Je souhaiterais que vous \textbf{(examiner)} ma candidature avec attention.
\item \textbf{(Avoir)} vous l'amabilité de me recevoir en entretien ?
\item Si mon profil \textbf{(correspondre)} à vos attentes, je me tiendrais à votre disposition.
\item Il est indispensable que le candidat \textbf{(maîtriser)} les logiciels de CAO.
\item Nous \textbf{(apprécier)} de pouvoir collaborer avec votre structure.
\end{enumerate}
\end{exercice}

\courssub{Je m'entraîne}

\begin{exercice}[Analyse de structure externe : Deux mises en page]
Voici la description schématique de deux CV (A et B).
\begin{itemize}
\item \textbf{CV A :} Police Comic Sans MS, taille 14, couleurs vives (rose/vert), marges très étroites, photo de vacances en grand format, rubriques mélangées (expériences entrecoupées de loisirs), 3 pages.
\item \textbf{CV B :} Police Arial/Calibri, taille 11, noir sur blanc, marges normales, photo d'identité professionnelle (3,5x4,5 cm) en haut à droite, rubriques clairement séparées par des interlignes et des titres en gras, ordre antéchronologique respecté, 1 page.
\end{itemize}
\begin{enumerate}
\item Identifiez 4 défauts majeurs de la structure externe du CV A.
\item Identifiez 4 qualités de la structure externe du CV B.
\item Quel CV un recruteur lira-t-il en priorité ? Justifiez en 2 lignes.
\end{enumerate}
\end{exercice}

\begin{exercice}[Rédaction de rubriques : Expériences et Compétences]
Vous êtes \textbf{Ngo Bitjek Jean}, titulaire d'un Baccalauréat Technique F3 (Génie Civil) obtenu en 2024 au Lycée Technique de Douala-Bonabéri. Vous avez effectué :
\begin{itemize}
\item Un stage de 3 mois (juillet-septembre 2023) à l'entreprise \emph{BTP Cameroun} : suivi de chantier, lecture de plans, gestion des approvisionnements.
\item Un job d'été (juillet-août 2022) comme magasinier à \emph{Quincaillerie du Port} : réception marchandises, inventaire, service client.
\end{itemize}
Vous maîtrisez : AutoCAD (niveau intermédiaire), Suite Office (Word, Excel avancé), lecture de plans d'exécution, français (langue maternelle), anglais (niveau B1), dibom (courant).
\begin{enumerate}
\item Rédigez la rubrique \textbf{« Expériences professionnelles »} (stages inclus) selon la structure interne standard (Date, Intitulé, Structure, Localisation, 3-4 activités principales par puce).
\item Rédigez la rubrique \textbf{« Compétences techniques et informatiques »} et \textbf{« Langues »}.
\end{enumerate}
\end{exercice}

\begin{exercice}[Correction de valeurs modales : Subjonctif / Conditionnel]
Les phrases suivantes, extraites de lettres de candidature ou d'entretiens simulés, contiennent des erreurs de conjugaison (mode/temps). Réécrivez chaque phrase correctement.
\begin{enumerate}
\item Si j'aurais le diplôme, je postulerais chez CAMTEL.
\item Il faut que je suis très motivé pour ce poste.
\item J'aimerais que vous me donnez une chance.
\item Il est nécessaire que le candidat connaît les normes de sécurité.
\item Nous voudrions que vous êtes disponible immédiatement.
\item Si je serais vous, j'accepterais l'offre.
\item Il faut qu'il a de l'expérience.
\item Je souhaiterai que vous m'invitez à un entretien.
\end{enumerate}
\end{exercice}

\begin{exercice}[Adaptation d'un CV à une offre]
Lisez l'offre d'emploi ci-dessous. À partir du profil de \textbf{Ngo Bitjek Jean} (Exercice 2), indiquez quels éléments du CV vous \textbf{mettez en avant} (gras, position haute, détails) et quels éléments vous \textbf{réduisez ou supprimez} pour coller à l'offre.

\textbf{Offre :} \emph{Entreprise « Routes \& Ouvrages d'Art » recherche un Technicien Bureau d'Études (BTP). Missions : Dessin AutoCAD, métrés, chiffrage, suivi dossiers d'appel d'offres. Profil : Bac+2/3 Génie Civil ou Bac Pro F3 avec expérience stage significative en BE. Maîtrise AutoCAD impérative. Anglais technique lu.}

\begin{enumerate}
\item Éléments à valoriser (3 minimum) :
\item Éléments à réduire/supprimer (2 minimum) :
\item Rédigez l'\textbf{Accroche (Profil / Objectif professionnel)} de 3 lignes maximum, placée en haut du CV, adaptée à cette offre.
\end{enumerate}
\end{exercice}

\courssub{Je me prépare au Bac}

\begin{exercice}[Sujet Type Bac : Rédaction complète d'un CV (1 page)]
\textbf{Situation :} Vous êtes \textbf{Mme Fouda Epse Mvondo Alice}, née le 12/03/2005 à Yaoundé. Vous êtes titulaire du \textbf{Baccalauréat Technique Série F4 (Électrotechnique)}, session 2024, mention Assez Bien, obtenu au Lycée Technique Industriel et Commercial de Yaoundé.
\textbf{Parcours :}
\begin{itemize}
\item Stage de fin d'études (3 mois, mars-mai 2024) à \textbf{ENEO} (Direction Régionale Centre) : Maintenance préventive/curative postes HT/BT, relevé de paramètres, rédaction de rapports techniques.
\item Stage de 1ère année (1 mois, juin 2022) à \textbf{CAMTEL} (Unité Lignes) : Tirage de câbles, épissure fibre optique, tests de continuité.
\item Job étudiant (vacances 2023) : Vendeuse en matériel électrique chez \textbf{Elec-Service} (Yaoundé, Mfoundi) : Conseil client, encaissement, rayonnage.
\end{itemize}
\textbf{Compétences :} Schémas unifilaires, Normes NFC 15-100, Multimètre/Oscilloscope, AutoCAD Electrical (notions), Word/Excel/PowerPoint, Maintenance préventive, Français (maternel), Anglais technique (lu/écrit B2), Ewondo (courant).
\textbf{Permis :} B (voiture) obtenu 2023.
\textbf{Centres d'intérêt :} Football féminin (club universitaire), Bénévolat Croix-Rouge (secourisme), Lecture technique.

\textbf{Consignes :}
\begin{enumerate}
\item Rédigez votre CV complet, tenant sur \textbf{une seule page A4}, structuré selon les normes professionnelles (Structure externe soignée, Structure interne complète et ordonnée).
\item Adaptez le CV pour une candidature au poste de \textbf{Technicien(ne) de Maintenance Électrique} chez \textbf{ENEO} (offre exigeant : rigueur, autonomie, connaissance normes sécurité, permis B, anglais technique).
\item Utilisez l'ordre antéchronologique. Soignez la présentation (rubriques, puces, alignements, police unique).
\end{enumerate}
\end{exercice}

\begin{exercice}[Sujet Type Bac : Analyse comparative et Langue]
\textbf{Document 1 : Extrait de CV de M. Essomba (Candidat 1)}
\textit{Rubrique Expériences :}
\begin{itemize}
\item 2023-2024 : Stagiaire Technicien Réseau - CAMTEL Yaoundé. Activités : Surveillance réseau, configuration routeurs, assistance utilisateurs.
\item 2022-2023 : Étudiant - Lycée Technique. Activités : Cours, devoirs, examens.
\item 2021 : Vendeur - Marché Central. Activités : Vente, argent, clients.
\end{itemize}
\textit{Rubrique Compétences :} Informatique (Word, Facebook, Jeux vidéo), Conduite, Football.
\textit{Présentation :} Police fantaisie, fond coloré, photo de groupe, 2 pages, fautes d'orthographe.

\textbf{Document 2 : Extrait de CV de Mlle Atangana (Candidate 2)}
\textit{Rubrique Expériences :}
\begin{itemize}
\item 03/2024 - 06/2024 : Stage Technicien Réseau - CAMTEL Yaoundé. Configuration routeurs Cisco (RIP, OSPF), Supervision réseau Nagios, Rédaction procédures, Support Niveau 1 utilisateurs (50 postes).
\item 06/2022 - 08/2022 : Stage Agent Accueil - CAMTEL Douala. Accueil physique/téléphonique, Orientation clients, Gestion réclamations, Reporting quotidien Excel.
\end{itemize}
\textit{Rubrique Compétences :} Réseaux : Cisco Packet Tracer, Nagios, Adressage IP/VLAN. Bureautique : Excel (TCD, RechercheV), Word, PowerPoint. Langues : Français (C1), Anglais technique (B2 - Certif. TOEIC 780). Permis B.
\textit{Présentation :} Police Calibri 11, sobriété, photo pro, 1 page, sans faute.

\textbf{Questions :}
\begin{enumerate}
\item \textbf{Structure externe :} Citez 3 défauts majeurs du CV de M. Essomba et 3 qualités majeures du CV de Mlle Atangana.
\item \textbf{Structure interne :} Analysez la rédaction des expériences (précision, verbes d'action, résultats) pour les deux candidats. Lequel est le plus convaincant ? Pourquoi ?
\item \textbf{Adéquation poste :} Pour un poste de \emph{Technicien Support Réseau Junior} chez CAMTEL, quel candidat retient le recruteur ? Justifiez par 3 arguments tirés des documents.
\item \textbf{Langue (Subjonctif/Conditionnel) :} M. Essomba écrit dans sa lettre : « \textit{Je voudrai que vous m'acceptez en stage. Si j'aurais la chance, je ferais de mon mieux.} » Corrigez les deux phrases en justifiant le mode utilisé.
\end{enumerate}
\end{exercice}

\begin{exercice}[Sujet Type Bac : Adaptation, Texte iconique et Oral]
\textbf{Profil unique :} \textbf{M. Kamga Paul}, Bac Tech G2 (Mécanique Auto) 2024. Stages : Garage Moderne (moteurs thermiques, diagnostic valise) 3 mois ; Concessionnaire Toyota (révision, électricité embarquée) 2 mois. Compétences : Diagnostic multimarque, Injection essence/diesel, Lecture schémas élec, Soudure, Outillage. Permis B, C. Anglais technique notice. Centres d'intérêt : Mécanique moto compétition, Tuning, Mécanique solidaire (asso).

\textbf{Offre A :} \emph{Garage « Rapid'Auto » (Centre-ville) : Mécanicien Polyvalent. Rapidité, relation client, polyvalence essence/diesel, diagnostic valise.}
\textbf{Offre B :} \emph{Usine « CamAssemblage » (Zone Industrielle) : Technicien Méthodes / Maintenance Préventive. Gammes de montage, plans, GMAO, rigueur, travail en équipe, horaires postés.}

\textbf{Questions :}
\begin{enumerate}
\item \textbf{Adaptation (Production) :} Produisez \textbf{deux versions de l'accroche (Profil)} et de la \textbf{rubrique Compétences} (ordonnancement des items) pour cibler l'Offre A puis l'Offre B. Présentez-les dans un tableau comparatif à deux colonnes.
\item \textbf{Texte iconique (Réception) :} Pour l'Offre B (Usine), le candidat choisit une mise en page : fond bleu marine très foncé, police blanche fine, icônes minimalistes, pas de photo, graphiques en barres pour les compétences (niveaux 80\%, 90\%), logo de l'entreprise visé en filigrane.
\begin{itemize}
\item Analysez ces choix graphiques (couleurs, typo, absence photo, graphiques, filigrane).
\item Que révèlent-ils de la stratégie du candidat et de sa compréhension de la culture d'entreprise industrielle ?
\end{itemize}
\item \textbf{Oral (Expression) :} Simulez l'oral de 2 minutes face au jury de l'Usine. Vous devez :
\begin{itemize}
\item Vous présenter (identité, diplôme, poste visé).
\item Mettre en avant 3 atouts majeurs pour \emph{cet} emploi précis.
\item Conclure par une phrase d'ouverture à l'entretien utilisant le \textbf{conditionnel de politesse}.
\end{itemize}
\textit{(Rédigez le texte intégral de votre intervention orale).}
\end{enumerate}
\end{exercice}

\newpage

\coursec{Corrigés des exercices}

\coursec{Exercices d'application et corrigés de la leçon VI}

\courssub{Exercices sur la nature et la structure du CV (Réception / Production)}

\begin{corrige}[Exercice 1 — QCM : Nature et fonction du CV]
\textbf{Réponse correcte : proposition 2.}

\textbf{Justification :}
\begin{itemize}
\item \textbf{Proposition 1 (fausse) :} le CV n'est pas un texte narratif ; il ne raconte pas la vie personnelle (récit de vie = biographie). Il sélectionne uniquement les informations utiles à une candidature.
\item \textbf{Proposition 2 (exacte) :} le CV est bien un \cle{texte fonctionnel} : il a une visée pratique (obtenir un entretien), il est \cle{synthétique} (une page idéalement) et présente le parcours (diplômes, expériences) et les compétences du candidat.
\item \textbf{Proposition 3 (fausse) :} elle décrit la \emph{lettre de motivation}, document distinct qui accompagne le CV mais ne s'y confond pas.
\item \textbf{Proposition 4 (fausse) :} le CV n'est pas un formulaire administratif imposé ; sa forme est libre (même si des normes professionnelles existent) et il n'est pas exigé par une loi pour toute embauche.
\end{itemize}
\end{corrige}

\begin{corrige}[Exercice 2 — Vrai / Faux justifié : Structure du CV]
\begin{enumerate}
\item \textbf{FAUX.} La structure externe concerne l'ensemble de la présentation matérielle : mise en page, marges, aération, choix et taille de la police, hiérarchie visuelle des titres, longueur (une page), qualité du papier. La police et la couleur n'en sont que deux éléments parmi d'autres.
\item \textbf{FAUX.} La rubrique « Centres d'intérêt » est \emph{facultative} ; elle ne figure jamais en premier. La structure interne débute par l'état civil / les coordonnées, puis éventuellement l'accroche, la formation et les expériences. Les centres d'intérêt, s'ils sont pertinents, se placent en fin de CV.
\item \textbf{VRAI.} L'ordre \cle{antéchronologique} (du plus récent au plus ancien) est la norme professionnelle : le recruteur s'intéresse d'abord à la situation actuelle ou la plus récente du candidat.
\item \textbf{FAUX.} La photographie n'est pas interdite au Cameroun ; elle est même courante. Elle doit toutefois être une photo d'identité \emph{professionnelle} (format type 3,5 $\times$ 4,5 cm, fond neutre, tenue correcte), et non une photo de vacances ou de groupe.
\end{enumerate}
\end{corrige}

\begin{corrige}[Exercice 3 — Définitions et terminologie]
\begin{enumerate}
\item \textbf{Structure externe :} l'organisation matérielle et visuelle du CV : mise en page, marges, police de caractères, taille des caractères, titres en gras, alignements, interlignes, longueur (une page), éventuelle photographie. C'est la « forme » du document.
\item \textbf{Structure interne :} l'organisation du contenu du CV en rubriques ordonnées : état civil et coordonnées, accroche (objectif professionnel), formation/diplômes, expériences professionnelles, compétences, langues, informations complémentaires (permis, centres d'intérêt).
\item \textbf{Rubrique :} chacune des sections thématiques du CV, annoncée par un titre visible (ex. « Formation », « Expériences professionnelles »), regroupant des informations de même nature.
\item \textbf{Mots-clés (\emph{keywords}) :} termes précis correspondant aux compétences et aux exigences de l'offre (ex. « AutoCAD », « métrés », « maintenance préventive »), que le recruteur — ou un logiciel de tri — repère en priorité dans un CV.
\item \textbf{CV ciblé (ou adapté) :} CV dont le contenu et l'ordonnancement des rubriques sont ajustés à une offre d'emploi précise : on y met en avant les expériences et compétences demandées, et on réduit ou supprime ce qui est hors sujet.
\end{enumerate}
\end{corrige}

\begin{corrige}[Exercice 5 — Analyse de structure externe : Deux mises en page]
\textbf{1. Quatre défauts majeurs du CV A :}
\begin{itemize}
\item Police fantaisiste (Comic Sans MS) et couleurs vives : manque de sobriété et de sérieux professionnel.
\item Marges très étroites et rubriques mélangées : document confus, illisible, sans hiérarchie de l'information.
\item Photo de vacances en grand format : inadaptée au contexte professionnel (il faut une photo d'identité sobre, de petit format).
\item Longueur de 3 pages : un CV de jeune diplômé doit tenir sur une page ; le recruteur consacre peu de temps à chaque dossier.
\end{itemize}
\textbf{2. Quatre qualités du CV B :}
\begin{itemize}
\item Police classique et lisible (Arial/Calibri, taille 11), noir sur blanc : sobriété professionnelle.
\item Photo d'identité professionnelle au format réglementaire, bien placée.
\item Rubriques clairement séparées (titres en gras, interlignes) : lecture rapide et hiérarchisée.
\item Ordre antéchronologique respecté et format d'une page : conformité aux normes du recrutement.
\end{itemize}
\textbf{3. CV lu en priorité :} le CV B. Sa présentation claire et aérée permet au recruteur de repérer en quelques secondes les informations essentielles ; le CV A, confus et fantaisiste, décourage la lecture et discrédite le candidat avant même l'étude du contenu.
\end{corrige}

\begin{corrige}[Exercice 6 — Rédaction de rubriques : Expériences et Compétences]
\textbf{1. Rubrique « Expériences professionnelles » (ordre antéchronologique) :}

\textbf{EXPÉRIENCES PROFESSIONNELLES}
\begin{itemize}
\item \textbf{Juillet -- septembre 2023 :} Stagiaire Technicien Génie Civil — \emph{BTP Cameroun}, Douala.
\begin{itemize}
\item Suivi de chantier et contrôle de l'avancement des travaux ;
\item Lecture et interprétation de plans d'exécution ;
\item Gestion des approvisionnements en matériaux.
\end{itemize}
\item \textbf{Juillet -- août 2022 :} Magasinier (job d'été) — \emph{Quincaillerie du Port}, Douala.
\begin{itemize}
\item Réception et contrôle des marchandises ;
\item Réalisation des inventaires de stock ;
\item Accueil et service à la clientèle.
\end{itemize}
\end{itemize}

\textbf{2. Rubriques « Compétences » et « Langues » :}

\textbf{COMPÉTENCES TECHNIQUES ET INFORMATIQUES}
\begin{itemize}
\item DAO/CAO : AutoCAD (niveau intermédiaire) ;
\item Bureautique : Word, Excel (niveau avancé), PowerPoint ;
\item Lecture de plans d'exécution, suivi de chantier.
\end{itemize}

\textbf{LANGUES}
\begin{itemize}
\item Français : langue maternelle ;
\item Anglais : niveau B1 (intermédiaire) ;
\item Dibom : courant.
\end{itemize}

\textbf{Points de vigilance :} dates en premier et au plus récent d'abord ; intitulé du poste en gras ; nom de la structure et localisation ; activités formulées par des verbes d'action ou des noms d'action, en puces.
\end{corrige}

\courssub{Exercices de langue : valeurs du subjonctif et du conditionnel}

\begin{corrige}[Exercice 4 — Langue : Modes subjonctif et conditionnel]
\begin{enumerate}
\item Je souhaiterais que vous \textbf{examiniez} ma candidature avec attention. (Subjonctif présent après « souhaiter que » ; « examiniez », 2\textsuperscript{e} pers. du pluriel.)
\item \textbf{Auriez}-vous l'amabilité de me recevoir en entretien ? (Conditionnel présent de politesse.)
\item Si mon profil \textbf{correspondait} à vos attentes, je me tiendrais à votre disposition. (Imparfait de l'indicatif après « si » exprimant l'hypothèse ; jamais de conditionnel après « si ».)
\item Il est indispensable que le candidat \textbf{maîtrise} les logiciels de CAO. (Subjonctif présent après « il est indispensable que ».)
\item Nous \textbf{apprécierions} de pouvoir collaborer avec votre structure. (Conditionnel présent de politesse / atténuation.)
\end{enumerate}
\textbf{Point de vigilance :} retenir la correspondance « si + imparfait $\rightarrow$ conditionnel présent » et l'emploi du subjonctif après les verbes de souhait et les expressions de nécessité.
\end{corrige}

\begin{corrige}[Exercice 7 — Correction de valeurs modales : Subjonctif / Conditionnel]
\begin{enumerate}
\item Si j'\textbf{avais} le diplôme, je postulerais chez CAMTEL. (Après « si » exprimant l'hypothèse : imparfait, jamais le conditionnel.)
\item Il faut que je \textbf{sois} très motivé pour ce poste. (Subjonctif présent d'\emph{être} après « il faut que ».)
\item J'aimerais que vous me \textbf{donniez} une chance. (Subjonctif présent après « aimer que ».)
\item Il est nécessaire que le candidat \textbf{connaisse} les normes de sécurité. (Subjonctif présent après « il est nécessaire que ».)
\item Nous voudrions que vous \textbf{soyez} disponible immédiatement. (Subjonctif présent d'\emph{être} après « vouloir que ».)
\item Si j'\textbf{étais} vous, j'accepterais l'offre. (Imparfait après « si » ; locution figée « si j'étais vous ».)
\item Il faut qu'il \textbf{ait} de l'expérience. (Subjonctif présent d'\emph{avoir}.)
\item Je \textbf{souhaiterais} que vous m'\textbf{invitiez} à un entretien. (Conditionnel de politesse « souhaiterais », puis subjonctif présent « invitiez ».)
\end{enumerate}
\textbf{Point de vigilance :} deux pièges classiques à l'examen — le conditionnel interdit après « si » hypothétique, et l'indicatif interdit après « il faut que » / les verbes de volonté.
\end{corrige}

\courssub{Exercices d'adaptation et de rédaction (Production)}

\begin{corrige}[Exercice 8 — Adaptation d'un CV à une offre]
\textbf{1. Éléments à valoriser (mots-clés de l'offre : AutoCAD, BE, génie civil) :}
\begin{itemize}
\item La maîtrise d'\textbf{AutoCAD} : la faire passer en tête des compétences, en gras, en précisant « dessin 2D, plans d'exécution » ;
\item Le stage à \emph{BTP Cameroun} : détailler les activités proches du bureau d'études (lecture de plans, suivi technique, approvisionnements = début du chiffrage) ;
\item Le diplôme \textbf{Bac Technique F3 Génie Civil} : mentionné dès l'accroche, car l'offre accepte ce profil « avec expérience stage significative » ;
\item L'anglais : préciser « anglais technique lu » pour répondre à l'exigence de l'offre.
\end{itemize}
\textbf{2. Éléments à réduire ou supprimer :}
\begin{itemize}
\item Le job de magasinier : le réduire à une ligne (ou le supprimer si la place manque), car sans rapport avec le poste ;
\item Les compétences trop générales (Word, PowerPoint) : les abréger au profit d'Excel (utile pour les métrés et chiffrages).
\end{itemize}
\textbf{3. Accroche (3 lignes maximum) :}

\emph{« Jeune titulaire du Baccalauréat Technique F3 (Génie Civil), formé à AutoCAD et à la lecture de plans lors d'un stage en entreprise de BTP. Rigoureux et disponible, je souhaite mettre mes compétences en dessin et en suivi technique au service de votre bureau d'études. »}
\end{corrige}

\courssub{Sujets type Baccalauréat / Probatoire (Intégration)}

\begin{corrige}[Exercice 9 — Sujet Type Bac : Rédaction complète d'un CV]
\textbf{Corrigé type (contenu attendu ; la mise en page réelle sera soignée sur une page A4) :}

\begin{center}
\textbf{FOUDA EPSE MVONDO Alice}\\
Née le 12/03/2005 à Yaoundé — Yaoundé, Cameroun\\
Tél. : 6XX XX XX XX — Courriel : alice.fouda@exemple.cm — Permis B
\end{center}

\textbf{PROFIL}\\
\emph{Jeune technicienne en électrotechnique (Bac Technique F4, mention Assez Bien), formée à la maintenance HT/BT chez ENEO. Rigoureuse, autonome et familiarisée avec les normes de sécurité, je candidate au poste de Technicienne de Maintenance Électrique.}

\textbf{FORMATION}
\begin{itemize}
\item \textbf{2024 :} Baccalauréat Technique F4 (Électrotechnique), mention Assez Bien — Lycée Technique Industriel et Commercial de Yaoundé.
\end{itemize}

\textbf{EXPÉRIENCES PROFESSIONNELLES}
\begin{itemize}
\item \textbf{Mars -- mai 2024 :} Stagiaire Technicienne — \emph{ENEO}, Direction Régionale Centre, Yaoundé. Maintenance préventive et curative de postes HT/BT ; relevé de paramètres électriques ; rédaction de rapports techniques.
\item \textbf{Vacances 2023 :} Vendeuse en matériel électrique — \emph{Elec-Service}, Yaoundé. Conseil client ; encaissement ; rayonnage.
\item \textbf{Juin 2022 :} Stagiaire — \emph{CAMTEL}, Unité Lignes. Tirage de câbles ; épissure de fibre optique ; tests de continuité.
\end{itemize}

\textbf{COMPÉTENCES}
\begin{itemize}
\item Lecture de schémas unifilaires ; normes NF C 15-100 ; maintenance préventive ;
\item Mesures : multimètre, oscilloscope ; AutoCAD Electrical (notions) ;
\item Bureautique : Word, Excel, PowerPoint.
\end{itemize}

\textbf{LANGUES} : Français (maternel) ; Anglais technique (lu/écrit, B2) ; Ewondo (courant).

\textbf{INFORMATIONS COMPLÉMENTAIRES} : Permis B (2023) ; secourisme (bénévolat Croix-Rouge) ; football féminin en club universitaire ; lecture technique.

\textbf{Adaptation à l'offre ENEO :} le stage ENEO est placé en tête et détaillé ; les mots-clés de l'offre (« maintenance », « normes », « autonomie », « permis B », « anglais technique ») figurent dans le profil et les compétences ; le job de vendeuse est réduit à une ligne.

\textbf{Barème indicatif (sur 20) :} structure externe soignée (une page, rubriques, sobriété) : 4 pts ; structure interne complète et ordonnée (état civil, profil, formation, expériences antéchronologiques, compétences, langues) : 8 pts ; adaptation à l'offre (mots-clés, mise en avant du stage ENEO) : 4 pts ; correction de la langue : 4 pts.

\textbf{Points de vigilance :} ordre antéchronologique respecté ; verbes d'action ; aucune information fantaisiste ; une seule page ; cohérence entre le profil annoncé et l'offre.
\end{corrige}

\begin{corrige}[Exercice 10 — Sujet Type Bac : Analyse comparative et Langue]
\textbf{1. Structure externe.}
\begin{itemize}
\item Défauts de M. Essomba : police fantaisiste et fond coloré (manque de sérieux) ; photo de groupe (non professionnelle) ; 2 pages avec fautes d'orthographe (négligence rédhibitoire).
\item Qualités de Mlle Atangana : police Calibri 11 sobre et lisible ; photo professionnelle ; une page sans faute (rigueur et respect du lecteur).
\end{itemize}
\textbf{2. Structure interne.} M. Essomba rédige des expériences vagues (« Vente, argent, clients »), sans dates précises, sans verbes d'action ni résultats ; il mélange études et expériences et cite des « compétences » hors sujet (Facebook, jeux vidéo). Mlle Atangana date précisément chaque expérience (03/2024 -- 06/2024), utilise des verbes d'action et des données techniques (configuration routeurs Cisco RIP/OSPF, supervision Nagios) et quantifie ses résultats (support de 50 postes). \textbf{Mlle Atangana est la plus convaincante} : son CV prouve les compétences au lieu de les affirmer.
\textbf{3. Adéquation au poste.} Le recruteur retient \textbf{Mlle Atangana} : (a) ses stages correspondent exactement au poste (réseau, support utilisateurs) ; (b) ses compétences techniques sont précises et certifiées (TOEIC 780, Excel TCD) ; (c) sa présentation irréprochable témoigne de la rigueur attendue d'un technicien support.
\textbf{4. Langue.} Phrase correcte : « Je \textbf{voudrais} que vous m'\textbf{acceptiez} en stage. Si j'\textbf{avais} la chance, je \textbf{ferais} de mon mieux. » Justification : conditionnel présent de politesse (\emph{voudrais}) ; subjonctif présent (\emph{acceptiez}) après « vouloir que » ; imparfait (\emph{avais}) après « si » hypothétique, suivi du conditionnel (\emph{ferais}) dans la proposition principale.

\textbf{Barème indicatif (sur 20) :} structure externe : 4 pts ; analyse des expériences : 5 pts ; adéquation au poste (3 arguments tirés des documents) : 6 pts ; correction modale justifiée : 5 pts.

\textbf{Points de vigilance :} les arguments doivent être \emph{tirés des documents} (citations ou références précises) ; la justification grammaticale (nom du mode et règle d'emploi) est exigée, pas seulement la correction.
\end{corrige}

\begin{corrige}[Exercice 11 — Sujet Type Bac : Adaptation, Texte iconique et Oral]
\textbf{1. Tableau comparatif des deux versions ciblées :}

\begin{center}
\begin{tabularx}{\linewidth}{|X|X|}
\hline
\rowcolor{popblueL} \textbf{Version Offre A — Rapid'Auto (garage)} & \textbf{Version Offre B — CamAssemblage (usine)} \\
\hline
\textbf{Accroche :} « Mécanicien polyvalent, Bac Technique G2, rapide et à l'écoute du client. Expérience en diagnostic valise multimarque et réparation essence/diesel acquise en garage et en concession. » & \textbf{Accroche :} « Technicien mécanique, Bac Technique G2, rigoureux et habitué au travail en équipe. Formé à la lecture de plans et schémas, je vise le poste de Technicien Méthodes / Maintenance préventive. » \\
\hline
\textbf{Compétences (ordre) :} 1. Diagnostic valise multimarque ; 2. Injection essence et diesel ; 3. Révision et électricité embarquée ; 4. Relation client ; 5. Soudure, outillage. & \textbf{Compétences (ordre) :} 1. Lecture de plans et schémas ; 2. Maintenance préventive et gammes d'entretien ; 3. Diagnostic et contrôle qualité ; 4. Soudure, outillage ; 5. Travail en équipe, respect des horaires postés. \\
\hline
\end{tabularx}
\end{center}

\textbf{2. Analyse du texte iconique (Offre B).}
\begin{itemize}
\item \emph{Couleurs et typographie :} le fond bleu marine et la police blanche fine évoquent la technicité et le sérieux industriel, mais nuisent à la lisibilité et à l'impression (un CV doit rester imprimable en noir et blanc) : choix esthétique risqué.
\item \emph{Absence de photo :} pertinente dans l'industrie, où l'on évalue les compétences plutôt que l'apparence ; elle évite aussi toute discrimination.
\item \emph{Graphiques en barres (80 \%, 90 \%) :} ils rendent les compétences visibles d'un coup d'œil, mais les pourcentages sont subjectifs et invérifiables ; mieux vaut des niveaux explicites (débutant, intermédiaire, avancé).
\item \emph{Logo de l'entreprise en filigrane :} il montre la personnalisation du CV pour cette entreprise précise, mais peut paraître présomptueux (usage du logo sans autorisation) : à manier avec prudence.
\end{itemize}
\emph{Stratégie révélée :} le candidat a compris que l'entreprise industrielle valorise la technicité, la sobriété et l'efficacité visuelle ; il cherche à montrer qu'il « parle le langage » de l'usine. Sa démarche de ciblage est bonne, mais certains choix (lisibilité, filigrane) doivent être corrigés.

\textbf{3. Texte intégral de l'intervention orale (2 minutes) :}

« Messieurs et Mesdames les membres du jury, bonjour. Je me présente : Kamga Paul, titulaire du Baccalauréat Technique G2 en Mécanique Automobile, session 2024. Je candidate au poste de Technicien Méthodes et Maintenance préventive au sein de votre usine.

Trois atouts me semblent répondre directement à vos besoins. Premièrement, ma formation technique et mes stages en garage et en concession m'ont appris à lire des plans et des schémas électriques et à travailler avec méthode, compétences indispensables pour établir vos gammes de montage. Deuxièmement, je suis rigoureux et habitué au travail en équipe : mes expériences m'ont familiarisé avec les procédures, les consignes de sécurité et le respect des délais, ce qui correspond à votre organisation en horaires postés. Troisièmement, je maîtrise le diagnostic et la maintenance des équipements, et je suis prêt à me former rapidement à votre logiciel de GMAO.

Je vous remercie de votre attention. Je serais heureux que vous me donniez l'occasion de vous exposer plus en détail ma motivation, et je me tiendrais à votre disposition pour un essai sur poste. »

\textbf{Barème indicatif (sur 20) :} tableau comparatif adapté aux deux offres : 6 pts ; analyse iconique argumentée (au moins 4 choix commentés) : 6 pts ; oral complet (présentation, 3 atouts ciblés, conclusion au conditionnel de politesse) : 6 pts ; langue correcte : 2 pts.

\textbf{Points de vigilance :} l'accroche doit citer le poste visé et reprendre les mots-clés de l'offre ; l'ordre des compétences change selon la cible ; la phrase finale doit employer le conditionnel de politesse (\emph{je serais heureux}, \emph{je me tiendrais}) et le subjonctif après « souhaiter que » si ce verbe est utilisé.
\end{corrige}

\newpage

\coursec{Fiche bilan}

\begin{popfigure}
\begin{tikzpicture}[mindmap, grow cyclic, every node/.style={concept, text=white, font=\small\bfseries},
root concept/.append style={concept color=popdark, font=\large\bfseries},
level 1/.append style={level distance=3.6cm, sibling angle=51.4},
level 2/.append style={level distance=2.2cm, sibling angle=40, font=\scriptsize}]
\node[root concept]{Rédiger un CV}
  child[concept color=popblue]{ node {Texte fonctionnel}
    child { node {But : obtenir un entretien} }
    child { node {Lecture rapide du recruteur} }
  }
  child[concept color=poporange]{ node {Structure externe}
    child { node {1 page, aéré, titres visibles} }
    child { node {Photo et mise en page} }
  }
  child[concept color=popgreen]{ node {Structure interne}
    child { node {État civil, formation, expérience} }
    child { node {Compétences, centres d'intérêt} }
  }
  child[concept color=poppurple]{ node {Langue}
    child { node {Subjonctif : souhait, doute, obligation} }
    child { node {Conditionnel : politesse, hypothèse} }
  }
  child[concept color=poppink]{ node {Rédaction}
    child { node {Ordre antichronologique} }
    child { node {Vérité et précision} }
  }
  child[concept color=popgold]{ node {Adaptation}
    child { node {Cibler l'offre d'emploi} }
    child { node {Mots-clés du poste} }
  }
  child[concept color=popmuted]{ node {Présentation}
    child { node {Relire, corriger} }
    child { node {Savoir le commenter à l'oral} }
  };
\end{tikzpicture}
\legende{Carte mentale de la leçon : le CV, un texte fonctionnel à maîtriser pour le Probatoire et le Baccalauréat technique.}
\end{popfigure}

\begin{bilan}
\textbf{Définitions clés}
\begin{itemize}
\item Le \cle{CV} (curriculum vitae) : texte \cle{fonctionnel} qui présente de manière brève et ordonnée le parcours, les compétences et l'identité d'un candidat, en vue d'obtenir un emploi, un stage ou une formation.
\item \cle{Structure externe} : l'aspect matériel du document (format, mise en page, blocs, photo).
\item \cle{Structure interne} : le contenu informationnel organisé en rubriques.
\item Texte \cle{iconique} : élément visuel (photo, logo, pictogramme) qui accompagne ou remplace l'écrit.
\end{itemize}

\textbf{Les rubriques du CV (structure interne)}
\begin{center}
\begin{tabularx}{\linewidth}{|l|X|}\hline
\rowcolor{popblueL} \textbf{Rubrique} & \textbf{Contenu attendu} \\ \hline
État civil / identité & Nom, prénom, date et lieu de naissance, nationalité, contacts (téléphone, courriel, adresse) \\ \hline
Objectif professionnel & Une phrase ciblée sur le poste visé \\ \hline
Formation / diplômes & Du plus récent au plus ancien (ordre \cle{antichronologique}) \\ \hline
Expérience professionnelle & Postes, stages, dates, tâches réalisées \\ \hline
Compétences & Techniques (informatique, métier) et linguistiques \\ \hline
Centres d'intérêt & Loisirs valorisants, activités associatives \\ \hline
Références (facultatif) & Personnes à contacter pour attester du profil \\ \hline
\end{tabularx}
\end{center}

\textbf{Langue : valeurs des modes à connaître par cœur}
\begin{center}
\begin{tabularx}{\linewidth}{|l|X|X|}\hline
\rowcolor{popblueL} \textbf{Mode} & \textbf{Valeurs principales} & \textbf{Exemple} \\ \hline
Subjonctif & Souhait, volonté, obligation, doute, sentiment, hypothèse (après \emph{il faut que}, \emph{je veux que}, \emph{bien que}) & \emph{Il faut que je \textbf{sois} ponctuel.} \\ \hline
Conditionnel & Politesse, hypothèse (si + imparfait), souhait atténué, information non confirmée, conseil & \emph{Je \textbf{voudrais} postuler.} / \emph{Si j'avais un emploi, je \textbf{serais} autonome.} \\ \hline
\end{tabularx}
\end{center}

\textbf{Méthodes express}
\begin{itemize}
\item \cle{Rédiger} : 1) collecter ses informations ; 2) choisir les rubriques ; 3) classer en ordre antichronologique ; 4) rédiger court (phrases nominales, verbes d'action) ; 5) relire et corriger.
\item \cle{Adapter} : lire l'offre, repérer les mots-clés du poste, faire correspondre compétences et expériences, modifier l'objectif professionnel.
\item \cle{Présenter à l'oral} : se présenter en 2 minutes (identité, formation, atouts, motivation), employer le conditionnel de politesse.
\end{itemize}
\end{bilan}

\begin{aretenir}[Checklist avant le Bac]
$\square$ Je sais définir le CV comme un texte fonctionnel et citer sa finalité.\\
$\square$ Je connais les rubriques obligatoires : état civil, formation, expérience, compétences.\\
$\square$ Je classe les diplômes et expériences du plus récent au plus ancien.\\
$\square$ Je rédige des formulations brèves, avec des verbes d'action (\emph{gérer, réaliser, concevoir}).\\
$\square$ Je n'écris que des informations vraies et vérifiables.\\
$\square$ Je soigne la structure externe : une seule page, aérée, sans faute, titres visibles.\\
$\square$ Je sais insérer et commenter un élément iconique (photo d'identité).\\
$\square$ Je reconnais les valeurs du subjonctif (souhait, obligation, doute) dans une phrase.\\
$\square$ J'emploie le conditionnel de politesse et d'hypothèse (\emph{je voudrais, si j'avais\ldots}).\\
$\square$ Je sais adapter mon CV à une offre d'emploi précise (mots-clés, objectif).\\
$\square$ Je relis pour éliminer les fautes d'orthographe et de conjugaison.\\
$\square$ Je suis capable de présenter mon CV à l'oral en deux minutes.
\end{aretenir}

\findecours

\end{document}
